% Lesson 16 — Speaking: Exchanges information about using banking services — Anglais Tle A — Cours signé OnBuch+
% Programme officiel MINESEC. Compiler avec : tectonic EN16-speaking-exchanges-information-about-using-banking-serv.tex
\newcommand{\DOCMATIERE}{Anglais}
\newcommand{\DOCNIVEAU}{Tle A}
\newcommand{\DOCMODULE}{Chapter 2 : Economic Life and Occupations}
\newcommand{\DOCLECON}{EN16}
\newcommand{\DOCTITRE}{Lesson 16 — Speaking: Exchanges information about using banking services}
% OnBuch+ — préambule « cours signés » ENRICHI (compilé avec tectonic / XeLaTeX)
% Système visuel « Pop » d'OnBuch+ : fond crème (jamais blanc), bordures encre
% épaisses, ombres « dures » décalées, titres Archivo Black en capitales.
% Version MATHÉMATIQUES : graphes (pgfplots/tikz), boîtes pédagogiques
% (définition, propriété, méthode, attention, le savais-tu, exercice résolu,
% exercice, TP, bilan), page de garde et signature OnBuch+ ancrée en bas.
%
% Macros attendues dans main.tex AVANT \input{preamble} :
%   \DOCMATIERE  \DOCNIVEAU  \DOCMODULE  \DOCLECON (numéro)  \DOCTITRE
\documentclass[11pt]{article}

\usepackage{fontspec}
\usepackage[english,french]{babel} % français = langue principale ; l'anglais via \eng{..} / textebox
\usepackage[a4paper,margin=2cm,top=3.4cm,bottom=2.8cm,headheight=1.4cm,footskip=1.3cm]{geometry}
\usepackage{amsmath,amssymb,amsfonts}
\usepackage{mathtools} % \xmapsto, \coloneqq... souvent utilisés par les modèles
\usepackage{cancel} % simplifications barrées (bilans de forces)
% \abs{z} (module, valeur absolue) et \norm{u} (norme), utilisés par les modèles
\providecommand{\abs}{}\let\abs\relax\DeclarePairedDelimiter{\abs}{\lvert}{\rvert}
\providecommand{\norm}{}\let\norm\relax\DeclarePairedDelimiter{\norm}{\lVert}{\rVert}
\usepackage[table]{xcolor} % [table] : fournit aussi \rowcolors (alternance de lignes)
\usepackage{pagecolor}
\usepackage{tikz}
\usetikzlibrary{arrows.meta,positioning,calc,shapes.geometric,shapes.misc,shapes.arrows,decorations.pathmorphing,decorations.pathreplacing,decorations.markings,patterns,shadows,mindmap,trees,fit,backgrounds,matrix,intersections,angles,quotes}
\usepackage{pgfplots}
\usepackage[version=4]{mhchem} % équations nucléaires \ce{^{235}_{92}U}
\usepackage[european,siunitx]{circuitikz} % schémas électriques
\pgfplotsset{compat=1.18}
\usepgfplotslibrary{fillbetween,groupplots} % aires entre courbes (calcul intégral)
\usepackage{enumitem}
\usepackage{fancyhdr}
% [most] charge toutes les bibliothèques tcolorbox (listings, minted, raster,
% documentation...) et gonfle inutilement la mémoire TeX sur un document long
% (~300 boîtes) : on ne charge que ce qui est réellement utilisé.
\usepackage[skins,breakable]{tcolorbox}
\usepackage{graphicx}
\usepackage{colortbl}
\usepackage{booktabs}
\usepackage{tabularx}
\usepackage{diagbox} % case d'en-tête diagonale des tableaux à double entrée
% Colonnes extensibles centrée / gauche / droite (C, L, R), souvent utilisées par les modèles
\newcolumntype{C}{>{\centering\arraybackslash}X}
\newcolumntype{L}{>{\raggedright\arraybackslash}X}
\newcolumntype{R}{>{\raggedleft\arraybackslash}X}
\usepackage{array}
\usepackage{multicol}
\usepackage{multirow}
\usepackage{siunitx}
% parse-numbers=false : accepte aussi \SI{2,5 \times 10^{-3}}{...}
\sisetup{locale=FR,output-decimal-marker={,},group-minimum-digits=4,parse-numbers=false}
\DeclareSIUnit\ohm{\ensuremath{\symup{\Omega}}} % U+2126 absent de la police
\DeclareSIUnit\division{div} % réglages d'oscilloscope (ms/div, V/div)
\DeclareSIUnit\tr{tr} % tours (vitesse de rotation, ex. \SI{1500}{\tr\per\minute})
\DeclareSIUnit\molar{M} % concentration molaire (ex. \SI{3}{\molar}, \SI{1}{\micro\molar})
\DeclareSIUnit\an{an} % année (le modèle confond parfois avec \year, un compteur TeX) — ex. \SI{5}{\kilo\meter\per\an}
\DeclareSIUnit\FCFA{FCFA} % franc CFA (ex. \SI{500}{\FCFA\per\kilo\gram})
\DeclareSIUnit\degreeBrix{\textdegree Brix} % teneur en sucre d'un sirop/jus (réfractométrie)
\DeclareSIUnit\month{mois} % le modèle utilise parfois \month, un compteur TeX, comme unité de durée
\usepackage{qrcode}
% \degree : commande fréquemment utilisée par les modèles pour un angle ou une
% température en dehors de \SI{}{} (non fournie par les paquets chargés).
\providecommand{\degree}{^{\circ}}
% \qed (fin de démonstration), utilisé par les modèles sans amsthm
\providecommand{\qed}{\hfill$\square$}
% \barre : double barre des valeurs interdites dans un tableau de variations (tkz-tab)
\providecommand{\barre}{\ensuremath{\|}}
% environnement proof (amsthm non chargé) utilisé par les modèles
\ifdefined\proof\else\newenvironment{proof}[1][Démonstration]{\par\noindent\textit{#1.}\ }{\qed\par}\fi
\usepackage{lastpage}
\usepackage{hyperref}
\hypersetup{hidelinks}

% ============================================================
% Palette « Pop » OnBuch+ — mêmes hex que la classe P (plus_ui.dart)
% ============================================================
\definecolor{poporange}{HTML}{F59321}
\definecolor{popdark}{HTML}{E07A0C}
\definecolor{popcream}{HTML}{FFF6E8}
\definecolor{popink}{HTML}{1C1714}
\definecolor{poppurple}{HTML}{7A5AE0}
\definecolor{popgreen}{HTML}{2E9E6A}
\definecolor{poppink}{HTML}{E0457B}
\definecolor{popblue}{HTML}{2F7DE1}
\definecolor{popgold}{HTML}{C98A12}
\definecolor{popmuted}{HTML}{8A7F76}
\definecolor{popline}{HTML}{EADFD2}
\definecolor{popgray}{HTML}{8A7F76}
\colorlet{popgrey}{popgray}
% Alias tolérants (le modèle nomme parfois les couleurs différemment)
\colorlet{popwhite}{white}
\colorlet{popblack}{popink}
\colorlet{popred}{poppink}
\colorlet{popyellow}{popgold}
% Teintes claires (fonds de tableaux/encadrés) — le modèle nomme parfois
% les couleurs avec un suffixe L (ex. popblueL) sans les avoir définies.
\colorlet{poporangeL}{poporange!20}
\colorlet{popdarkL}{popdark!20}
\colorlet{poppurpleL}{poppurple!20}
\colorlet{popgreenL}{popgreen!20}
\colorlet{poppinkL}{poppink!20}
\colorlet{popblueL}{popblue!20}
\colorlet{popgoldL}{popgold!20}
\colorlet{popredL}{popred!20}
\colorlet{popgrayL}{popgray!20}
% Teintes claires pour fonds de tableaux / zones
\colorlet{popblueL}{popblue!10!white}
\colorlet{poporangeL}{poporange!14!white}
\colorlet{popgreenL}{popgreen!12!white}
\colorlet{poppurpleL}{poppurple!12!white}
\colorlet{poppinkL}{poppink!12!white}
\colorlet{popgoldL}{popgold!14!white}

\pagecolor{popcream}

% ============================================================
% Typographie
% ============================================================
\defaultfontfeatures{Ligatures=TeX}
\setmainfont{PlusJakartaSans-Medium.ttf}[Path=../../../fonts/,BoldFont=PlusJakartaSans-Bold.ttf]
% Symboles, API et emojis absents de la police principale : repli sur DejaVu Sans (caractères actifs)
\newfontface\symfont{DejaVuSans.ttf}[Path=../../../fonts/]
\makeatletter
\newcommand\symactive[1]{\catcode#1=13 \begingroup\lccode`\~=#1 \lowercase{\endgroup\def~}{{\symfont\char#1}}}
\newcommand\symblank[1]{\catcode#1=13 \begingroup\lccode`\~=#1 \lowercase{\endgroup\def~}{}}
\newcommand\symhyph[1]{\catcode#1=13 \begingroup\lccode`\~=#1 \lowercase{\endgroup\def~}{\mbox{-}}}
\newcount\symc
\newcommand\symrange[2]{\symc=#1 \@whilenum\symc<#2 \do{\expandafter\symactive\expandafter{\the\symc}\advance\symc by 1 }}
\newcommand\symblankrange[2]{\symc=#1 \@whilenum\symc<#2 \do{\expandafter\symblank\expandafter{\the\symc}\advance\symc by 1 }}
\makeatother
\symrange{"0250}{"02FF}\symactive{"2713}\symactive{"2714}\symactive{"2717}\symactive{"2718}\symactive{"2610}\symactive{"2611}\symactive{"2612}
\symhyph{"2011}\symhyph{"2E1A}\symblankrange{"1F300}{"1FAFF}\symblank{"FE0F}
% Maths en sans-serif (Fira Math) avec chiffres et lettres droites de Plus
% Jakarta, pour que les formules se fondent dans le texte.
\usepackage[math-style=ISO,bold-style=ISO]{unicode-math}
\setmathfont{FiraMath-Regular.otf}[Path=../../../fonts/]
\setmathfont{PlusJakartaSans-Medium.ttf}[Path=../../../fonts/,range={up/num,up/{Latin,latin},\mathup}]
\usepackage{icomma} % virgule décimale sans espace en mode math
% Caractères mathématiques Unicode absents des polices de texte (Plus Jakarta,
% Archivo Black) : rendus par leur équivalent mathématique, en texte comme en
% formule — sinon ils s'affichent en carré vide (ex. « Arithmétique dans ℤ »).
\usepackage{newunicodechar}
\newunicodechar{ℝ}{\ensuremath{\mathbb{R}}}
\newunicodechar{ℤ}{\ensuremath{\mathbb{Z}}}
\newunicodechar{ℂ}{\ensuremath{\mathbb{C}}}
\newunicodechar{ℕ}{\ensuremath{\mathbb{N}}}
\newunicodechar{ℚ}{\ensuremath{\mathbb{Q}}}
\newunicodechar{𝔻}{\ensuremath{\mathbb{D}}}
\newunicodechar{α}{\ensuremath{\alpha}}
\newunicodechar{β}{\ensuremath{\beta}}
\newunicodechar{γ}{\ensuremath{\gamma}}
\newunicodechar{δ}{\ensuremath{\delta}}
\newunicodechar{ε}{\ensuremath{\varepsilon}}
\newunicodechar{θ}{\ensuremath{\theta}}
\newunicodechar{λ}{\ensuremath{\lambda}}
\newunicodechar{μ}{\ensuremath{\mu}}
\newunicodechar{σ}{\ensuremath{\sigma}}
\newunicodechar{τ}{\ensuremath{\tau}}
\newunicodechar{φ}{\ensuremath{\varphi}}
\newunicodechar{ψ}{\ensuremath{\psi}}
\newunicodechar{ω}{\ensuremath{\omega}}
\newunicodechar{Σ}{\ensuremath{\Sigma}}
% majuscules produites par \MakeUppercase dans les titres (α-aminés -> Α)
\newunicodechar{Α}{\ensuremath{\alpha}}
\newunicodechar{Β}{\ensuremath{\beta}}
\newunicodechar{Γ}{\ensuremath{\Gamma}}
\newunicodechar{Δ}{\ensuremath{\Delta}}
\newunicodechar{Ε}{\ensuremath{\varepsilon}}
\newunicodechar{Θ}{\ensuremath{\Theta}}
\newunicodechar{Λ}{\ensuremath{\Lambda}}
\newunicodechar{Μ}{\ensuremath{\mu}}
\newunicodechar{Τ}{\ensuremath{\tau}}
\newunicodechar{Φ}{\ensuremath{\Phi}}
\newunicodechar{Ψ}{\ensuremath{\Psi}}
\newunicodechar{Ω}{\ensuremath{\Omega}}
\newunicodechar{↦}{\ensuremath{\mapsto}}
\newunicodechar{∈}{\ensuremath{\in}}
\newunicodechar{∉}{\ensuremath{\notin}}
\newunicodechar{∧}{\ensuremath{\wedge}}
\newunicodechar{⇔}{\ensuremath{\Leftrightarrow}}
\newunicodechar{⟺}{\ensuremath{\Longleftrightarrow}}
\newunicodechar{⇒}{\ensuremath{\Rightarrow}}
\newunicodechar{⟹}{\ensuremath{\Longrightarrow}}
\newunicodechar{∘}{\ensuremath{\circ}}
\newunicodechar{∪}{\ensuremath{\cup}}
\newunicodechar{∩}{\ensuremath{\cap}}
\newunicodechar{≡}{\ensuremath{\equiv}}
\newunicodechar{⊂}{\ensuremath{\subset}}
\newunicodechar{⊥}{\ensuremath{\perp}}
\newunicodechar{∃}{\ensuremath{\exists}}
\newunicodechar{∀}{\ensuremath{\forall}}
\newunicodechar{∛}{\ensuremath{\sqrt[3]{\,}}}
\newunicodechar{∜}{\ensuremath{\sqrt[4]{\,}}}
\newunicodechar{⃗}{\ensuremath{{}^{\to}}}
\newunicodechar{ⁿ}{\ifmmode^{n}\else\textsuperscript{\ensuremath{n}}\fi}
\newunicodechar{ˣ}{\ifmmode^{x}\else\textsuperscript{\ensuremath{x}}\fi}
\newunicodechar{ᵇ}{\ifmmode^{b}\else\textsuperscript{\ensuremath{b}}\fi}
\newunicodechar{ʳ}{\ifmmode^{r}\else\textsuperscript{\ensuremath{r}}\fi}
\newunicodechar{ᵘ}{\ifmmode^{u}\else\textsuperscript{\ensuremath{u}}\fi}
\newunicodechar{ʸ}{\ifmmode^{y}\else\textsuperscript{\ensuremath{y}}\fi}
\newunicodechar{ᵃ}{\ifmmode^{a}\else\textsuperscript{\ensuremath{a}}\fi}
\newunicodechar{ᵉ}{\ifmmode^{e}\else\textsuperscript{\ensuremath{e}}\fi}
\newunicodechar{ᵏ}{\ifmmode^{k}\else\textsuperscript{\ensuremath{k}}\fi}
\newunicodechar{ᵖ}{\ifmmode^{p}\else\textsuperscript{\ensuremath{p}}\fi}
\newunicodechar{ᴺ}{\ifmmode^{N}\else\textsuperscript{\ensuremath{N}}\fi}
\newunicodechar{ᶜ}{\ifmmode^{c}\else\textsuperscript{\ensuremath{c}}\fi}
\newunicodechar{ʰ}{\ifmmode^{h}\else\textsuperscript{\ensuremath{h}}\fi}
\newunicodechar{ᵠ}{\ifmmode^{q}\else\textsuperscript{\ensuremath{q}}\fi}
\newunicodechar{ₙ}{\ifmmode_{n}\else\textsubscript{\ensuremath{n}}\fi}
\newunicodechar{ₐ}{\ifmmode_{a}\else\textsubscript{\ensuremath{a}}\fi}
\newunicodechar{ᵢ}{\ifmmode_{i}\else\textsubscript{\ensuremath{i}}\fi}
\newunicodechar{ᵣ}{\ifmmode_{r}\else\textsubscript{\ensuremath{r}}\fi}
\newunicodechar{ₖ}{\ifmmode_{k}\else\textsubscript{\ensuremath{k}}\fi}
\newunicodechar{ᵦ}{\ifmmode_{\beta}\else\textsubscript{\ensuremath{\beta}}\fi}
\newunicodechar{ₓ}{\ifmmode_{x}\else\textsubscript{\ensuremath{x}}\fi}
\newfontfamily\popdisplay{ArchivoBlack-Regular.ttf}[Path=../../../fonts/]
\newfontfamily\poplogo{SpaceGrotesk-Bold.ttf}[Path=../../../fonts/]
\newfontfamily\popsemi{PlusJakartaSans-SemiBold.ttf}[Path=../../../fonts/]
\newfontfamily\popxbold{PlusJakartaSans-ExtraBold.ttf}[Path=../../../fonts/]
\newfontfamily\popxboldls{PlusJakartaSans-ExtraBold.ttf}[Path=../../../fonts/,LetterSpace=3.0]

\setlist[enumerate]{leftmargin=1.7em,itemsep=5pt,topsep=4pt}
\setlist[itemize]{leftmargin=1.7em,itemsep=4pt,topsep=4pt,label={\color{poporange}\textbullet}}
\setlength{\parindent}{0pt}
\setlength{\parskip}{5pt}
\renewcommand{\arraystretch}{1.25}

% pgfplots : style Pop par défaut
\pgfplotsset{
  every axis/.append style={axis line style={popink,line width=1pt},tick style={popink},
    grid style={popline,line width=0.5pt},label style={font=\small},tick label style={font=\footnotesize}},
  popcurve/.style={poporange,line width=1.8pt,no markers,smooth},
}
% Même style disponible dans un \draw TikZ simple (hors pgfplots)
\tikzset{popcurve/.style={poporange,line width=1.8pt}}
\tikzset{popgrid/.style={popline,very thin}}
\colorlet{popcurve}{poporange} % le nom du style est parfois utilisé comme couleur

% ============================================================
% Pill / squiggle
% ============================================================
\newcommand{\poppill}[3]{%
  \tikz[baseline=-0.55ex]\node[draw=popink,line width=1.2pt,rounded corners=2.6pt,%
    fill=#2,inner xsep=6.5pt,inner ysep=3pt]{%
    \popxboldls\fontsize{7}{7}\selectfont\color{#3}\MakeUppercase{#1}};%
}
\newcommand{\popsquiggle}[1]{%
  \tikz[baseline=-0.4ex]{
    \draw[#1,line width=2.6pt,line cap=round]
      plot [smooth] coordinates {(0,0.09) (0.34,0.01) (0.68,0.21) (1.02,0.01) (1.36,0.10)};
  }%
}
% Mot-clé surligné (marqueur orange sous le texte)
\newcommand{\cle}[1]{{\popxbold\color{popdark}#1}}

% ============================================================
% En-tête / pied
% ============================================================
\pagestyle{fancy}
\fancyhf{}
\renewcommand{\headrulewidth}{0pt}
\renewcommand{\footrulewidth}{0pt}
\lhead{\raisebox{-0.15ex}{\poplogo\fontsize{13}{13}\selectfont\color{popink}OnBuch\color{poporange}+}}
\rhead{\poppill{\DOCMATIERE\ \textbullet\ \DOCNIVEAU\ \textbullet\ Leçon \DOCLECON}{white}{popink}}
\lfoot{\popxbold\fontsize{7.6}{7.6}\selectfont\color{popmuted}\MakeUppercase{Cours signé OnBuch+}\ \textbullet\ {\popsemi\DOCTITRE}}
\rfoot{\tikz[baseline=-0.6ex]\node[rounded corners=8.5pt,fill=popink,minimum height=17pt,minimum width=17pt,inner xsep=5pt,inner sep=0]{\popxbold\fontsize{8}{8}\selectfont\color{popcream}\thepage\,/\,\pageref*{LastPage}};}

% ============================================================
% Titres
% ============================================================
\newcommand{\chaptitre}[2]{%
  \begin{tcolorbox}[enhanced,colback=poporange,colframe=popink,boxrule=2.6pt,arc=11pt,
      drop shadow=popink,left=15pt,right=15pt,top=12pt,bottom=11pt,before skip=3pt,after skip=18pt]
    \poppill{#2}{popink}{popcream}\par\vspace{9pt}
    {\raggedright\hyphenpenalty=10000\exhyphenpenalty=10000\popdisplay\fontsize{22}{24}\selectfont\color{popcream}\MakeUppercase{#1}\par}
  \end{tcolorbox}
}
% Section numérotée : pastille orange avec le numéro + titre Archivo
\newcounter{popsec}
\newcommand{\coursec}[1]{%
  \stepcounter{popsec}\setcounter{popsub}{0}%
  \par\vspace{16pt}\noindent
  \begin{minipage}{\linewidth}
  \tikz[baseline=-0.7ex]\node[draw=popink,line width=1.6pt,fill=poporange,rounded corners=4pt,minimum size=22pt,inner sep=0]{\popdisplay\fontsize{12}{12}\selectfont\color{popcream}\thepopsec};%
  \hspace{8pt}{\popdisplay\fontsize{15}{17}\selectfont\color{popink}\MakeUppercase{#1}}\par
  \vspace{4pt}\noindent{\color{poporange}\rule{36pt}{3.6pt}}
  \end{minipage}\par\nopagebreak\vspace{8pt}
}
\newcounter{popsub}
\newcommand{\courssub}[1]{%
  \stepcounter{popsub}%
  \par\vspace{9pt}\noindent{\popxbold\fontsize{11.5}{13}\selectfont\color{popink}{\color{poporange}\thepopsec.\thepopsub}\hspace{6pt}#1}\par\nopagebreak\vspace{4pt}
}

% ============================================================
% Boîtes pédagogiques — même recette : fond blanc, bordure épaisse colorée,
% ombre dure encre, badge plein. Titre optionnel [#1].
% ============================================================
\tcbset{popbase/.style={breakable,enhanced,colback=white,boxrule=2.3pt,arc=9pt,
  shadow={3pt}{-3pt}{0pt}{popink},left=14pt,right=14pt,top=11pt,bottom=11pt,before skip=11pt,after skip=15pt,
  fontupper=\popsemi\fontsize{10.6}{14.5}\selectfont\color{popink},
  fontlower=\popsemi\fontsize{10.6}{14.5}\selectfont\color{popink}}}
% Coche dessinée (le glyphe ✓ n'existe pas dans les polices du document)
\newcommand{\popcheck}[1]{\tikz[baseline=-0.5ex]\draw[#1,line width=1.6pt,line cap=round,line join=round](-0.13,0.0)--(-0.04,-0.09)--(0.13,0.1);}
\providecommand{\coche}{\popcheck{popgreen}}
\providecommand{\resultat}[1]{\par\smallskip\noindent\fcolorbox{popgreen}{popgreenL}{\parbox{\dimexpr\linewidth-2\fboxsep-2\fboxrule\relax}{\textbf{Résultat.} #1}}\par\smallskip}
\newcommand{\popboxhead}[3]{\poppill{#1}{#2}{white}\ifx\relax#3\relax\else\hspace{6pt}{\popxbold\fontsize{10.6}{12}\selectfont\color{popink}#3}\fi\par\vspace{7pt}}

\newtcolorbox{aretenir}[1][]{popbase,colframe=popgreen,before upper={\popboxhead{À retenir}{popgreen}{#1}}}
% Texte en anglais (document de lecture, dialogue, extrait) : cadre
% d'étude. Un titre optionnel (source, auteur) s'affiche dans l'en-tête.
\newtcolorbox{textebox}[1][]{popbase,colframe=popblue,before upper={\popboxhead{Text}{popblue}{#1}\begin{otherlanguage}{english}},after upper={\end{otherlanguage}}}
% Marques ✓ / ✗ (forme correcte / fautive) souvent utilisées par les modèles
\providecommand{\cross}{\textcolor{poppink}{\ensuremath{\times}}}
\providecommand{\xmark}{\cross}\providecommand{\crossmark}{\cross}\providecommand{\cmark}{\ensuremath{\checkmark}}
% Ligne de réponse / justification à compléter (inventée par les modèles)
\providecommand{\justification}[1]{\par\noindent\textit{Justification~:} \dotfill\par}
% \textipa (package tipa non chargé) : la police ne couvre pas l'API -> simple passage
\providecommand{\textipa}[1]{#1}
% Soulignement (ulem non chargé) : \uline, \ul
\providecommand{\uline}[1]{\underline{#1}}\providecommand{\ul}[1]{\underline{#1}}
% \glossaire{\item ...} inventée par les modèles : liste de glossaire
\providecommand{\glossaire}[1]{\begin{itemize}#1\end{itemize}}
% \alert (beamer) inventée par les modèles : texte mis en évidence
\providecommand{\alert}[1]{\textbf{\textcolor{poppink}{#1}}}
% variantes ASCII d'ordinaux écrites en macro
\providecommand{\iere}{\textsuperscript{re}}\providecommand{\ieme}{\textsuperscript{e}}\providecommand{\ere}{\textsuperscript{re}}\providecommand{\eme}{\textsuperscript{e}}\providecommand{\er}{\textsuperscript{er}}\providecommand{\ie}{\textsuperscript{e}}
% Ordinaux français écrits en macro par les modèles : 1\ière, 3\ième
\newcommand{\ière}{\textsuperscript{re}}\newcommand{\ième}{\textsuperscript{e}}
% \exemple (inventée par les modèles) : étiquette « Exemple : »
\providecommand{\exemple}{\textbf{Exemple~:}\ }
% \square utilisable aussi hors mode math (cases à cocher)
\AtBeginDocument{\let\squareMath\square\renewcommand{\square}{\ensuremath{\squareMath}}}
% Mot ou phrase en anglais à l'intérieur d'un paragraphe français
\newcommand{\eng}[1]{\ifmmode\text{\textit{#1}}\else\begingroup\selectlanguage{english}\itshape #1\endgroup\fi}
\let\esp\eng % alias toléré
% flèches d'intonation inventées par les modèles
\providecommand{\textsearrow}{}\renewcommand{\textsearrow}{\ensuremath{\searrow}}\providecommand{\textnearrow}{}\renewcommand{\textnearrow}{\ensuremath{\nearrow}}
% \fr{..} : mot français (inventé par les modèles)
\providecommand{\fr}[1]{\textit{#1}}
\newtcolorbox{exemplebox}[1][]{popbase,colframe=poporange,before upper={\popboxhead{Exemple}{poporange}{#1}}}
\newtcolorbox{definition}[1][]{popbase,colframe=popblue,before upper={\popboxhead{Définition}{popblue}{#1}}}
\newtcolorbox{propriete}[1][]{popbase,colframe=poppurple,before upper={\popboxhead{Propriété}{poppurple}{#1}}}
\newtcolorbox{methode}[1][]{popbase,colframe=popgold,before upper={\popboxhead{Méthode}{popgold}{#1}}}
\newtcolorbox{attention}[1][]{popbase,colframe=poppink,before upper={\popboxhead{Attention}{poppink}{#1}}}
\newtcolorbox{experience}[1][]{popbase,colframe=popblue,colback=popblueL,before upper={\popboxhead{Expérience}{popblue}{#1}}}
% « Le savais-tu ? » : carte violette pleine (contexte, culture, Cameroun)
\newtcolorbox{savaistu}[1][]{popbase,colback=poppurple,colframe=popink,
  fontupper=\popsemi\fontsize{10.4}{14.2}\selectfont\color{white},
  before upper={\poppill{Le savais-tu\,?}{white}{poppurple}\ifx\relax#1\relax\else\hspace{6pt}{\popxbold\color{white}#1}\fi\par\vspace{7pt}}}
% Exercice résolu : énoncé en haut, \tcblower puis la solution sur fond crème
\newcounter{popexo}\newcounter{popexr}
\newtcolorbox{exoresolu}[1][]{popbase,colframe=poporange,colbacklower=popcream,
  segmentation style={popink,line width=1.2pt,dashed},
  before upper={\stepcounter{popexr}\popboxhead{Exercice résolu \thepopexr}{poporange}{#1}},
  before lower={\poppill{Solution}{popink}{popcream}\par\vspace{6pt}}}
% Exercice (non résolu en place) — la correction est dans la partie Corrigés
\newtcolorbox{exercice}[1][]{popbase,colframe=popink,
  before upper={\stepcounter{popexo}\popboxhead{Exercice \thepopexo}{popink}{#1}}}
% Corrigé
\newtcolorbox{corrige}[1][]{popbase,colframe=popgreen,colback=popgreenL,
  before upper={\popboxhead{Corrigé}{popgreen}{#1}}}
% Objectifs / prérequis / bilan
\newtcolorbox{objectifs}{popbase,colframe=popink,colback=poporangeL,
  before upper={\popboxhead{Objectifs de la leçon}{popink}{}}}
\newtcolorbox{prerequis}{popbase,colframe=popink,colback=popblueL,
  before upper={\popboxhead{Prérequis}{popblue}{}}}
\newtcolorbox{bilan}{popbase,colframe=popink,colback=white,boxrule=2.8pt,
  before upper={\popboxhead{Fiche bilan}{poporange}{L'essentiel à connaître pour l'examen}}}
% Figure encadrée avec légende
\newtcolorbox{popfigure}[1][]{enhanced,breakable=false,colback=white,colframe=popline,boxrule=1.4pt,arc=8pt,
  left=10pt,right=10pt,top=10pt,bottom=8pt,before skip=10pt,after skip=12pt,halign=center,
  lower separated=false,halign lower=center,
  fontlower=\popsemi\fontsize{9}{11.5}\selectfont\color{popmuted},
  #1}
\newcommand{\legende}[1]{\par\vspace{6pt}{\popsemi\fontsize{9}{11.5}\selectfont\color{popmuted}#1}}

% ============================================================
% Signature OnBuch+
% ============================================================
\newcommand{\poplogoinline}[1]{{\poplogo\fontsize{#1}{#1}\selectfont\color{popink}OnBuch\color{poporange}+}}
\newcommand{\signatureonbuch}{%
  \begin{tcolorbox}[enhanced,colback=popink,colframe=popink,boxrule=2.2pt,arc=10pt,
      drop shadow=poporange,left=14pt,right=14pt,top=10pt,bottom=10pt,before skip=0pt,after skip=0pt]
    \tikz[baseline=-0.7ex]\node[circle,fill=popgreen,draw=popcream,line width=1.3pt,minimum size=22pt,inner sep=0]{\popcheck{white}};%
    \hspace{9pt}\begin{minipage}[c]{0.62\linewidth}
      {\popxbold\fontsize{12}{13}\selectfont\color{popcream}Signé\ \textbullet\ {\poplogo OnBuch\color{poporange}+}}\par\vspace{2pt}
      {\raggedright\popsemi\fontsize{8}{10.5}\selectfont\color{popline}\MakeUppercase{Équipe pédagogique OnBuch+}\\\MakeUppercase{Conforme au programme officiel MINESEC}\par}
    \end{minipage}\hfill
    {\poplogo\fontsize{22}{22}\selectfont\color{popcream}OnBuch\color{poporange}+}
  \end{tcolorbox}%
}

% ============================================================
% Page de garde
% #1 titre · #2 matière · #3 niveau · #4 module · #5 n° leçon · #6 durée ·
% #7 description / objectifs (texte ou itemize)
% ============================================================
\newcommand{\pagegarde}[7]{%
  \thispagestyle{empty}
  \newgeometry{a4paper,margin=1.7cm,top=1.6cm,bottom=1.4cm}%
  \begin{tikzpicture}[remember picture,overlay]
    \fill[poporange!18] ([xshift=-3.2cm,yshift=-3.6cm]current page.north east) circle (4.6cm);
    \fill[poppurple!14] ([xshift=2.4cm,yshift=7.6cm]current page.south west) circle (3.8cm);
  \end{tikzpicture}%
  {\poplogo\fontsize{30}{30}\selectfont\color{popink}OnBuch\color{poporange}+}\hfill
  \raisebox{0.6ex}{\poppill{Cours signé \textbullet\ Premium}{popink}{popcream}}\par
  \vspace{1pt}\popsquiggle{poppurple}\par
  \vspace{8pt}
  \begin{tcolorbox}[enhanced,colback=poporange,colframe=popink,boxrule=2.8pt,arc=13pt,
      drop shadow=popink,left=18pt,right=18pt,top=12pt,bottom=13pt,after skip=10pt]
    \poppill{#2 \textbullet\ #3}{popink}{popcream}\hspace{5pt}\poppill{#4}{white}{popink}\par\vspace{8pt}
    {\popdisplay\fontsize{14}{14}\selectfont\color{popink}LEÇON #5}\par\vspace{4pt}
    {\raggedright\hyphenpenalty=10000\exhyphenpenalty=10000\popdisplay\fontsize{25}{27}\selectfont\color{popcream}\MakeUppercase{#1}\par}
    \vspace{8pt}
    {\popxbold\fontsize{9.5}{9.5}\selectfont\color{popink}\MakeUppercase{Durée conseillée : #6}}
  \end{tcolorbox}
  \begin{tcolorbox}[enhanced,colback=white,colframe=popink,boxrule=2.2pt,arc=10pt,
      drop shadow=popink,left=16pt,right=16pt,top=10pt,bottom=10pt,after skip=8pt]
    \poppill{Dans cette leçon}{poppurple}{white}\par\vspace{5pt}
    \popsemi\fontsize{9.3}{12.6}\selectfont\color{popink}#7
  \end{tcolorbox}
  \noindent\begin{minipage}[t]{0.487\linewidth}
    \begin{tcolorbox}[enhanced,colback=popgreen,colframe=popink,boxrule=2.2pt,arc=10pt,
        drop shadow=popink,left=11pt,right=11pt,top=10pt,bottom=10pt,height=3.1cm,valign=center]
      \tcbox[colback=white,colframe=popink,boxrule=1.3pt,arc=5pt,boxsep=0pt,nobeforeafter,
        left=3pt,right=3pt,top=3pt,bottom=3pt,valign=center]{\qrcode[height=1.9cm]{https://play.google.com/store/apps/details?id=com.luvvix.onbuch&hl=fr}}%
      \hspace{8pt}\begin{minipage}[c]{0.5\linewidth}
        {\popdisplay\fontsize{11}{12.5}\selectfont\color{white}\MakeUppercase{Télécharge OnBuch}}\par\vspace{4pt}
        {\raggedright\popsemi\fontsize{8.4}{11}\selectfont\color{white}Gratuit \textbullet\ Google Play\\Tuteur IA, annales, concours, résultats\par}
      \end{minipage}
    \end{tcolorbox}
  \end{minipage}\hfill
  \begin{minipage}[t]{0.487\linewidth}
    \begin{tcolorbox}[enhanced,colback=poppurple,colframe=popink,boxrule=2.2pt,arc=10pt,
        drop shadow=popink,left=11pt,right=11pt,top=10pt,bottom=10pt,height=3.1cm,valign=center]
      \tikz[baseline=-0.7ex]\node[circle,fill=white,draw=popink,line width=1.3pt,minimum size=1.6cm,inner sep=0]{\popdisplay\fontsize{24}{24}\selectfont\color{poppurple}+};%
      \hspace{8pt}\begin{minipage}[c]{0.6\linewidth}
        {\popdisplay\fontsize{11}{12.5}\selectfont\color{white}\MakeUppercase{Passe à OnBuch+}}\par\vspace{4pt}
        {\raggedright\popsemi\fontsize{8.4}{11}\selectfont\color{white}Tous les cours signés, Tuteur IA illimité, fiches PDF — onglet OnBuch+ de l'app\par}
      \end{minipage}
    \end{tcolorbox}
  \end{minipage}
  \vspace{8pt}
  \popcontenu
  \vfill
  \signatureonbuch
  \restoregeometry
  \newpage
}

% Rangée « Ce cours contient » (page de garde)
\newcommand{\poptile}[3]{% #1 couleur #2 gros texte #3 légende
  \begin{tcolorbox}[enhanced,colback=#1,colframe=popink,boxrule=2pt,arc=9pt,shadow={2.5pt}{-2.5pt}{0pt}{popink},
      left=6pt,right=6pt,top=9pt,bottom=9pt,halign=center,valign=center,height=1.75cm,width=\linewidth,nobeforeafter]
    {\popdisplay\fontsize{12.5}{13}\selectfont\color{white}\MakeUppercase{#2}}\par\vspace{3pt}
    {\centering\popsemi\fontsize{7.8}{9.5}\selectfont\color{white}#3\par}
  \end{tcolorbox}}
\newcommand{\popcontenu}{%
  \par\noindent{\popxboldls\fontsize{7.5}{7.5}\selectfont\color{popmuted}CE COURS SIGNÉ CONTIENT}\par\vspace{7pt}
  \noindent\begin{minipage}[t]{0.235\linewidth}\poptile{popblue}{Cours}{complet, illustré, pas à pas}\end{minipage}\hfill
  \begin{minipage}[t]{0.235\linewidth}\poptile{poporange}{Méthodes}{et exercices résolus}\end{minipage}\hfill
  \begin{minipage}[t]{0.235\linewidth}\poptile{poppink}{Exercices}{type examen + corrigés}\end{minipage}\hfill
  \begin{minipage}[t]{0.235\linewidth}\poptile{popgreen}{Fiche bilan}{l'essentiel à retenir}\end{minipage}\par}

% Sommaire
\newtcolorbox{sommaire}{enhanced,colback=white,colframe=popink,boxrule=2.2pt,arc=10pt,
  drop shadow=popink,left=18pt,right=18pt,top=13pt,bottom=13pt,before skip=6pt,after skip=20pt,
  before upper={\poppill{Sommaire}{poppurple}{white}\par\vspace{9pt}\popsemi\fontsize{10.6}{14.5}\selectfont\color{popink}}}

% Signature de fin de cours — poussée tout en bas de la dernière page
\newcommand{\findecours}{%
  \par\vfill\nopagebreak
  \begin{center}\popsquiggle{poporange}\end{center}\vspace{4pt}
  \signatureonbuch
}
\AtBeginDocument{\def\llbracket{\mathopen{[\mkern-3.5mu[}}\def\rrbracket{\mathclose{]\mkern-3.5mu]}}} % intervalles d'entiers (glyphe ⟦ absent de la police)
% Glyphes absents de Fira Math : redéfinis à partir de glyphes disponibles
\AtBeginDocument{%
  \def\setminus{\mathbin{\backslash}}\let\smallsetminus\setminus
  \def\vdots{\mathord{\vbox{\baselineskip3.2pt\lineskiplimit0pt\kern4pt\hbox{.}\hbox{.}\hbox{.}}}}%
  \def\ddots{\mathinner{\mkern1mu\raise6.5pt\hbox{.}\mkern2mu\raise3.6pt\hbox{.}\mkern2mu\raise0.7pt\hbox{.}\mkern1mu}}%
  \def\triangle{\mathord{\scalebox{0.9}{$\symup{\Delta}$}}}%
  \def\nabla{\mathord{\raisebox{\depth}{\scalebox{1}[-1]{$\symup{\Delta}$}}}}%
  \def\checkmark{\popcheck{popdark}}%
  \def\star{\mathbin{\ast}}\def\bigstar{\mathord{\ast}}%
  \def\diamond{\mathbin{\scalebox{0.6}{$\Diamond$}}}%
  \def\bigcup{\mathop{\mathchoice{\raisebox{-0.3em}{\scalebox{1.7}{$\cup$}}}{\scalebox{1.25}{$\cup$}}{\cup}{\cup}}\displaylimits}%
  \def\bigcap{\mathop{\mathchoice{\raisebox{-0.3em}{\scalebox{1.7}{$\cap$}}}{\scalebox{1.25}{$\cap$}}{\cap}{\cap}}\displaylimits}%
  \def\lesseqqgtr{\mathrel{\lessgtr}}\def\bigoplus{\mathop{\oplus}\displaylimits}%
}
\providecommand{\ssi}{si et seulement si} % abréviation fréquente
\AtBeginDocument{\providecommand{\wideparen}{\overparen}} % arc de cercle

\begin{document}

\pagegarde{Lesson 16 — Speaking: Exchanges information about using banking services}{Anglais}{Tle A}{Chapter 2 : Economic Life and Occupations}{EN16}{2 h}{Cette leçon d'expression orale apprend à l'élève de Terminale A à échanger des informations sur l'utilisation des services bancaires : ouvrir un compte, déposer ou retirer de l'argent, demander un renseignement. À travers des dialogues modèles, des formules fonctionnelles et des jeux de rôle guidés, l'élève développe son aisance à l'oral avec une prononciation et une intonation correctes.
\vspace{4pt}
\begin{multicols}{2}\small
\begin{itemize}
\item À la fin de cette leçon, je sais utiliser le vocabulaire de base des services bancaires (account, deposit, withdrawal, balance, transfer...)
\item À la fin de cette leçon, je sais demander une information poliment à un employé de banque
\item À la fin de cette leçon, je sais donner une information claire et précise sur une opération bancaire
\item À la fin de cette leçon, je sais relancer la conversation quand je n'ai pas compris (Could you repeat that, please?)
\item À la fin de cette leçon, je sais ouvrir et conclure poliment un échange au guichet
\item À la fin de cette leçon, je sais mener un court dialogue complet en anglais sur les services bancaires
\end{itemize}\end{multicols}}

\begin{sommaire}
\begin{multicols}{2}
\begin{enumerate}
\item Découverte : un dialogue à la banque
\item Lexique des services bancaires
\item Formules fonctionnelles
\item Prononciation, accentuation et intonation
\item Dialogues modèles et jeux de rôle guidés
\item Entraînement au Bac et production finale
\item Activité d'intégration
\item Méthodes et exercices résolus
\item Exercices
\item Corrigés des exercices
\item Fiche bilan
\end{enumerate}
\end{multicols}
\end{sommaire}

\begin{objectifs}
\begin{itemize}
\item À la fin de cette leçon, je sais utiliser le vocabulaire de base des services bancaires (account, deposit, withdrawal, balance, transfer...)
\item À la fin de cette leçon, je sais demander une information poliment à un employé de banque
\item À la fin de cette leçon, je sais donner une information claire et précise sur une opération bancaire
\item À la fin de cette leçon, je sais relancer la conversation quand je n'ai pas compris (Could you repeat that, please?)
\item À la fin de cette leçon, je sais ouvrir et conclure poliment un échange au guichet
\item À la fin de cette leçon, je sais mener un court dialogue complet en anglais sur les services bancaires
\item À la fin de cette leçon, je sais prononcer correctement le lexique bancaire avec la bonne accentuation et intonation
\item À la fin de cette leçon, je sais jouer un rôle (client ou employé de banque) dans une situation authentique
\end{itemize}
\end{objectifs}

\begin{prerequis}
\begin{itemize}
\item Vocabulaire de l'argent et de la vie économique vu en Première (money, price, pay, buy, sell)
\item Questions en Wh- et en Yes/No (What...? How much...? Do you...?)
\item Formules de politesse de base (please, thank you, excuse me)
\item Présent simple et present continuous pour décrire des actions habituelles ou en cours
\item Modaux can/could/may pour demander poliment
\end{itemize}
\end{prerequis}

\newpage

\coursec{Découverte : un dialogue à la banque}

Imagine : tu viens de recevoir ta bourse d'études et tu dois ouvrir un compte à la banque de ton quartier à Yaoundé. Le guichetier te pose des questions en anglais : \eng{What type of account would you like to open?} Sauras-tu répondre, demander des précisions sur les frais, et conclure poliment ? Au Nigeria, au Ghana ou à Buea, ces échanges se font naturellement en anglais. Cette leçon t'apprend à mener ce dialogue avec aisance, comme le jour du Bac où tu devras échanger des informations sur les services bancaires.

\textbf{Problématique :} comment un échange à un guichet de banque est-il organisé, et quelles formules permettent de demander, donner et vérifier une information en anglais ?

\courssub{Première écoute : le dialogue modèle}

Avant toute explication, écoute (ou lis) ce dialogue sans chercher à tout traduire. Pose-toi trois questions simples : \emph{qui parle ? où ? pourquoi ?}

\begin{textebox}[Dialogue modèle 1 — At the SCB Bank in Bamenda]
\textbf{Clerk:} Good morning, Madam. Welcome to SCB Bank. How may I help you?

\textbf{Customer:} Good morning. I would like to open a savings account, please.

\textbf{Clerk:} Certainly. Do you have your ID card and a passport photograph?

\textbf{Customer:} Yes, here they are. How much money do I need to open the account?

\textbf{Clerk:} The minimum deposit is 10,000 francs CFA.

\textbf{Customer:} I see. And what documents will I receive?

\textbf{Clerk:} You will receive an account number and a cheque book within a week.

\textbf{Customer:} Thank you very much for your help.

\textbf{Clerk:} You're welcome. Have a nice day!
\end{textebox}

\begin{definition}[Glossaire du dialogue]
\begin{itemize}
\item \cle{bank clerk / teller} (nom) : guichetier, employé de banque.
\item \cle{customer} (nom) : client.
\item \cle{savings account} (nom) : compte d'épargne.
\item \cle{ID card} (nom) : carte nationale d'identité.
\item \cle{passport photograph} (nom) : photo d'identité (format passeport).
\item \cle{minimum deposit} (nom) : dépôt minimum (somme exigée à l'ouverture).
\item \cle{account number} (nom) : numéro de compte.
\item \cle{cheque book} (nom) : chéquier.
\item \cle{within a week} (locution) : d'ici une semaine, en moins d'une semaine.
\end{itemize}
\end{definition}

\begin{methode}[Comment écouter un dialogue en trois étapes]
\begin{enumerate}
\item \textbf{Première écoute — le sens global.} Ne traduis rien. Réponds seulement à trois questions : \emph{Who is speaking? Where are they? Why is the customer there?} Ici : un guichetier et une cliente, à la SCB de Bamenda, pour ouvrir un compte.
\item \textbf{Deuxième écoute — les informations précises.} Note les chiffres, les montants, les documents : \eng{10,000 francs CFA}, \eng{ID card}, \eng{passport photograph}, \eng{cheque book}, \eng{within a week}.
\item \textbf{Troisième écoute — les formules de politesse.} Repère comment on salue, demande, remercie et conclut : \eng{How may I help you?}, \eng{I would like...}, \eng{Thank you very much}, \eng{You're welcome}.
\end{enumerate}
\end{methode}

\courssub{Observation : les étapes du dialogue}

Relisons le dialogue en découpant chaque réplique. On s'aperçoit qu'un échange au guichet n'est pas un désordre : il suit presque toujours la même architecture en cinq étapes.

\begin{exemplebox}[Découpage du dialogue modèle]
\begin{itemize}
\item \textbf{Salutation (Greeting) :} \eng{Good morning, Madam. Welcome to SCB Bank.} — « Bonjour Madame. Bienvenue à la SCB Bank. »
\item \textbf{Demande (Request) :} \eng{I would like to open a savings account, please.} — « Je voudrais ouvrir un compte d'épargne, s'il vous plaît. »
\item \textbf{Informations (Information) :} \eng{The minimum deposit is 10,000 francs CFA.} — « Le dépôt minimum est de 10 000 francs CFA. »
\item \textbf{Questions de précision (Clarification) :} \eng{How much money do I need? And what documents will I receive?} — « Combien me faut-il ? Et quels documents vais-je recevoir ? »
\item \textbf{Conclusion (Closing) :} \eng{Thank you very much. — You're welcome. Have a nice day!} — « Merci beaucoup. — Je vous en prie. Bonne journée ! »
\end{itemize}
\end{exemplebox}

\begin{popfigure}
\begin{tikzpicture}[node distance=0.45cm,
  etape/.style={rectangle, rounded corners, draw=popblue, fill=popblueL, text width=4.6cm, align=center, minimum height=0.9cm, font=\small},
  fleche/.style={-{Stealth[length=2.5mm]}, thick, popdark}]
\node[etape, align=center] (e1){\textbf{1. Greeting}\\ \eng{Good morning! How may I help you?}};
\node[etape, below=of e1, align=center] (e2){\textbf{2. Request}\\ \eng{I would like to open an account.}};
\node[etape, below=of e2, align=center] (e3){\textbf{3. Information}\\ \eng{The minimum deposit is 10,000 FCFA.}};
\node[etape, below=of e3, align=center] (e4){\textbf{4. Clarification}\\ \eng{How much do I need? What documents?}};
\node[etape, below=of e4, align=center] (e5){\textbf{5. Closing}\\ \eng{Thank you. — You're welcome!}};
\draw[fleche] (e1) -- (e2);
\draw[fleche] (e2) -- (e3);
\draw[fleche] (e3) -- (e4);
\draw[fleche] (e4) -- (e5);
\end{tikzpicture}
\legende{Les cinq étapes d'un échange au guichet : Greeting, Request, Information, Clarification, Closing.}
\end{popfigure}

\begin{aretenir}[La structure d'un dialogue au guichet]
Tout échange de service en anglais suit cinq étapes : \cle{Greeting} (salutation), \cle{Request} (demande), \cle{Information} (donner les informations), \cle{Clarification} (poser des questions de précision), \cle{Closing} (conclure poliment). Retenir cette trame, c'est pouvoir construire n'importe quel dialogue au Bac, même avec un vocabulaire limité.
\end{aretenir}

\courssub{Identifier les rôles : customer et clerk}

Dans ce dialogue, deux personnages ont des rôles bien distincts, et le langage de chacun est différent.

\begin{center}
\begin{tabularx}{\linewidth}{|X|X|X|}
\hline
\rowcolor{popblueL} \textbf{Customer (client)} & \textbf{Bank clerk / teller (guichetier)} & \textbf{Remarque} \\
\hline
\eng{I would like to open an account.} & \eng{How may I help you?} & Le client demande, l'employé propose son aide. \\
\hline
\eng{How much money do I need?} & \eng{The minimum deposit is 10,000 francs.} & Le client pose des questions, l'employé donne les informations. \\
\hline
\eng{Thank you for your help.} & \eng{You're welcome. Have a nice day!} & Le client remercie, l'employé conclut. \\
\hline
\end{tabularx}
\end{center}

\begin{attention}[Pièges de traduction des francophones]
\begin{itemize}
\item \eng{customer} se prononce « keu-ste-meure » ; ne dis pas \emph{« costumer »} (un \eng{costumer} est un costumier de théâtre !).
\item \eng{I would like} est la forme polie obligatoire au guichet. \eng{I want to open an account} (« je veux ») est trop direct, presque impoli en anglais : préfère toujours \eng{I would like...} ou \eng{I'd like...}.
\item \eng{You're welcome} signifie « je vous en prie » (réponse à un remerciement), et non « vous êtes le bienvenu » dans ce contexte.
\item En anglais, on écrit les montants avec une virgule pour les milliers : \eng{10,000 francs} = 10 000 francs. Attention à ne pas lire « dix virgule zéro » !
\end{itemize}
\end{attention}

\courssub{Reformulation globale : de quoi parle-t-on ?}

Après l'écoute, l'examinateur du Bac te demandera souvent de reformuler le dialogue en une ou deux phrases. C'est l'épreuve de la compréhension globale. Voici le modèle à imiter.

\begin{exemplebox}[Reformuler le sens global]
\eng{The customer is at the bank. She wants to open a savings account. The clerk asks for her ID card and a passport photograph, and tells her the minimum deposit is 10,000 francs CFA. She will receive an account number and a cheque book within a week.}

« La cliente est à la banque. Elle veut ouvrir un compte d'épargne. Le guichetier lui demande sa carte d'identité et une photo d'identité, et lui indique que le dépôt minimum est de 10 000 francs CFA. Elle recevra un numéro de compte et un chéquier d'ici une semaine. »
\end{exemplebox}

\begin{propriete}[Reformuler : la règle d'or]
Pour reformuler un dialogue, utilise le présent simple (\eng{the customer wants...}, \eng{the clerk asks...}) et les verbes de parole au présent : \eng{says}, \eng{asks}, \eng{answers}, \eng{tells}. N'essaie jamais de traduire mot à mot : garde seulement \emph{qui}, \emph{où}, \emph{quoi}, \emph{combien}. Contrairement au français, l'anglais n'exige pas de subjonctif ni de concordance des temps compliquée dans ce type de résumé au présent.
\end{propriete}

\begin{exoresolu}[Questions de compréhension sur le dialogue]
Énoncé — Réponds en anglais, par des phrases complètes :
\begin{enumerate}
\item \eng{Where does the dialogue take place?}
\item \eng{What does the customer want to do?}
\item \eng{How much is the minimum deposit?}
\item \eng{What will the customer receive within a week?}
\end{enumerate}
\tcblower
Solution détaillée :
\begin{enumerate}
\item \eng{The dialogue takes place at SCB Bank in Bamenda.} — On repère l'indice \eng{Welcome to SCB Bank}.
\item \eng{She wants to open a savings account.} — Reformulation de \eng{I would like to open...} avec le verbe \eng{want} au présent.
\item \eng{The minimum deposit is 10,000 francs CFA.} — Information précise repérée à la deuxième écoute.
\item \eng{She will receive an account number and a cheque book.} — Attention au futur \eng{will receive}, repris tel quel du dialogue.
\end{enumerate}
\end{exoresolu}

\courssub{Première prise de parole : la répétition chorale}

La dernière étape de cette découverte est orale : tu vas dire le dialogue, d'abord tous ensemble, puis à deux. La répétition chorale (\eng{choral repetition}) permet de parler sans stress, car personne ne parle seul.

\begin{experience}[Répétition chorale guidée]
\begin{enumerate}
\item \textbf{Écoute et répète} chaque réplique après le professeur, toute la classe ensemble, en marquant l'intonation : voix qui monte pour les questions en \eng{How much...?}, voix qui descend pour les réponses.
\item \textbf{Deux groupes} : le groupe A joue le \eng{clerk}, le groupe B joue le \eng{customer}. Puis on échange les rôles.
\item \textbf{En binôme} : rejoue le dialogue sans regarder le texte, en gardant les cinq étapes (Greeting, Request, Information, Clarification, Closing). Tu peux changer le montant ou le type de compte.
\end{enumerate}
\end{experience}

\begin{attention}[Prononciation : les mots pièges du dialogue]
\begin{itemize}
\item \eng{clerk} : « kleurk » (en anglais britannique) — ne prononce pas le \emph{e} comme en français.
\item \eng{account} : « e-kaounte » — l'accent tombe sur la deuxième syllabe.
\item \eng{deposit} : « di-po-zite » — accent sur la deuxième syllabe.
\item \eng{cheque book} : « tchèque bouke » — \eng{cheque} se prononce exactement comme \eng{check}.
\item \eng{photograph} : « fô-teu-grafe » — accent sur la première syllabe.
\end{itemize}
\end{attention}

\begin{savaistu}[L'anglais bancaire, une compétence de vie réelle]
Au Cameroun anglophone — à Bamenda, Buea ou Limbe — tous les échanges bancaires se font en anglais : formulaires, guichets, distributeurs automatiques. Maîtriser ce dialogue n'est donc pas seulement un exercice scolaire pour le Bac : c'est une compétence de vie réelle. Le jour où tu ouvriras ton premier compte pour recevoir ton salaire ou ta bourse, tu utiliseras exactement les mêmes phrases : \eng{I would like to open an account. How much is the minimum deposit?}
\end{savaistu}

\begin{exercice}[À toi de jouer]
\begin{enumerate}
\item Remets les cinq étapes du dialogue dans l'ordre : \emph{Closing — Request — Greeting — Clarification — Information}.
\item Traduis en anglais : « Je voudrais ouvrir un compte, s'il vous plaît. Quel est le dépôt minimum ? »
\item Reformule le dialogue en deux phrases commençant par \eng{The customer...} et \eng{The clerk...}.
\end{enumerate}
\end{exercice}

\begin{corrige}[Exercice : À toi de jouer]
\begin{enumerate}
\item Ordre correct : \emph{Greeting — Request — Information — Clarification — Closing}.
\item \eng{I would like to open an account, please. How much is the minimum deposit?} (et non \emph{I want}, trop direct).
\item \eng{The customer wants to open a savings account. The clerk asks for her ID card and gives her the information about the deposit.}
\end{enumerate}
\end{corrige}

Tu sais maintenant écouter un dialogue bancaire, repérer ses cinq étapes, identifier les rôles et le reformuler. Dans la prochaine partie, nous enrichirons ton vocabulaire : tous les mots des services bancaires pour varier tes dialogues.

\coursec{Lexique des services bancaires}

Dans la section précédente, nous avons découvert un dialogue entre une cliente et un employé de banque à Douala. Ce dialogue contenait de nombreux mots du monde bancaire : \eng{account}, \eng{deposit}, \eng{balance}, \eng{withdraw}... Pour pouvoir échanger des informations sur les services bancaires — c'est la compétence visée par cette leçon — il faut d'abord maîtriser ce vocabulaire. C'est l'objet de cette section : observer, classer, mémoriser et utiliser le lexique essentiel de la banque.

\courssub{Texte support : At the bank in Yaoundé}

Lisons d'abord ce court texte. Souligne mentalement tous les mots qui appartiennent au champ lexical de la banque.

\begin{textebox}[At the bank in Yaoundé]
Last Monday, Mrs Ngono went to the main branch of her bank in the centre of Yaoundé. She waited in the queue for a few minutes, then she walked to the counter. A polite bank clerk greeted her and asked how he could help her.

Mrs Ngono explained that she wanted to open a savings account for her daughter, who had just started university in Buea. The clerk gave her a form to fill in and asked for her identity card. He also told her about the bank charges: the account was free for the first year, but after that, customers had to pay a small fee every month.

While she was there, Mrs Ngono also decided to check the balance of her current account. She was happy to see that her salary had arrived. She then made a deposit of 50,000 francs into her daughter's new account and withdrew 20,000 francs in cash to pay the house rent. The clerk asked her if she wanted a bank card. With an ATM card, she could withdraw money from any machine in town, day and night, simply by typing her PIN code.

Before leaving, Mrs Ngono asked the clerk about loans. Her husband wanted to borrow some money to buy a second-hand car for his taxi business. The clerk explained that the interest rate was quite reasonable at the moment and that the bank could lend them the money if their application was accepted. He gave her an appointment with the bank manager for the following Friday. Mrs Ngono thanked him warmly and left the bank with a smile.
\end{textebox}

\textbf{Glossaire}

\begin{center}
\begin{tabularx}{\linewidth}{|l|l|X|}
\hline
\rowcolor{popblueL} \textbf{Mot anglais} & \textbf{Nature} & \textbf{Traduction} \\
\hline
branch & nom & agence (bancaire) \\
\hline
counter & nom & guichet, comptoir \\
\hline
bank clerk & nom & employé(e) de banque, guichetier \\
\hline
to fill in & verbe & remplir (un formulaire) \\
\hline
bank charges & nom pluriel & frais bancaires \\
\hline
balance & nom & solde (d'un compte) \\
\hline
to withdraw & verbe & retirer (de l'argent) \\
\hline
PIN code & nom & code confidentiel \\
\hline
interest rate & nom & taux d'intérêt \\
\hline
to lend & verbe & prêter \\
\hline
\end{tabularx}
\end{center}

\textbf{Questions de compréhension}

\begin{enumerate}
\item Why did Mrs Ngono go to the bank? (Give two reasons.)
\item How much money did she deposit, and how much did she withdraw?
\item What can she do with an ATM card?
\item Why does her husband want a loan?
\end{enumerate}

\courssub{Vue d'ensemble : la carte mentale du lexique bancaire}

Le vocabulaire de ce texte se classe naturellement en cinq familles. La carte mentale suivante les résume ; nous allons ensuite étudier chaque famille en détail.

\begin{popfigure}
\begin{tikzpicture}[mindmap, grow cyclic, every node/.style={concept, text=white}, root concept/.append style={concept color=popblue, font=\bfseries}, level 1/.append style={level distance=3.4cm, sibling angle=72}, level 1 concept/.append style={font=\small\bfseries}]
\node[root concept]{BANKING SERVICES}
  child[concept color=poporange]{ node{People \& Places} }
  child[concept color=popgreen]{ node{Account} }
  child[concept color=poppurple]{ node{Operations} }
  child[concept color=poppink]{ node{Payment} }
  child[concept color=popgold]{ node{Fees \& Loans} };
\end{tikzpicture}
\legende{Carte mentale : les cinq familles du lexique bancaire}
\end{popfigure}

\courssub{Champ lexical 1 : les personnes et les lieux}

\begin{itemize}
\item \cle{customer} : le client, la cliente. \eng{The customer is waiting at the counter.} « Le client attend au guichet. »
\item \cle{bank clerk} ou \cle{teller} : l'employé(e) de banque, le/la guichetier(ère). \eng{Teller} est surtout américain ; \eng{clerk} se prononce « klârk » en anglais britannique, « kleurk » en américain.
\item \cle{bank manager} : le directeur / la directrice d'agence. \eng{You need an appointment with the bank manager.} « Il vous faut un rendez-vous avec le directeur d'agence. »
\item \cle{counter} : le guichet, le comptoir. \eng{Please go to counter number three.} « Veuillez vous présenter au guichet numéro trois. »
\item \cle{branch} : l'agence (bancaire). \eng{This bank has branches in Douala, Buea and Bamenda.} « Cette banque a des agences à Douala, Buea et Bamenda. »
\end{itemize}

\begin{attention}[Faux ami : branch]
\eng{Branch} signifie « agence, succursale » dans le contexte bancaire, mais aussi « branche (d'un arbre) ». En revanche, ne traduis jamais « branche » (d'un arbre) par un mot inventé : c'est bien \eng{branch}. Le piège est inverse : « une agence » ne se dit jamais \eng{an agency} pour une banque — on dit \eng{a branch} ou \eng{a bank branch}. \eng{Agency} convient pour une agence de voyages (\eng{travel agency}) ou immobilière (\eng{estate agency}).
\end{attention}

\courssub{Champ lexical 2 : le compte}

\begin{itemize}
\item \cle{to open an account} : ouvrir un compte ; \cle{to close an account} : fermer un compte. \eng{I would like to open an account, please.} « Je voudrais ouvrir un compte, s'il vous plaît. »
\item \cle{savings account} : compte d'épargne ; \cle{current account} : compte courant (compte chèque). Aux États-Unis, on dit \eng{checking account}.
\item \cle{account number} : numéro de compte. \eng{Can you give me your account number?} « Pouvez-vous me donner votre numéro de compte ? »
\item \cle{balance} : le solde. \eng{What is my balance?} « Quel est mon solde ? »
\end{itemize}

\begin{attention}[Faux ami : balance]
Le mot anglais \eng{balance} ressemble au français « balance », mais attention : \eng{balance} signifie « solde » (l'argent qui reste sur un compte), jamais l'instrument de pesée. « Une balance » (instrument) se dit \eng{scales} ou \eng{a scale} en anglais. \eng{I checked my balance at the ATM.} « J'ai vérifié mon solde au distributeur. »
\end{attention}

\courssub{Champ lexical 3 : les opérations}

\begin{itemize}
\item \cle{to deposit money} : déposer de l'argent ; \cle{a deposit} : un dépôt, un versement.
\item \cle{to withdraw money} : retirer de l'argent ; \cle{a withdrawal} : un retrait.
\item \cle{to transfer money} : transférer de l'argent, faire un virement ; \cle{a transfer} : un virement, un transfert.
\end{itemize}

\begin{definition}[Withdrawal]
Un \cle{withdrawal} (prononcé « ouiz-drô-ol ») est l'action de retirer de l'argent de son compte : un retrait. Le verbe correspondant est \cle{to withdraw money from an account} (verbe irrégulier : \eng{withdraw — withdrew — withdrawn}). Exemple : \eng{She made a withdrawal of 20,000 francs.} « Elle a fait un retrait de 20 000 francs. »
\end{definition}

\begin{exemplebox}[Exemples traduits]
\eng{I would like to check my balance.} « Je voudrais vérifier mon solde. »

\eng{She made a deposit of 50,000 francs.} « Elle a fait un dépôt de 50 000 francs. »

\eng{He applied for a loan.} « Il a demandé un prêt. »

\eng{Can I transfer money to my brother's account in Bamenda?} « Puis-je transférer de l'argent sur le compte de mon frère à Bamenda ? »
\end{exemplebox}

\courssub{Champ lexical 4 : les moyens de paiement}

\begin{itemize}
\item \cle{cash} : l'argent liquide, les espèces. Attention : \eng{cash} est indénombrable — jamais de \eng{a cash} ni de \eng{cashes}. \eng{Can I pay in cash?} « Puis-je payer en liquide ? »
\item \cle{cheque} (américain : \eng{check}) : un chèque ; \cle{cheque book} : un chéquier, un carnet de chèques.
\item \cle{bank card} ou \cle{ATM card} : une carte bancaire. \eng{ATM} se prononce « éï-ti-èm » et désigne aussi le distributeur automatique (\eng{cash machine}, \eng{cashpoint}).
\item \cle{PIN code} : le code confidentiel. \eng{PIN} signifie \eng{Personal Identification Number} ; on dit souvent simplement \eng{my PIN} (« mon code »).
\end{itemize}

\courssub{Champ lexical 5 : frais et prêts}

\begin{itemize}
\item \cle{bank charges} ou \cle{fees} : les frais bancaires. \eng{The bank charges a monthly fee.} « La banque prélève des frais mensuels. »
\item \cle{interest rate} : le taux d'intérêt. \eng{The interest rate on savings is low.} « Le taux d'intérêt sur l'épargne est faible. »
\item \cle{loan} : un prêt ; \cle{to apply for a loan} : demander un prêt, faire une demande de prêt.
\item \cle{to borrow} : emprunter ; \cle{to lend} : prêter.
\end{itemize}

\begin{attention}[To borrow ou to lend ?]
Erreur très fréquente chez les francophones, car le français utilise parfois « prêter » dans les deux sens familièrement. En anglais, la distinction est stricte :
\begin{itemize}
\item \cle{to borrow} = emprunter (on reçoit) : \eng{Can I borrow 5,000 francs?} « Puis-je t'emprunter 5 000 francs ? » On emprunte \emph{à} quelqu'un : \eng{borrow something from somebody}.
\item \cle{to lend} = prêter (on donne) : \eng{The bank lent me money.} « La banque m'a prêté de l'argent. » On prête \emph{à} quelqu'un : \eng{lend something to somebody}. Verbe irrégulier : \eng{lend — lent — lent}.
\end{itemize}
Astuce : \emph{borrow} = \emph{b}anque ? Non — retiens : le \emph{borrower} « boit » (reçoit), le \emph{lender} « laisse partir » (donne).
\end{attention}

\begin{attention}[Ne dis pas « remove money »]
Pour dire « retirer de l'argent », les francophones traduisent souvent mot à mot : \eng{I want to remove money}. C'est incorrect : \eng{to remove} signifie « enlever, ôter » (un vêtement, un objet). Dis : \eng{I want to withdraw money} (soutenu) ou \eng{I want to take out some money} (courant). De même, « retirer de l'argent au distributeur » : \eng{to withdraw cash from an ATM}.
\end{attention}

\courssub{Les collocations clés : le verbe juste avec le bon nom}

En anglais, chaque nom s'associe à un verbe précis : ce sont les \cle{collocations}. Le français dit « faire un dépôt », l'anglais dit \eng{make a deposit} — mais « ouvrir un compte » se traduit \eng{open an account}, pas \eng{start an account}. Ces blocs s'apprennent entiers, comme des formules.

\begin{propriete}[Collocations bancaires à mémoriser en blocs]
\begin{itemize}
\item \eng{to open an account} (ouvrir un compte) — jamais \eng{to start an account}.
\item \eng{to make a deposit} (faire un dépôt) — jamais \eng{to do a deposit}.
\item \eng{to withdraw cash} (retirer des espèces).
\item \eng{to check one's balance} (vérifier son solde).
\item \eng{to pay a fee} (payer des frais) — jamais \eng{to give a fee}.
\item \eng{to apply for a loan} (demander un prêt) — attention à la préposition \eng{for}.
\item \eng{to transfer money to an account} (virer de l'argent sur un compte).
\end{itemize}
\end{propriete}

\courssub{Tableau récapitulatif}

\begin{center}
\begin{tabularx}{\linewidth}{|l|X|X|}
\hline
\rowcolor{popblueL} \textbf{English} & \textbf{Français} & \textbf{Exemple (traduction)} \\
\hline
customer & client(e) & \eng{The customer signed the form.} (Le client a signé le formulaire.) \\
\hline
bank clerk / teller & employé de banque & \eng{Ask the teller for help.} (Demandez de l'aide au guichetier.) \\
\hline
branch & agence & \eng{Where is the nearest branch?} (Où est l'agence la plus proche ?) \\
\hline
savings account & compte d'épargne & \eng{She opened a savings account.} (Elle a ouvert un compte d'épargne.) \\
\hline
current account & compte courant & \eng{My salary goes into my current account.} (Mon salaire arrive sur mon compte courant.) \\
\hline
balance & solde & \eng{Your balance is low.} (Votre solde est bas.) \\
\hline
deposit & dépôt & \eng{He made a deposit.} (Il a fait un dépôt.) \\
\hline
withdrawal & retrait & \eng{Withdrawals are free here.} (Les retraits sont gratuits ici.) \\
\hline
transfer & virement & \eng{The transfer takes two days.} (Le virement prend deux jours.) \\
\hline
cash & espèces, liquide & \eng{I only have cash.} (Je n'ai que du liquide.) \\
\hline
cheque book & chéquier & \eng{I need a new cheque book.} (J'ai besoin d'un nouveau chéquier.) \\
\hline
ATM card & carte bancaire & \eng{Insert your ATM card.} (Insérez votre carte bancaire.) \\
\hline
PIN code & code confidentiel & \eng{Never tell anyone your PIN code.} (Ne donnez jamais votre code à personne.) \\
\hline
fee / charge & frais & \eng{There is a fee for this service.} (Ce service est payant.) \\
\hline
interest rate & taux d'intérêt & \eng{The interest rate is 5\%.} (Le taux d'intérêt est de 5 \%.) \\
\hline
loan & prêt & \eng{They applied for a loan.} (Ils ont demandé un prêt.) \\
\hline
to borrow / to lend & emprunter / prêter & \eng{I borrowed money from the bank.} (J'ai emprunté de l'argent à la banque.) \\
\hline
\end{tabularx}
\end{center}

\courssub{Exercice résolu}

\begin{exoresolu}[Complète avec le mot bancaire qui convient]
Énoncé. Complète chaque phrase avec le mot ou l'expression suivante : \eng{balance — withdraw — loan — deposit — PIN code — branch — fees}.

\begin{enumerate}
\item I went to the \dots\ of my bank near the market in Bafoussam.
\item She wants to \dots\ 10,000 francs from the cash machine.
\item Never write your \dots\ on your bank card!
\item He made a \dots\ of 100,000 francs into his savings account.
\item What is the \dots\ of my current account, please?
\item My uncle asked the bank for a \dots\ to buy a taxi.
\item Students pay no bank \dots\ during the first year.
\end{enumerate}
\tcblower
Solution détaillée.
\begin{enumerate}
\item \eng{branch} — il s'agit du lieu : une agence bancaire.
\item \eng{withdraw} — « retirer » de l'argent au distributeur ; \eng{remove} serait une erreur de francophone.
\item \eng{PIN code} — le code confidentiel, qu'il ne faut jamais révéler.
\item \eng{deposit} — « faire un dépôt » : collocation \eng{make a deposit}.
\item \eng{balance} — le solde du compte.
\item \eng{loan} — « demander un prêt » : \eng{apply for a loan}.
\item \eng{fees} — les frais bancaires ; on dit aussi \eng{bank charges}.
\end{enumerate}
\end{exoresolu}

\begin{exercice}[À toi de jouer]
Traduis en anglais :
\begin{enumerate}
\item Je voudrais ouvrir un compte d'épargne.
\item Puis-je vérifier mon solde, s'il vous plaît ?
\item La banque m'a prêté de l'argent pour acheter une moto.
\item Elle a retiré 25 000 francs au guichet.
\item Quel est le taux d'intérêt sur ce prêt ?
\end{enumerate}
\end{exercice}

\begin{corrige}[Exercice « À toi de jouer »]
\begin{enumerate}
\item \eng{I would like to open a savings account.}
\item \eng{Can I check my balance, please?} (ou \eng{May I check my balance, please?})
\item \eng{The bank lent me money to buy a motorbike.} (attention : \eng{lent}, passé de \eng{lend} — la banque \emph{prête}.)
\item \eng{She withdrew 25,000 francs at the counter.} (passé irrégulier : \eng{withdrew}.)
\item \eng{What is the interest rate on this loan?}
\end{enumerate}
\end{corrige}

\begin{savaistu}[Le lexique bancaire au Bac]
Au baccalauréat, le vocabulaire de la banque et de la vie économique apparaît régulièrement : dans les textes de compréhension écrite (\eng{reading}), dans les dialogues à compléter et dans les sujets d'expression sur la vie économique (par exemple : \eng{Discuss the advantages of saving money in a bank}). Apprendre ce lexique en blocs — avec ses collocations — te fait gagner des points en compréhension \emph{et} en production. Au Cameroun anglophone comme au Cameroun francophone, les services bancaires et le \eng{mobile money} font partie de la vie quotidienne : c'est un thème « réaliste » que les examinateurs aiment exploiter.
\end{savaistu}

Maintenant que le lexique est en place, passons à la section suivante : les formules fonctionnelles qui te permettront d'utiliser ce vocabulaire dans un dialogue réel à la banque.

\coursec{Formules fonctionnelles}

Dans la section précédente, nous avons construit le \cle{vocabulaire} des services bancaires : \eng{to open an account}, \eng{to withdraw money}, \eng{interest rate}, \eng{bank charges}... Mais des mots isolés ne suffisent pas pour parler. Il faut maintenant les assembler en \cle{phrases complètes} et, surtout, choisir la \cle{formule} adaptée à chaque moment du dialogue : demander un service, demander une information, donner une information, relancer, vérifier, conclure. C'est exactement ce que l'examinateur du Bac évalue dans l'épreuve d'expression orale : non pas la quantité de mots connus, mais la capacité à \emph{fonctionner} en anglais dans une situation réelle.

\courssub{Observer avant de mémoriser : un extrait de dialogue}

Lisons d'abord ce court échange entre un client et un employé de banque à Douala. Soulignons mentalement les formules qui reviennent.

\begin{textebox}[At the bank — first exchange]
\textbf{Clerk:} “Good morning, sir. How can I help you?”

\textbf{Customer:} “Good morning. I would like to open a savings account, please.”

\textbf{Clerk:} “Certainly. Do you have your ID card with you?”

\textbf{Customer:} “Yes, here it is. Could you tell me the minimum deposit, please?”

\textbf{Clerk:} “The minimum deposit is 10,000 francs.”

\textbf{Customer:} “Sorry, could you repeat that, please?”

\textbf{Clerk:} “Of course. Ten thousand francs.”

\textbf{Customer:} “So, if I understand correctly, I need my ID card and 10,000 francs?”

\textbf{Clerk:} “That's right. The account will be ready in a week.”

\textbf{Customer:} “Thank you very much for your help.”

\textbf{Clerk:} “You're welcome. Have a nice day!”
\end{textebox}

Observons : le client n'utilise jamais \eng{I want}. Il dit \eng{I would like to open...}, \eng{Could you tell me...?}, \eng{Sorry, could you repeat that?}. Chaque formule correspond à une \cle{fonction de communication}. C'est cette correspondance que nous allons systématiser.

\courssub{Demander un service : la règle d'or de la politesse}

\begin{propriete}[La formule standard pour demander un service]
Pour demander un service en anglais, la forme standard et polie est :

\textbf{\eng{I would like + nom} ou \eng{I would like + to + infinitif}}

La contraction \eng{I'd like} est très fréquente à l'oral.

\eng{I would like to make a transfer.} « Je voudrais faire un virement. »

\eng{I'd like to open an account, please.} « Je voudrais ouvrir un compte, s'il vous plaît. »

\eng{I would like some information about loans.} « Je voudrais des renseignements sur les prêts. »

\textbf{Attention :} \eng{I want...} sonne \emph{impoli} ou enfantin en anglais. À l'écrit formel, on peut monter d'un degré : \eng{I would like to...} (poli) ou \eng{I would be grateful if you could...} (très formel).
\end{propriete}

\begin{attention}[Le piège de la traduction mot à mot]
Les francophones traduisent souvent « je voudrais savoir » par \eng{I would want to know} : c'est une \cle{erreur}. En anglais, on dit :

\eng{I would like to know...} ou mieux encore \eng{Could you tell me...?}

De même, « je veux retirer de l'argent » ne se dit pas \eng{I want to withdraw money} dans une banque, mais \eng{I would like to withdraw some money} ou \eng{Could I withdraw some money?}
\end{attention}

\courssub{Demander une information : les questions utiles}

Pour obtenir un renseignement, deux grandes familles de formules :

\begin{enumerate}
\item \textbf{Les questions directes} avec un mot interrogatif : \eng{How much is the minimum deposit?} « Quel est le dépôt minimum ? » ; \eng{What documents do I need?} « Quels documents me faut-il ? » ; \eng{When will my card be ready?} « Quand ma carte sera-t-elle prête ? »
\item \textbf{Les questions indirectes}, plus polies, introduites par \eng{Could you tell me...?} : remarque que l'ordre des mots redevient celui d'une phrase affirmative (sujet + verbe), sans inversion.
\end{enumerate}

\begin{propriete}[Question directe ou indirecte ?]
\textbf{Directe :} \eng{What \textbf{are} the bank charges?} (inversion : verbe avant le sujet)

\textbf{Indirecte :} \eng{Could you tell me what the bank charges \textbf{are}?} (pas d'inversion après le mot interrogatif)

\textbf{Indirecte sans mot interrogatif :} \eng{Could you tell me \textbf{if} the account is free?} (on utilise \eng{if} quand la question directe serait une question fermée : \eng{Is the account free?})
\end{propriete}

\begin{exemplebox}[Questions fréquentes à la banque]
\eng{“Could you tell me the interest rate, please?”} « Pourriez-vous me dire le taux d'intérêt, s'il vous plaît ? »

\eng{“What are the bank charges for a transfer?”} « Quels sont les frais bancaires pour un virement ? »

\eng{“How long does it take to get a cheque book?”} « Combien de temps faut-il pour obtenir un chéquier ? »

\eng{“Could you tell me if I can use my card abroad?”} « Pourriez-vous me dire si je peux utiliser ma carte à l'étranger ? »
\end{exemplebox}

\courssub{Donner une information : répondre clairement}

Côté employé de banque (ou pour répondre à un camarade dans un jeu de rôle), on donne l'information avec des phrases simples et directes :

\begin{itemize}
\item \eng{“The minimum deposit is 10,000 francs.”} « Le dépôt minimum est de 10 000 francs. »
\item \eng{“You will need your ID card and two passport photos.”} « Il vous faudra votre carte d'identité et deux photos d'identité. »
\item \eng{“The account will be ready in a week.”} « Le compte sera prêt dans une semaine. »
\item \eng{“The interest rate is four per cent per year.”} « Le taux d'intérêt est de quatre pour cent par an. »
\item \eng{“There is a small charge for this service.”} « Il y a un petit frais pour ce service. »
\end{itemize}

Note l'emploi du futur \eng{will need} / \eng{will be ready} pour annoncer une information certaine, et de \eng{there is / there are} pour signaler l'existence d'un frais ou d'un document.

\courssub{Relancer, demander de répéter, vérifier : survivre et réussir}

C'est la partie la plus importante pour l'oral du Bac. Un bon candidat n'est pas celui qui comprend tout, mais celui qui \cle{gère la communication} quand il ne comprend pas.

\begin{methode}[Relancer sans paniquer]
\begin{enumerate}
\item Tu n'as pas compris un mot ou une phrase ? Ne reste \textbf{jamais silencieux}. Dis immédiatement : \eng{“Pardon?”} ou \eng{“Sorry, could you repeat that, please?”} « Pardon, pourriez-vous répéter, s'il vous plaît ? »
\item La personne parle trop vite ? Dis : \eng{“Could you speak more slowly, please?”} « Pourriez-vous parler plus lentement, s'il vous plaît ? »
\item Un mot précis te bloque ? Demande une explication : \eng{“What do you mean by ‘charges’?”} « Qu'entendez-vous par “charges” ? »
\item Tu crois avoir compris ? Vérifie en reformulant : \eng{“So, if I understand correctly, I need two photos?”} « Donc, si je comprends bien, il me faut deux photos ? » ou \eng{“Do you mean that the card is free?”} « Voulez-vous dire que la carte est gratuite ? »
\item À l'examen, ces stratégies sont \textbf{valorisées} : elles prouvent que tu sais maintenir un échange réel, comme dans la vraie vie.
\end{enumerate}
\end{methode}

\begin{attention}[Répéter ≠ répéter n'importe comment]
Ne dis pas \eng{“What?”} seul : c'est brutal en anglais. Préfère \eng{“Pardon?”} ou \eng{“Sorry?”} (avec l'intonation montante). Et n'oublie pas le \eng{please} final dans \eng{“Could you repeat that, please?”} : sans lui, la phrase peut sembler sèche.
\end{attention}

\courssub{Conclure poliment l'échange}

Un dialogue réussi se termine toujours par des formules de remerciement :

\begin{itemize}
\item \eng{“Thank you very much for your help.”} « Merci beaucoup pour votre aide. »
\item \eng{“That's very kind of you.”} « C'est très gentil de votre part. »
\item \eng{“You're welcome. Have a nice day!”} « Je vous en prie. Bonne journée ! » (réponse standard au remerciement)
\end{itemize}

\begin{attention}[Faux ami : « you're welcome »]
\eng{You're welcome} ne signifie pas « vous êtes le bienvenu » dans ce contexte : c'est la réponse polie à \eng{thank you}, l'équivalent de « je vous en prie » ou « de rien ».
\end{attention}

\courssub{Les degrés de politesse : Can I, Could I, May I ?}

L'anglais module la politesse comme le français module le tutoiement et le vouvoiement. Voici l'échelle, du plus direct au plus formel :

\begin{center}
\begin{tabularx}{\linewidth}{|l|X|X|}
\hline
\rowcolor{popblueL} Formule & Niveau & Exemple \\
\hline
\eng{I want...} & trop direct, à éviter & \eng{I want some money.} \\
\hline
\eng{Can I...?} & courant, familier & \eng{Can I withdraw some money?} \\
\hline
\eng{Could I...?} & poli, recommandé & \eng{Could I withdraw some money?} \\
\hline
\eng{May I...?} & très formel & \eng{May I see the manager, please?} \\
\hline
\eng{I would like to...} & poli, standard & \eng{I would like to make a transfer.} \\
\hline
\end{tabularx}
\end{center}

\begin{savaistu}[Pourquoi tant de politesse ?]
Dans la culture anglophone, la politesse linguistique est considérée comme un signe de respect professionnel. Un client qui dit \eng{I want money} dans une banque de Londres ou de Buea sera perçu comme rude, même si sa grammaire est correcte. À l'inverse, un simple \eng{Could I...?} ouvre toutes les portes. Retiens : en anglais, la forme compte autant que le fond.
\end{savaistu}

\courssub{Tableau récapitulatif : qui dit quoi ?}

\begin{popfigure}
\begin{tikzpicture}[
box/.style={draw=popblue, thick, rounded corners=3pt, fill=popblueL, text width=6.6cm, align=left, inner sep=5pt, font=\small},
fbox/.style={draw=poporange, thick, rounded corners=3pt, fill=poporangeL, text width=6.6cm, align=left, inner sep=5pt, font=\small},
head/.style={draw=popdark, thick, rounded corners=3pt, fill=popdark, text=white, text width=6.6cm, align=center, inner sep=5pt, font=\small\bfseries},
func/.style={draw=popgreen, thick, circle, fill=popgreenL, inner sep=3pt, font=\small\bfseries, text=popdark}
]
\node[head] (h1) at (0,0) {Customer says};
\node[head] (h2) at (7.6,0) {Clerk says};
\node[box] (a1) at (0,-1.9) {\textbf{Demander :} “I'd like to open an account, please.” “Could I withdraw some money?”};
\node[fbox] (a2) at (7.6,-1.9) {\textbf{Accueillir :} “Good morning! How can I help you?” “Certainly.”};
\node[box] (b1) at (0,-4.1) {\textbf{S'informer :} “Could you tell me the interest rate?” “What documents do I need?”};
\node[fbox] (b2) at (7.6,-4.1) {\textbf{Informer :} “The minimum deposit is 10,000 francs.” “You will need your ID card.”};
\node[box] (c1) at (0,-6.3) {\textbf{Relancer :} “Sorry, could you repeat that, please?” “What do you mean by ‘charges’?”};
\node[fbox] (c2) at (7.6,-6.3) {\textbf{Répéter :} “Of course.” “Ten thousand francs.” “It means the fees you pay.”};
\node[box] (d1) at (0,-8.5) {\textbf{Conclure :} “Thank you very much for your help.”};
\node[fbox] (d2) at (7.6,-8.5) {\textbf{Conclure :} “You're welcome. Have a nice day!”};
\node[func] (f1) at (3.8,-1.9) {1};
\node[func] (f2) at (3.8,-4.1) {2};
\node[func] (f3) at (3.8,-6.3) {3};
\node[func] (f4) at (3.8,-8.5) {4};
\draw[-{Stealth}, thick, popmuted] (h1.south) -- (a1.north);
\draw[-{Stealth}, thick, popmuted] (h2.south) -- (a2.north);
\end{tikzpicture}
\legende{Les quatre étapes du dialogue à la banque : 1. demander, 2. s'informer, 3. relancer, 4. conclure.}
\end{popfigure}

\courssub{Entraînement}

\begin{textebox}[Mini-échanges à compléter]
\textbf{Exchange 1}

A: “I would like to withdraw 20,000 francs.”

B: “Certainly. Do you have your \_\_\_\_\_\_\_\_ ?”

A: “Sorry, could you \_\_\_\_\_\_\_\_ that, please?”

B: “Your ID card. Do you have it?”

\textbf{Exchange 2}

A: “Could you tell me the \_\_\_\_\_\_\_\_ rate, please?”

B: “It is four per cent per year.”

A: “So, if I \_\_\_\_\_\_\_\_ correctly, the account is free?”

B: “That's right. There are no charges.”

\textbf{Exchange 3}

A: “I would like to \_\_\_\_\_\_\_\_ a savings account.”

B: “Of course. You will need two passport photos and your ID card.”

A: “Thank you very much for your \_\_\_\_\_\_\_\_.”

B: “You're welcome. Have a nice day!”
\end{textebox}

\begin{exoresolu}[Compléter et reformuler poliment]
\textbf{Énoncé.} (a) Complète les mini-échanges ci-dessus avec les mots qui conviennent. (b) Reformule ces demandes trop directes pour les rendre polies : 1. \eng{“I want to open an account.”} 2. \eng{“Give me the interest rate.”} 3. \eng{“What? Repeat!”}
\tcblower
\textbf{Solution.}

(a) Exchange 1 : \eng{ID card} ; \eng{repeat}. Exchange 2 : \eng{interest} ; \eng{understand}. Exchange 3 : \eng{open} ; \eng{help}.

(b) 1. \eng{“I would like to open an account, please.”} — on remplace \eng{I want} par \eng{I would like}, la formule standard de politesse.

2. \eng{“Could you tell me the interest rate, please?”} — un ordre (\eng{Give me...}) devient une question indirecte polie ; attention, pas d'inversion : on dit \eng{the interest rate is} et non \eng{is the interest rate} après \eng{Could you tell me}.

3. \eng{“Sorry, could you repeat that, please?”} — \eng{What?} seul est impoli ; on ajoute \eng{Sorry} et \eng{please} pour adoucir la demande.
\end{exoresolu}

\begin{exercice}[À ton tour]
Transforme chaque phrase française en formule anglaise polie, comme si tu étais dans une agence bancaire à Yaoundé :

1. « Je voudrais faire un virement, s'il vous plaît. »
2. « Pourriez-vous me dire quels documents il me faut ? »
3. « Pardon, pourriez-vous parler plus lentement ? »
4. « Donc, si je comprends bien, la carte sera prête dans une semaine ? »
5. « Merci beaucoup pour votre aide. »
\end{exercice}

\begin{corrige}[Exercice — À ton tour]
1. \eng{“I would like to make a transfer, please.”}

2. \eng{“Could you tell me what documents I need?”} (question indirecte : pas d'inversion)

3. \eng{“Sorry, could you speak more slowly, please?”}

4. \eng{“So, if I understand correctly, the card will be ready in a week?”}

5. \eng{“Thank you very much for your help.”}
\end{corrige}

\begin{aretenir}[L'essentiel des formules fonctionnelles]
\begin{itemize}
\item Demander un service : \eng{I would like to...} / \eng{Could I...?} — jamais \eng{I want}.
\item Demander une information : \eng{Could you tell me...?} (sans inversion) ou question directe avec \eng{How much / What / When}.
\item Donner une information : phrases simples avec \eng{is}, \eng{will need}, \eng{there is}.
\item Relancer : \eng{Sorry, could you repeat that, please?} ; vérifier : \eng{So, if I understand correctly...?}
\item Conclure : \eng{Thank you for your help.} / \eng{You're welcome. Have a nice day!}
\item Échelle de politesse : \eng{Can I} $<$ \eng{Could I} $<$ \eng{May I} ; à l'oral du Bac, vise toujours le niveau \eng{Could}.
\end{itemize}
\end{aretenir}

Ces formules sont maintenant tes outils. Dans la prochaine section, nous travaillerons leur \cle{prononciation}, leur \cle{accentuation} et leur \cle{intonation}, car une formule polie dite avec une mauvaise intonation peut perdre toute sa force.

\coursec{Prononciation, accentuation et intonation}

Dans la section précédente, tu as appris les formules fonctionnelles pour demander, donner et vérifier une information à la banque : \eng{I would like to open an account}, \eng{How much is the minimum deposit?}, \eng{Could you repeat that, please?} Mais à l'oral — et surtout à l'épreuve orale du Baccalauréat — il ne suffit pas de connaître les bonnes phrases : il faut les \cle{prononcer} correctement. Un examinateur juge trois choses : l'\cle{accentuation} (sur quelle syllabe insister), les \cle{sons difficiles} (comme le fameux \eng{th}), et l'\cle{intonation} (la mélodie de la phrase). Cette section t'entraîne à ces trois points, pas à pas.

\courssub{Observer d'abord : le même dialogue, deux mélodies}

\begin{textebox}[At the bank — listen to the melody]
\textbf{Customer:} Good morning. I would like to withdraw some cash, please. ↘\\
\textbf{Clerk:} Certainly. Do you have your ID card? ↗\\
\textbf{Customer:} Yes, here it is. ↘\\
\textbf{Clerk:} Thank you. How much would you like to withdraw? ↘\\
\textbf{Customer:} One hundred thousand francs, please. ↘\\
\textbf{Clerk:} Do you have an account with us? ↗\\
\textbf{Customer:} Yes, I do. ↘\\
\textbf{Clerk:} What is your account number? ↘\\
\textbf{Customer:} It is 045 221 387. ↘\\
\textbf{Clerk:} Perfect. Please sign here. Have a nice day! ↘
\end{textebox}

Observe bien les flèches. Les questions qui attendent \eng{yes} ou \eng{no} (\eng{Do you have your ID card?}, \eng{Do you have an account with us?}) se terminent par une voix qui \textbf{monte} ↗. Les questions qui commencent par \eng{How much}, \eng{What}, \eng{Where} se terminent par une voix qui \textbf{descend} ↘. Ce n'est pas un hasard : c'est une règle fondamentale de l'anglais parlé, que nous allons détailler plus bas.

\courssub{L'accentuation : frapper la bonne syllabe}

En français, toutes les syllabes ont à peu près la même force : \emph{u-ni-ver-si-té}. En anglais, chaque mot a une \cle{syllabe forte} (\eng{stressed syllable}) prononcée plus fort, plus longtemps et plus haut. Si tu frappes la mauvaise syllabe, ton interlocuteur peut ne pas te comprendre : \eng{dePOsit} est correct, \eng{DEposit} sonne bizarre.

\begin{propriete}[Règle d'accentuation des mots bancaires]
Dans un mot anglais, une seule syllabe est accentuée. Pour les mots de la banque, retiens par cœur la syllabe forte (écrite ici en MAJUSCULES) :
\begin{itemize}
\item \eng{acCOUNT} — e-KAOUNT — « un compte » (fort sur la 2\textsuperscript{e} syllabe)
\item \eng{dePOsit} — di-PO-zite — « un dépôt » (fort sur la 2\textsuperscript{e})
\item \eng{withDRAwal} — with-DRO-ol — « un retrait » (fort sur la 2\textsuperscript{e})
\item \eng{BALance} — BA-lensse — « le solde » (fort sur la 1\textsuperscript{re})
\item \eng{TRANSfer} — TRANSS-feure — « un virement » (fort sur la 1\textsuperscript{re} pour le nom)
\item \eng{INterest} — IN-teu-resste — « l'intérêt » (fort sur la 1\textsuperscript{re})
\item \eng{CUStomer} — KEUSS-teu-meure — « le client » (fort sur la 1\textsuperscript{re})
\end{itemize}
\end{propriete}

\begin{attention}[Nom ou verbe : l'accent change !]
Certains mots anglais changent d'accentuation selon leur nature grammaticale. Ainsi, le nom \eng{a TRANSfer} (« un virement ») s'accentue sur la première syllabe, mais le verbe \eng{to transFER} (« virer, transférer ») s'accentue sur la deuxième : \eng{I want to transFER money.} — « Je veux transférer de l'argent. » Même phénomène avec \eng{an INcrease} (nom) et \eng{to inCREASE} (verbe). Retiens : souvent, \textbf{nom = accent au début, verbe = accent à la fin}.
\end{attention}

\begin{methode}[Bien placer l'accent : la méthode des mains]
\begin{enumerate}
\item Découpe le mot en syllabes : \eng{de-po-sit} (3 syllabes).
\item Frappe dans tes mains à chaque syllabe, en frappant \textbf{plus fort} sur la syllabe accentuée : de (doux) — \textbf{PO} (fort) — sit (doux).
\item Répète le mot cinq fois avec les mains, puis sans les mains, en gardant le même rythme.
\item Place le mot dans une phrase et accentue-le toujours au même endroit : \eng{I want to make a dePOsit.}
\end{enumerate}
Pour le son \eng{th}, place ta main devant ta bouche : tu dois sentir un filet d'air quand la langue passe entre les dents. \eng{Thousand} se dit « saou-zend », jamais « tou-zend » ni « fou-zend ».
\end{methode}

\courssub{Les sons difficiles pour les francophones}

Trois sons reviennent constamment dans les dialogues bancaires et posent problème aux élèves francophones.

\begin{attention}[Les pièges de prononciation à la banque]
\begin{itemize}
\item \textbf{Le \eng{th}} : la langue sort légèrement entre les dents, l'air passe. \eng{Thousand} = « saou-zend », \eng{three} = « sri », \eng{the} = « ze » (th doux, voixé). Ne dis jamais « zeu bank » : dis « ze bènkeu », avec un \eng{th} doux et bref.
\item \textbf{Le \eng{a} de \eng{bank}} : ce n'est ni « a » ni « an » français. Dis « bènkeu », avec un son entre « è » et « e ». Même son dans \eng{cash} (« kèche »), \eng{transaction} (« tran-ZAK-cheune »).
\item \textbf{\eng{cheque} = \eng{check}} : les deux orthographes se prononcent exactement pareil, « tchèkeu ». Le \eng{ch} anglais se dit « tch » comme dans « match ».
\item \textbf{Le \eng{h}} : il est \textbf{muet} dans \eng{hour} (« aoueur ») mais \textbf{soufflé} dans \eng{how} (« haou »), \eng{have} (« hav »), \eng{hundred} (« heun-dred »). Ne confonds pas !
\end{itemize}
\end{attention}

\begin{savaistu}[\eng{Cheque} ou \eng{check} ?]
En Grande-Bretagne, on écrit \eng{cheque} ; aux États-Unis, on écrit \eng{check}. Mais les deux mots se prononcent exactement de la même façon : « tchèkeu ». Seule l'orthographe change ! Au Cameroun, pays bilingue où circulent les deux normes, tu rencontreras les deux écritures sur les documents bancaires. À l'écrit, choisis une norme et reste cohérent ; à l'oral, aucun problème : le son est identique.
\end{savaistu}

\courssub{L'intonation : la mélodie de la phrase}

L'\cle{intonation} est le mouvement de la voix à la fin de la phrase. En anglais, elle est codifiée et \textbf{systématiquement évaluée à l'oral du Bac}.

\begin{propriete}[Règle d'intonation des questions]
\begin{itemize}
\item \textbf{Question fermée} (réponse \eng{yes} ou \eng{no}, commençant par \eng{Do}, \eng{Can}, \eng{Have}, \eng{Is}...) : la voix \textbf{monte} ↗.\\ \eng{Do you have your ID card? ↗} (diou you hav your aï-DI kard?) — « Avez-vous votre carte d'identité ? »
\item \textbf{Question ouverte} (commençant par \eng{What}, \eng{Where}, \eng{How much}, \eng{When}...) : la voix \textbf{descend} ↘.\\ \eng{How much is the minimum deposit? ↘} (haou meutch iz ze MI-ni-meum di-PO-zite?) — « Quel est le dépôt minimum ? »
\item \textbf{Phrase affirmative ou demande polie} : la voix \textbf{descend} ↘.\\ \eng{I would like to open an account. ↘} (aï woud laïke tou ou-peune an a-KAOUNT) — « Je voudrais ouvrir un compte. »
\end{itemize}
\end{propriete}

\begin{popfigure}
\begin{tikzpicture}[scale=1.0]
% Flèche montante
\node[font=\bfseries, color=popblue] at (2.2,2.6) {Do you have an account?};
\draw[-{Stealth[length=4mm]}, line width=2.2pt, popblue] (0,0.3) .. controls (1.5,0.4) and (3,0.9) .. (4.4,2.0);
\node[color=popblue, font=\small, align=center] at (2.2,-0.5) {question fermée (yes/no) : la voix monte ↗};
% Flèche descendante
\node[font=\bfseries, color=poporange] at (9.8,2.6) {What is your account number?};
\draw[-{Stealth[length=4mm]}, line width=2.2pt, poporange] (7.6,2.0) .. controls (9,1.4) and (10.5,0.7) .. (12,0.3);
\node[color=poporange, font=\small, align=center] at (9.8,-0.5) {question ouverte (Wh-) : la voix descend ↘};
\end{tikzpicture}
\legende{Les deux mélodies fondamentales de la question en anglais : montante pour les questions fermées, descendante pour les questions ouvertes.}
\end{popfigure}

\courssub{Le rythme anglais : les mots pleins et les petits mots}

L'anglais est une langue à \cle{rythme accentué} : on insiste sur les \textbf{mots porteurs de sens} (noms, verbes, adjectifs, mots interrogatifs) et on « avale » les \textbf{petits mots} grammaticaux (\eng{the}, \eng{a}, \eng{of}, \eng{to}, \eng{is}). C'est le contraire du français, où chaque syllabe a le même poids.

\begin{exemplebox}[Phrases rythmées — les mots forts en MAJUSCULES]
\begin{itemize}
\item \eng{I would LIKE to withDRAW CASH. ↘} (aï woud \textbf{laïke} tou with-\textbf{DRO} \textbf{kèche}) — « Je voudrais retirer des espèces. »
\item \eng{WHAT are the BANK CHARGES? ↘} (\textbf{ouate} ar ze \textbf{bènkeu} \textbf{TCHAr-djiz}?) — « Quels sont les frais bancaires ? »
\item \eng{CAN I OPEN an acCOUNT HERE? ↗} (\textbf{kèn} aï \textbf{ou-peune} an a-\textbf{KAOUNT} \textbf{hi-eur}?) — « Puis-je ouvrir un compte ici ? »
\item \eng{The BANK OPENS at EIGHT. ↘} (ze \textbf{bènkeu} \textbf{ou-peunze} at \textbf{èite}) — « La banque ouvre à huit heures. »
\end{itemize}
Remarque : \eng{the} se dit « ze » (bref), \eng{a} se dit « e » (bref), \eng{to} se dit « teu » (bref). Ne les prononce jamais en entier dans une phrase rapide.
\end{exemplebox}

\begin{center}
\begin{tabularx}{\linewidth}{|X|X|X|}
\hline
\rowcolor{popblueL} Français & English & Prononciation figurée \\
\hline
Je voudrais ouvrir un compte. ↘ & I would like to open an account. ↘ & aï woud laïke tou ou-peune an a-KAOUNT \\
\hline
Avez-vous votre carte d'identité ? ↗ & Do you have your ID card? ↗ & diou you hav your aï-DI kard? \\
\hline
Quel est le dépôt minimum ? ↘ & What is the minimum deposit? ↘ & ouate iz ze MI-ni-meum di-PO-zite? \\
\hline
Je voudrais retirer des espèces. ↘ & I would like to withdraw cash. ↘ & aï woud laïke tou with-DRO kèche \\
\hline
Quels sont les frais bancaires ? ↘ & What are the bank charges? ↘ & ouate ar ze bènkeu TCHAr-djiz? \\
\hline
Puis-je faire un virement ? ↗ & Can I make a transfer? ↗ & kèn aï mèike e TRANSS-feure? \\
\hline
Quel est le solde de mon compte ? ↘ & What is the balance of my account? ↘ & ouate iz ze BA-lensse ov maï a-KAOUNT? \\
\hline
Acceptez-vous les chèques ? ↗ & Do you accept cheques? ↗ & diou you ak-SEPTE tchèkes? \\
\hline
\end{tabularx}
\end{center}

\courssub{Entraînement guidé : de la classe au binôme}

\begin{experience}[Répétition chorale puis travail en binôme]
\begin{enumerate}
\item \textbf{Répétition chorale} : le professeur lit chaque phrase du tableau ci-dessus ; la classe répète en chœur, en levant la main quand la voix monte ↗ et en la baissant quand elle descend ↘.
\item \textbf{Frappe des mains} : la classe répète les mots clés en frappant fort sur la syllabe accentuée : \eng{de-PO-sit}, \eng{with-DRA-wal}, \eng{BAL-ance}, \eng{CUS-to-mer}.
\item \textbf{Binôme} : l'élève A joue le client, l'élève B l'employé de banque. A pose deux questions fermées (voix montante) et deux questions ouvertes (voix descendante) ; B répond avec des phrases affirmatives (voix descendante). Puis on échange les rôles.
\item \textbf{Contrôle mutuel} : chaque partenaire vérifie trois points — le \eng{th}, la syllabe forte, la flèche d'intonation — et signale les erreurs.
\end{enumerate}
\end{experience}

\begin{exoresolu}[Mets la bonne flèche et la bonne prononciation]
Pour chaque phrase, indique si l'intonation finale monte ↗ ou descend ↘, et écris la prononciation figurée en lettres françaises.\\
1. \eng{Can I open a savings account?}\\
2. \eng{Where is the nearest bank?}\\
3. \eng{I would like to deposit a cheque.}\\
4. \eng{Do you accept cash?}
\tcblower
\textbf{Solution détaillée :}\\
1. Question fermée (commence par \eng{Can}, réponse oui/non) → voix \textbf{montante} ↗ : « kèn aï ou-peune e SÈI-vin-guez a-KAOUNT? »\\
2. Question ouverte (commence par \eng{Where}) → voix \textbf{descendante} ↘ : « ouère iz ze NI-resste bènkeu? »\\
3. Phrase affirmative (demande polie) → voix \textbf{descendante} ↘ : « aï woud laïke tou di-PO-zite e tchèke » — attention, \eng{cheque} = « tchèkeu ».\\
4. Question fermée (commence par \eng{Do}) → voix \textbf{montante} ↗ : « diou you ak-SEPTE kèche? » — attention au \eng{a} de \eng{cash} : « kèche », pas « kache ».
\end{exoresolu}

\begin{exercice}[À toi de jouer]
Ajoute la flèche d'intonation correcte (↗ ou ↘) à la fin de chaque phrase, puis souligne la syllabe forte des mots en gras.\\
1. \eng{Could you tell me the \textbf{balance} of my account?}\\
2. \eng{Is there a \textbf{minimum} deposit?}\\
3. \eng{How much \textbf{interest} does this account pay?}\\
4. \eng{I would like to \textbf{transfer} some money to Bamenda.}\\
5. \eng{Do you \textbf{accept} cheques here?}
\end{exercice}

\begin{corrige}[Exercice « À toi de jouer »]
1. Demande polie → ↘ ; \eng{BALance} (1\textsuperscript{re} syllabe).\\
2. Question fermée (\eng{Is there...?}) → ↗ ; \eng{MI-nimum} (1\textsuperscript{re} syllabe).\\
3. Question ouverte (\eng{How much...?}) → ↘ ; \eng{INterest} (1\textsuperscript{re} syllabe).\\
4. Phrase affirmative → ↘ ; ici \eng{transfer} est un \textbf{verbe} : \eng{transFER} (2\textsuperscript{e} syllabe).\\
5. Question fermée (\eng{Do you...?}) → ↗ ; \eng{acCEPT} (2\textsuperscript{e} syllabe).
\end{corrige}

\begin{aretenir}[L'essentiel de la section]
\begin{itemize}
\item Chaque mot anglais a une syllabe forte : \eng{acCOUNT}, \eng{dePOsit}, \eng{withDRAwal}, \eng{BALance}, \eng{TRANSfer} (nom), \eng{INterest}, \eng{CUStomer}.
\item Le \eng{th} se prononce la langue entre les dents : \eng{thousand} = « saou-zend », \eng{the} = « ze ».
\item \eng{cheque} (GB) et \eng{check} (US) se prononcent tous deux « tchèkeu ».
\item Question fermée (oui/non) → voix montante ↗ ; question ouverte (Wh-) et phrase affirmative → voix descendante ↘.
\item Accentue les mots porteurs de sens, avale les petits mots (\eng{the}, \eng{a}, \eng{of}, \eng{to}).
\end{itemize}
\end{aretenir}

Maintenant que ta prononciation, ton accentuation et ton intonation sont prêtes, tu peux passer à la section suivante : les dialogues modèles et les jeux de rôle guidés, où tu mettras en pratique tout ce que tu viens d'apprendre dans de vraies conversations à la banque.

\coursec{Dialogues modèles et jeux de rôle guidés}

Maintenant que vous savez prononcer correctement les mots-clés du monde bancaire (\eng{withdrawal} « ouiz-drô-ol », \eng{balance} « ba-lens », \eng{account} « a-kaount ») et placer l'intonation montante sur les questions polies, il est temps de passer à la pratique. C'est ici que tout se joue : au Bac, l'épreuve d'expression orale ne vous demande pas de réciter, mais de \cle{réagir} et d'\cle{échanger}. Nous allons d'abord observer deux dialogues modèles, en souligner les formules fonctionnelles, puis progresser en trois étapes : dialogue à trous, cartes de rôle guidées, improvisation libre.

\courssub{Dialogue modèle 2 : retirer de l'argent au guichet}

Lisez d'abord le dialogue en silence, puis à voix haute avec un camarade. Observez comment le client annonce sa demande, comment l'employé vérifie son identité, et comment l'échange se conclut poliment.

\begin{textebox}[Dialogue 2 — At the counter: withdrawing money]
Customer: Good afternoon. I would like to withdraw 30,000 francs from my account.

Clerk: Certainly, Sir. May I see your ID card and your account number?

Customer: Here you are.

Clerk: Thank you. Please sign here... Here is your money.

Customer: Could you tell me my balance, please?

Clerk: Your balance is 75,500 francs.

Customer: Thank you. Goodbye!

Clerk: Goodbye, Sir!
\end{textebox}

\textbf{Glossaire.} \emph{to withdraw} (verbe irrégulier : withdraw, withdrew, withdrawn) : retirer (de l'argent) ; \emph{ID card} (nom) : carte d'identité ; \emph{account number} (nom) : numéro de compte ; \emph{Here you are} (expression) : voici, je vous le passe ; \emph{to sign} (verbe) : signer ; \emph{balance} (nom) : solde.

\textbf{Analyse guidée : soulignez les formules fonctionnelles.} Repérez dans le dialogue les répliques qui remplissent chaque fonction :

\begin{center}
\begin{tabularx}{\linewidth}{|X|X|}\hline
\rowcolor{popblueL} Fonction de communication & Réplique du dialogue \\ \hline
Saluer & \eng{Good afternoon.} \\ \hline
Annoncer sa demande & \eng{I would like to withdraw 30,000 francs.} \\ \hline
Demander poliment (employé) & \eng{May I see your ID card?} \\ \hline
Donner un document & \eng{Here you are.} \\ \hline
Donner une instruction & \eng{Please sign here.} \\ \hline
Demander une information & \eng{Could you tell me my balance, please?} \\ \hline
Donner une information & \eng{Your balance is 75,500 francs.} \\ \hline
Remercier et conclure & \eng{Thank you. Goodbye!} \\ \hline
\end{tabularx}
\end{center}

\begin{exemplebox}[Trois répliques à connaître par cœur]
\begin{itemize}
\item \eng{May I see your ID card?} = « Puis-je voir votre carte d'identité ? » — demande polie avec \cle{May I...?}
\item \eng{Please sign here.} = « Veuillez signer ici. » — instruction polie avec \cle{Please} + base verbale.
\item \eng{Your balance is 75,500 francs.} = « Votre solde est de 75 500 francs. » — donner une information chiffrée.
\end{itemize}
\end{exemplebox}

\begin{attention}[Here you are ≠ « vous y êtes »]
\eng{Here you are} signifie « voici, je vous le donne » quand on tend un objet. Ne traduisez jamais mot à mot. De même, \eng{I would like to withdraw} est bien plus poli que \eng{I want to withdraw} (« je veux »), qui sonne brutal en anglais. Au guichet, utilisez toujours \cle{I would like...} ou \cle{Could I...?} Attention aussi aux nombres : en anglais, on lit les montants chiffre par chiffre ou par groupes — \eng{75,500 francs} se dit « seventy-five thousand five hundred francs ».
\end{attention}

\courssub{Dialogue modèle 3 : demander des informations sur une carte bancaire}

Ce troisième dialogue montre une situation plus riche : le client pose plusieurs questions, l'employé répond, et le client \cle{relance} quand il veut une précision. C'est exactement le comportement noté au Bac.

\begin{textebox}[Dialogue 3 — Asking about a bank card]
Customer: Good morning, Madam. I would like some information about your bank cards.

Clerk: Of course, Sir. We offer a Visa card and a prepaid card.

Customer: How much does the Visa card cost per year?

Clerk: It costs 10,000 francs per year.

Customer: Oh, really? And can I use it at any cash machine in Douala?

Clerk: Yes, you can use it at any ATM, and also abroad.

Customer: That's interesting. How long does it take to get the card?

Clerk: About one week. You just need your ID card and a passport photo.

Customer: Only one week? That's fine! Could I apply today?

Clerk: Certainly. Please fill in this form.

Customer: Thank you very much for your help. Goodbye!

Clerk: You're welcome, Sir. Goodbye!
\end{textebox}

\textbf{Glossaire.} \emph{prepaid card} : carte prépayée ; \emph{cash machine / ATM} : distributeur automatique de billets ; \emph{abroad} : à l'étranger ; \emph{to apply} : faire une demande, poser sa candidature ; \emph{to fill in a form} : remplir un formulaire.

\textbf{Analyse guidée.} Soulignez dans ce dialogue : deux questions avec \eng{How much...?} et \eng{How long...?}, une question fermée avec \eng{Can I...?}, et surtout les réactions du client : \eng{Oh, really?}, \eng{That's interesting.}, \eng{Only one week? That's fine!}. Ces petites phrases montrent à l'examinateur que vous \cle{écoutez} votre partenaire.

\begin{propriete}[Les mots interrogatifs utiles à la banque]
\begin{itemize}
\item \eng{How much does it cost?} = « Combien cela coûte-t-il ? » (prix, montant indénombrable).
\item \eng{How long does it take?} = « Combien de temps cela prend-il ? » (durée).
\item \eng{What documents do I need?} = « Quels documents me faut-il ? »
\item \eng{Can I use it abroad?} = « Puis-je l'utiliser à l'étranger ? » (question fermée : réponse \eng{Yes, you can. / No, you can't.}).
\end{itemize}
Rappel : après \eng{How much / How long / What}, on inverse avec l'auxiliaire \cle{do/does} : \eng{How much \textbf{does} the card \textbf{cost}?} et jamais \eng{How much costs the card?}
\end{propriete}

\begin{attention}[L'interaction, pas la récitation]
Au Bac, l'examinateur note l'\cle{interaction}, pas la récitation. Écoute ton partenaire et réagis à ce qu'il dit réellement : \eng{Oh, really? Only one week? That's fine!} Un dialogue appris par cœur où chacun débite sa réplique sans tenir compte de la réponse de l'autre est sanctionné. Si tu ne comprends pas, relance : \eng{Sorry, could you repeat that, please?} (« Désolé, pourriez-vous répéter, s'il vous plaît ? ») ou \eng{What do you mean by “fees”?} (« Qu'entendez-vous par “frais” ? »).
\end{attention}

\courssub{Progression des jeux de rôle : du guidé au libre}

Pour passer du dialogue appris au dialogue improvisé, nous suivons trois étapes de difficulté croissante. Chaque étape retire un peu plus de support, comme quand on enlève les petites roues d'un vélo.

\begin{popfigure}
\begin{center}
\begin{tikzpicture}[node distance=0.6cm]
\node[draw=popblue, fill=popblueL, rounded corners, align=center, minimum width=3.6cm, minimum height=1.4cm] (a) {\textbf{Étape 1}\\ Dialogue à trous\\ (très guidé)};
\node[draw=poporange, fill=poporangeL, rounded corners, align=center, minimum width=3.6cm, minimum height=1.4cm, right=1.2cm of a] (b) {\textbf{Étape 2}\\ Cartes de rôle\\ (semi-guidé)};
\node[draw=popgreen, fill=popgreenL, rounded corners, align=center, minimum width=3.6cm, minimum height=1.4cm, right=1.2cm of b] (c) {\textbf{Étape 3}\\ Improvisation\\ (libre)};
\draw[-{Stealth}, very thick, popdark] (a) -- (b);
\draw[-{Stealth}, very thick, popdark] (b) -- (c);
\node[above=0.15cm of a, popmuted, align=center] {\small support maximal};
\node[above=0.15cm of c, popmuted, align=center] {\small support nul};
\end{tikzpicture}
\end{center}
\legende{Frise de progression des jeux de rôle : du dialogue à trous à l'improvisation libre.}
\end{popfigure}

\courssub{Jeu de rôle 1 (guidé) : dialogue à trous}

Complétez les répliques du client avec les formules fonctionnelles étudiées (\eng{I would like...}, \eng{Could you tell me...?}, \eng{Thank you}). L'employé est déjà écrit ; à vous de faire vivre le client.

\begin{exoresolu}[Dialogue à trous — au guichet]
Énoncé. Complétez les répliques du client : Clerk: Good morning, Sir. Can I help you? — Customer: (1) ............ (saluer + demander un retrait de 20,000 francs). — Clerk: Certainly. May I see your ID card? — Customer: (2) ............ (donner la carte). — Clerk: Please sign here. — Customer: (3) ............ (demander le solde). — Clerk: Your balance is 54,000 francs. — Customer: (4) ............ (remercier et conclure).
\tcblower
Solution détaillée. (1) \eng{Good morning. I would like to withdraw 20,000 francs from my account, please.} — salutation + \cle{I would like} pour la politesse. (2) \eng{Here you are.} — formule fixe pour tendre un document. (3) \eng{Could you tell me my balance, please?} — demande d'information polie avec \cle{Could you}. (4) \eng{Thank you very much. Goodbye!} — remerciement + conclusion. Quatre répliques, quatre fonctions : c'est la structure minimale d'un échange réussi.
\end{exoresolu}

\courssub{Jeu de rôle 2 (semi-guidé) : cartes de rôle}

Travaillez par deux. Chaque élève lit sa carte, puis le dialogue se construit librement à partir des consignes. Ne rédigez pas le dialogue à l'avance : parlez !

\begin{experience}[Cartes de rôle — Opening a savings account]
\textbf{Carte A — Client.} You want to open a savings account. Your budget is 25,000 francs. Greet the clerk, explain what you want, ask at least two questions (interest rate? minimum deposit? documents?), react to the answers, thank the clerk and conclude.

\textbf{Carte B — Bank clerk.} A customer wants to open a savings account. Greet the customer and give the conditions: minimum deposit 10,000 francs, interest 3\% per year, no monthly fees. Ask for the required documents: ID card, passport photo, proof of address. Answer the customer's questions politely and conclude.
\end{experience}

\textbf{Glossaire.} \emph{savings account} : compte d'épargne ; \emph{interest rate} : taux d'intérêt ; \emph{minimum deposit} : dépôt minimum ; \emph{monthly fees} : frais mensuels ; \emph{proof of address} : justificatif de domicile.

\textbf{Exemple de début possible.} \eng{Good afternoon. I would like to open a savings account, please.} — \eng{Certainly, Sir. The minimum deposit is 10,000 francs.} — \eng{I see. And what is the interest rate?} — \eng{It is 3\% per year.} — \eng{Oh, really? That's good. What documents do I need?}

\courssub{Jeu de rôle 3 (libre) : au distributeur ou réclamation au guichet}

Sans aucune carte, choisissez l'une des deux situations et improvisez un dialogue d'au moins huit répliques :

\begin{itemize}
\item \textbf{Situation 1.} The cash machine (ATM) kept your card. Go to the counter, explain the problem, answer the clerk's questions (ID card, account number) and ask when you can get your card back.
\item \textbf{Situation 2.} You withdrew 50,000 francs but the machine gave you only 40,000 francs. Make a complaint at the counter politely: \eng{Excuse me, there is a problem...}
\end{itemize}

\begin{methode}[Réussir un jeu de rôle au Bac : 4 étapes = 4 points faciles]
\begin{enumerate}
\item \textbf{Salue et présente ta demande} : \eng{Good morning. I would like to open an account.}
\item \textbf{Pose au moins deux questions} : \eng{How much does it cost? What documents do I need?}
\item \textbf{Relance si tu ne comprends pas} : \eng{Sorry, could you repeat that, please? What do you mean by “fees”?}
\item \textbf{Remercie et conclus} : \eng{Thank you very much for your help. Goodbye!}
\end{enumerate}
Chaque étape visible rapporte des points dans la grille de l'examinateur : ne sautez jamais la salutation ni la conclusion.
\end{methode}

\courssub{Grille d'auto-évaluation}

Après chaque jeu de rôle, évaluez-vous honnêtement avec cette grille. Cochez « oui » seulement si vous l'avez vraiment fait à voix haute.

\begin{center}
\begin{tabularx}{\linewidth}{|X|c|c|}\hline
\rowcolor{popblueL} Critère & Oui & Pas encore \\ \hline
Ai-je salué mon partenaire (\eng{Good morning / Good afternoon}) ? & & \\ \hline
Ai-je présenté clairement ma demande (\eng{I would like...}) ? & & \\ \hline
Ai-je posé au moins deux questions (\eng{How much...? Can I...?}) ? & & \\ \hline
Ai-je réagi aux réponses (\eng{Oh, really? That's fine!}) ? & & \\ \hline
Ai-je relancé quand je ne comprenais pas (\eng{Could you repeat, please?}) ? & & \\ \hline
Ai-je remercié et conclu (\eng{Thank you. Goodbye!}) ? & & \\ \hline
Ai-je bien prononcé \eng{withdrawal}, \eng{balance}, \eng{account} ? & & \\ \hline
\end{tabularx}
\end{center}

Objectif : sept « oui ». Si un critère manque, rejouez la scène en vous concentrant uniquement sur ce point.

\begin{exercice}[À vous de jouer]
Par groupes de deux, improvisez ce dialogue : un élève de Terminale va à la banque de Buea pour demander comment payer ses frais d'inscription à l'université par virement (\eng{transfer}). Le client pose au moins trois questions ; l'employé explique les frais et le délai. Terminez par la grille d'auto-évaluation.
\end{exercice}

\begin{corrige}[Exercice]
Corrigé de méthode : il n'existe pas une seule bonne réponse, mais un bon dialogue contient obligatoirement : une salutation (\eng{Good afternoon}), une demande claire (\eng{I would like to pay my university fees by transfer}), au moins trois questions (\eng{How much are the fees? How long does the transfer take? What information do you need?}), des réactions (\eng{Oh, really? Two days? That's fine.}), un remerciement et une conclusion (\eng{Thank you for your help. Goodbye!}). Vérifiez chaque point avec la grille d'auto-évaluation et recommencez si un élément manque.
\end{corrige}

\begin{savaistu}[Un complément utile pour le Bac]
Les sujets d'expression orale du Baccalauréat portent très souvent sur des situations de la vie économique quotidienne : la banque, le marché, la poste, l'agence de voyage, la station-service. Si vous maîtrisez le dialogue bancaire — saluer, demander, relancer, conclure — vous savez déjà gérer la plupart de ces sujets, car la structure de l'échange reste exactement la même. Entraînez-vous à transposer : \eng{I would like to send a parcel} (poste), \eng{I would like to book a ticket to Lagos} (agence de voyage).
\end{savaistu}

\coursec{Entraînement au Bac et production finale}

\courssub{Transition et objectifs de la séance}

Après avoir travaillé les \textbf{dialogues modèles} et pratiqué des \textbf{jeux de rôle guidés} (section précédente), vous disposez maintenant du vocabulaire, des formules fonctionnelles et de la prononciation nécessaires pour interagir dans une situation bancaire authentique. Cette dernière section a pour but de vous mettre en condition d'examen : l'épreuve orale du Baccalauréat évalue votre capacité à \textbf{communiquer}, non à réciter. Nous allons donc simuler un sujet type, organiser une évaluation par les pairs à l'aide d'une grille précise, et corriger collectivement les erreurs récurrentes pour ancrer les bons réflexes.

\begin{methode}[Gérer son stress à l'oral — 3 réflexes clés]
\begin{enumerate}
    \item \textbf{Respirez} : inspirez profondément par le nez 3 secondes, expirez par la bouche 5 secondes. Cela ralentit le rythme cardiaque et pose la voix.
    \item \textbf{Parlez lentement} : un débit maîtrisé masque l'hésitation et laisse le temps à votre cerveau de trouver les mots.
    \item \textbf{Utilisez des \cle{fillers} (mots-outils) pour gagner du temps sans silence} : \eng{Well...}, \eng{Let me see...}, \eng{Actually...}, \eng{That's a good question...}. Ces marqueurs montrent que vous gérez l'interaction, ce qui est noté positivement.
\end{enumerate}
\end{methode}

\courssub{Critères d'évaluation de l'oral au Baccalauréat}

L'examinateur note sur une grille officielle. Comprendre ces critères, c'est savoir où placer vos efforts. Voici les quatre piliers, déclinés pour notre situation bancaire :

\begin{propriete}[Grille d'évaluation officielle — Adaptée « Services bancaires »]
\begin{center}
\begin{tabularx}{\linewidth}{|X|X|X|}
\hline
\rowcolor{popblueL}
\textbf{Critère} & \textbf{Indicateurs de réussite (Niveau A2/B1 visé)} & \textbf{Points forts attendus} \\
\hline
\cle{Aisance / Fluidité} & Parle sans pauses trop longues ; utilise des connecteurs logiques (\eng{First, Then, Also, However}) ; enchaîne les tours de parole naturellement. & Pas de « blanc » de plus de 5 secondes ; débit régulier. \\
\hline
\cle{Exactitude / Correction} & Structures grammaticales correctes (présent, conditionnel \eng{would like}, modaux \eng{can/could}, prépositions) ; vocabulaire bancaire précis (\eng{account, charges, deposit, withdrawal, statement}). & Erreurs mineures non bloquantes ; auto-correction possible. \\
\hline
\cle{Prononciation / Intonation} & Accent tonique sur les mots forts ; intonation montante pour les questions \eng{yes/no}, descendante pour \eng{wh-} ; \cle{th} / \cle{r} / \cle{h} aspiré corrects. & Compréhensible du premier coup ; rythme anglais (accentuel). \\
\hline
\cle{Interaction / Stratégies} & Salue, remercie ; pose des questions (\eng{Could you tell me...?}) ; \textbf{relance} en cas de non-compréhension (\eng{Sorry? Could you repeat?}) ; conclut poliment. & Preuve de compétence communicative : on ne subit pas, on gère l'échange. \\
\hline
\end{tabularx}
\end{center}
\end{propriete}

\begin{popfigure}
\begin{tikzpicture}[
    cell/.style={rectangle, draw=popblue, thick, rounded corners, minimum height=1.8cm, align=left, fill=popblue!5, font=\small, text width=3.8cm},
    header/.style={rectangle, draw=popdark, thick, rounded corners, minimum height=1.2cm, align=center, fill=popblue!30, font=\small\bfseries, text width=3.8cm},
    crit/.style={rectangle, draw=popdark, thick, rounded corners, minimum height=1.8cm, align=center, fill=poppurple!15, font=\small\bfseries, text width=2.5cm},
    score/.style={rectangle, draw=poporange, thick, rounded corners, minimum height=1.8cm, align=left, fill=poporange!15, font=\small\bfseries, text width=4.2cm}
]
% Header row
\node[header] (h1) at (0,0) {Critère};
\node[header, align=center] (h2) at (3.3,0){0 point \\ (Absent / Insuffisant)};
\node[header, align=center] (h3) at (7.1,0){1 point \\ (Partiel / Perfectible)};
\node[header, align=center] (h4) at (10.9,0){2 points \\ (Maîtrisé / Réussi)};

% Data rows
\foreach \crit/\zero/\un/\deux [count=\i from 1] in {
    Politesse / Pas de salutation ni remerciement / Salutation OU remerciement seul / Salutation + Formules de politesse + Remerciement,
    Questions / Aucune question posée / 1 question simple (ex: How much?) / 2-3 questions variées (charges, documents, carte),
    Relances / Silence face à l'incompréhension / "Pardon?" ou "What?" (familier) / "Could you repeat please?" / "Sorry, I didn't catch that.",
    Prononciation / Incompréhensible (accent français très fort) / Compréhensible avec effort / Claire, accentuation correcte, "th" / "h" ok,
    Conclusion / Part sans dire au revoir / "Bye" seul / "Thank you, goodbye. Have a nice day."
}{
    \pgfmathsetmacro{\y}{-\i*2.3};
    \node[crit] (r\i) at (0,\y) {\crit};
    \node[cell] at (3.3,\y) {\zero};
    \node[cell] at (7.1,\y) {\un};
    \node[score] at (10.9,\y) {\deux};
}
\end{tikzpicture}
\legende{Grille d'évaluation par les pairs — Critères × Niveaux (0 à 2 points). Total sur 10.}
\end{popfigure}

\courssub{Simulation d'examen : Sujet type}

\begin{textebox}[Sujet d'examen simulé — Afriland Bank]
\textbf{Instructions :} “You are a customer at Afriland Bank, Yaoundé. You want to open a savings account. Ask the bank clerk about :
\begin{itemize}
    \item the documents required,
    \item the monthly bank charges,
    \item the possibility of getting a debit card.
\end{itemize}
The examiner plays the role of the bank clerk. Start the conversation.”
\end{textebox}

\begin{methode}[Préparation en binôme — 2 minutes chrono]
\begin{enumerate}
    \item \textbf{Répartissez les rôles} : Élève A = Client (candidat), Élève B = Employé de banque (examinateur). Changez de rôle au second passage.
    \item \textbf{Notez 3 formules « passe-partout »} à placer impérativement dans votre dialogue (ex. : \eng{I would like to...}, \eng{Could you tell me...?}, \eng{Thank you for your help.}).
    \item \textbf{Anticipez les réponses de l'employé} : il va vous demander votre pièce d'identité, un justificatif de domicile, un dépôt initial. Préparez des réponses courtes : \eng{Here is my ID card.}, \eng{Yes, I have a utility bill.}, \eng{How much is the initial deposit?}
    \item \textbf{Répétez mentalement} la structure : Salutation $\rightarrow$ Demande polie $\rightarrow$ 3 questions $\rightarrow$ Relance si besoin $\rightarrow$ Conclusion.
\end{enumerate}
\end{methode}

\begin{exemplebox}[Modèle de production attendue (Niveau B1)]
\begin{tabular}{@{}p{0.15\linewidth}p{0.8\linewidth}@{}}
\textbf{Client} & “Good morning, Sir. \eng{I would like to open a savings account.}” (Salutation + Demande polie) \\
\textbf{Clerk} & “Good morning. Certainly. \eng{Do you have your ID card and a proof of address?}” \\
\textbf{Client} & “Yes, here they are. \eng{Could you tell me what the monthly charges are, please?}” (Question 1) \\
\textbf{Clerk} & “\eng{It's 500 CFA francs per month for a standard savings account.}” \\
\textbf{Client} & “\eng{Let me see...} \eng{And do I get a debit card with this account?}” (Question 2) \\
\textbf{Clerk} & “\eng{Yes, a Visa debit card is available for 2,000 francs issuance fee.}” \\
\textbf{Client} & “\eng{Okay. What is the minimum initial deposit?}” (Question 3) \\
\textbf{Clerk} & “\eng{10,000 francs.}” \\
\textbf{Client} & “\eng{Well, that sounds fine.} \eng{Thank you very much for your help. Goodbye, Sir.}” (Conclusion)
\end{tabular}
\end{exemplebox}

\begin{propriete}[Structure d'un échange réussi — Rappel synthétique]
Tout échange oral noté au Bac doit respecter cette architecture en 5 étapes. L'omission d'une étape fait perdre des points sur \textbf{Interaction} et \textbf{Aisance}.
\begin{center}
\begin{tabular}{|c|l|l|}
\hline
\rowcolor{popblueL}
\textbf{Étape} & \textbf{Fonction} & \textbf{Exemples de formules} \\
\hline
1 & \cle{Salutation} & \eng{Good morning, Sir/Madam.} \\
2 & \cle{Demande polie (Objectif)} & \eng{I would like to... / I want to...} \\
3 & \cle{2-3 Questions / Réponses} & \eng{Could you tell me...? / How much...? / What documents...?} \\
4 & \cle{Relance éventuelle} & \eng{Sorry, could you repeat? / Did you say 500 or 5,000?} \\
5 & \cle{Remerciement + Conclusion} & \eng{Thank you for your help. Goodbye.} \\
\hline
\end{tabular}
\end{center}
\end{propriete}

\courssub{Passage devant la classe \& Évaluation par les pairs}

\begin{experience}[Simulation Bac — Passage à l'oral]
\begin{enumerate}
    \item \textbf{Installation} : Un binôme vient au tableau. L'enseignant (ou un élève désigné) joue le rôle de l'examinateur / employé de banque pour standardiser l'interaction.
    \item \textbf{Observation} : Le reste de la classe remplit la \textbf{grille d'évaluation} (figure ci-dessus) pour le candidat. Chaque critère reçoit 0, 1 ou 2 points. Notez aussi une \textbf{observation qualitative} (ex. : « Bon "th", mais a oublié de saluer »).
    \item \textbf{Durée} : 1 minute 30 maximum par passage.
    \item \textbf{Rotation} : 4 à 5 binômes passent (selon l'effectif) pour couvrir les cas variés : compte courant, compte épargne, virement, carte bloquée.
\end{enumerate}
\end{experience}

\courssub{Correction collective : Reformulation des erreurs fréquentes}

Après les passages, l'enseignant mène une correction au tableau en reprenant les erreurs entendues. C'est le moment de \textbf{reformuler} (donner la forme correcte) sans humilier.

\begin{attention}[Erreur fréquente — Ne pas réciter un dialogue appris par cœur]
\begin{itemize}
    \item \textbf{Piège} : L'élève débite un texte mémorisé sans écouter la question de l'examinateur. Examinateur : “\eng{Do you have an ID card?}” — Élève : “\eng{I would like to open a savings account.}” (Hors sujet total).
    \item \textbf{Conséquence} : Note 0 en \textbf{Interaction} et \textbf{Aisance}.
    \item \textbf{Parade} : \textbf{Écoutez le mot-clé} dans la question (\eng{ID card, charges, card, deposit}) et y répondre directement avant d'enchaîner votre question préparée. Utilisez \eng{Yes, I do. / Here it is. / Let me check.}
\end{itemize}
\end{attention}

\begin{exoresolu}[Correction d'énoncés d'élèves — Reformulation guidée]
\textbf{Consigne :} Voici 5 productions d'élèves entendues lors des passages. Identifiez l'erreur (grammaire, vocabulaire, prononciation, interaction) et proposez la reformulation correcte.

\begin{enumerate}
\item “\eng{Good morning. I want open a account.}”
\item “\eng{How much the charges?}”
\item “\eng{I want a card for retire money.}”
\item “\eng{What?}” (après une réponse de l'employé non comprise)
\item “\eng{Okay. Bye.}” (Départ brutal)
\end{enumerate}
\tcblower
\textbf{Solutions :}
\begin{enumerate}
\item \textbf{Erreurs} : \eng{want} + infinitif sans \eng{to} (incorrect) ; \eng{a account} (article). \\ \textbf{Reformulation} : “\eng{Good morning. \textbf{I would like to open \underline{an} account.}}” (Politesse + grammaire).
\item \textbf{Erreur} : Absence d'auxiliaire / inversion (structure question). \\ \textbf{Reformulation} : “\eng{\textbf{How much \underline{are} the charges?} / \textbf{Could you tell me \underline{how much the charges \underline{are}}?}}” (Style indirect plus poli).
\item \textbf{Erreurs} : \eng{retire} (verbe) au lieu de \eng{withdraw} ; \eng{for} + V-ing (but here purpose $\rightarrow$ \eng{to withdraw}). \eng{Retire} = prendre sa retraite (faux ami). \\ \textbf{Reformulation} : “\eng{\textbf{I would like a debit card \underline{to withdraw money.}}}”
\item \textbf{Erreur} : \eng{What?} seul = familier, voire impoli. \\ \textbf{Reformulation} : “\eng{\textbf{Sorry, could you repeat that, please?} / \textbf{I didn't catch the amount.}}”
\item \textbf{Erreur} : Absence de remerciement + formule de politesse de clôture. Perte de 2 points sur Interaction. \\ \textbf{Reformulation} : “\eng{\textbf{Thank you very much for your help. Goodbye, Sir. Have a nice day.}}”
\end{enumerate}
\end{exoresolu}

\begin{exemplebox}[Exemples traduits — Formules de relance et de gain de temps]
\begin{tabular}{@{}p{0.45\linewidth}p{0.5\linewidth}@{}}
“\eng{Let me see... I would like a savings account.}” & \textit{Voyons... je voudrais un compte d'épargne.} \\
“\eng{Well, the charges are 500 francs per month.}” & \textit{Eh bien, les frais sont de 500 francs par mois.} \\
“\eng{Actually, I need a debit card for withdrawals.}” & \textit{En fait, j'ai besoin d'une carte de débit pour les retraits.} \\
“\eng{Sorry, could you speak more slowly?}” & \textit{Pardon, pourriez-vous parler plus lentement ?} \\
“\eng{Did you say the initial deposit is 10,000?}” & \textit{Avez-vous dit que le dépôt initial est de 10 000 ?}
\end{tabular}
\end{exemplebox}

\courssub{Bilan de la leçon : Auto-évaluation}

\begin{aretenir}[Bilan — Je sais maintenant...]
\begin{itemize}
    \item \textbf{Demander} un service bancaire poliment : \eng{I would like to... / Could I...?}
    \item \textbf{Informer} sur mes besoins / documents : \eng{Here is my ID. / I have a proof of address.}
    \item \textbf{Questionner} sur les frais, cartes, conditions : \eng{How much are the charges? / Is there a debit card?}
    \item \textbf{Relancer} en cas de problème : \eng{Sorry, could you repeat? / I didn't catch that.}
    \item \textbf{Conclure} l'échange : \eng{Thank you for your help. Goodbye.}
    \item \textbf{Prononcer} correctement : \eng{account} \textit{(a-kaount)}, \eng{withdrawal} \textit{(wi-zro-al)}, \eng{receipt} \textit{(ri-sit)}, \eng{Visa} \textit{(vi-za)}.
\end{itemize}
\end{aretenir}

\begin{savaistu}[L'astuce qui fait la différence au Bac]
À l'oral du Baccalauréat, \textbf{utiliser une relance naturelle} (\eng{Sorry, could you repeat that?} ou \eng{Excuse me, I didn't catch the last word.}) est \textbf{mieux noté qu'un long silence}. Pourquoi ? Parce que le silence montre une rupture de communication, alors que la relance prouve que vous possédez la \textbf{compétence stratégique} de réparer la conversation. C'est un indicateur fort de niveau B1. N'ayez pas peur de dire « Pardon » en anglais, c'est professionnel !
\end{savaistu}

\begin{exercice}[Entraînement individuel — Préparation Bac (Devoir maison)]
\begin{enumerate}
    \item Rédigez, en anglais, un dialogue complet (10-12 répliques) pour ce sujet : \textit{« You are at UBA Bank. You want to transfer 50,000 FCFA to your sister's account in Douala. Ask about fees, delay, and required info. »}
    \item Soulignez dans votre dialogue : les 5 étapes de la structure (Salutation, Demande, Questions, Relance, Conclusion).
    \item Enregistrez-vous (téléphone) en jouant les deux rôles. Réécoutez : entendez-vous vos \eng{th}, \eng{h}, l'accent tonique ? Notez 2 points à améliorer.
\end{enumerate}
\end{exercice}

\begin{corrige}[Exercice — Éléments de réponse]
\begin{enumerate}
    \item \textbf{Dialogue type} :
    “\eng{Client: Good morning. I would like to make a transfer.} \\
    \eng{Clerk: Good morning. Where to?} \\
    \eng{Client: To my sister's account in Douala. The amount is 50,000 francs.} \\
    \eng{Clerk: Okay. Do you have her account number and the SWIFT/BIC code?} \\
    \eng{Client: Let me see... Yes, I have the account number. What is the SWIFT code?} \\
    \eng{Clerk: It's UBAFCMCD for UBA Douala.} \\
    \eng{Client: \textbf{Sorry, could you spell that?}} (Relance) \\
    \eng{Clerk: U-B-A-F-C-M-C-D.} \\
    \eng{Client: Thank you. \textbf{How much are the transfer fees?}} (Question) \\
    \eng{Clerk: 1,500 francs commission + 500 francs stamp duty.} \\
    \eng{Client: \textbf{And how long does it take?}} (Question) \\
    \eng{Clerk: Usually 24 hours.} \\
    \eng{Client: Perfect. \textbf{Thank you very much for your help. Goodbye.}}” (Conclusion)
    \item Structure respectée : 1-Salutation, 2-Demande, 3-3 questions (SWIFT, frais, délai), 4-Relance (épellation), 5-Conclusion.
    \item Auto-correction personnelle attendue sur : \eng{transfer} \textit{(trans-fer)}, \eng{commission} \textit{(ko-mi-chon)}, \eng{usually} \textit{(you-jou-a-li)}.
\end{enumerate}
\end{corrige}

\newpage

\coursec{Activité d'intégration}

\coursec{Lesson 16 — Speaking: Exchanges information about using banking services}

\courssub{Découverte : A dialogue at the bank (Listening / Reading)}

\begin{textebox}[Listening script — At the counter]
\eng{— Good morning, Madam. Welcome to Buea Community Bank. How can I help you today?}\\
\eng{— Good morning. I would like to withdraw some money from my current account, please.}\\
\eng{— Certainly. Could I have your account number and your ID card?}\\
\eng{— Here is my ID card. My account number is 44029188.}\\
\eng{— Thank you. How much would you like to withdraw?}\\
\eng{— 50,000 FCFA, please.}\\
\eng{— Right. Do you need a receipt?}\\
\eng{— Yes, please. And could you tell me my current balance?}\\
\eng{— Of course. Your current balance is 125,500 FCFA.}\\
\eng{— Sorry, did you say one hundred twenty-five thousand five hundred?}\\
\eng{— That is correct. One hundred and twenty-five thousand five hundred FCFA.}\\
\eng{— Thank you very much. Goodbye.}\\
\eng{— You are welcome, Madam. Have a nice day.}
\end{textebox}

\begin{exercice}[Compréhension orale / écrite]
Répondez en anglais aux questions suivantes d'après le dialogue.
\begin{enumerate}
\item What service does the customer want? 
\item What two documents does the teller ask for? 
\item How much money does the customer withdraw? 
\item What is the customer's balance after the withdrawal? 
\item Which \cle{relance} (clarification strategy) does the customer use? 
\end{enumerate}
\end{exercice}

\begin{corrige}[Exercice 1]
\begin{enumerate}
\item She wants to \textbf{withdraw money} from her current account.
\item The teller asks for her \textbf{account number} and her \textbf{ID card}.
\item She withdraws \textbf{50,000 FCFA}.
\item Her balance is \textbf{125,500 FCFA} (one hundred and twenty-five thousand five hundred).
\item She says: \eng{“Sorry, did you say one hundred twenty-five thousand five hundred?”} (vérification des chiffres).
\end{enumerate}
\end{corrige}

\courssub{Lexique des services bancaires (Vocabulary)}

\begin{propriete}[Mots clés — Banking Services]
\begin{center}
\begin{tabularx}{\linewidth}{|X|X|X|}
\hline
\rowcolor{popblueL}
\textbf{English} & \textbf{Français} & \textbf{Remarque / Contexte} \\
\hline
\eng{savings account} & compte épargne & \eng{Also called \emph{deposit account}.} \\
\eng{current account} & compte courant & \eng{UK English ; US : \emph{checking account}.} \\
\eng{to open / close an account} & ouvrir / fermer un compte & \eng{Open an account \textbf{with} a deposit.} \\
\eng{deposit (n/v)} & dépôt / déposer & \eng{Make a deposit / Deposit money.} \\
\eng{withdrawal (n) / withdraw (v)} & retrait / retirer & \eng{Irregular verb: withdraw – withdrew – withdrawn.} \\
\eng{balance} & solde & \eng{Check / ask for \textbf{my balance}.} \\
\eng{bank card / debit card} & carte bancaire / carte de débit & \eng{Credit card = carte de crédit.} \\
\eng{fee / charge} & frais (bancaires) & \eng{Annual fee, transaction fee.} \\
\eng{interest rate} & taux d'intérêt & \eng{3\% per annum (p.a.).} \\
\eng{proof of address} & justificatif de domicile & \eng{Electricity bill, rent receipt.} \\
\eng{passport photo} & photo d'identité & \eng{Recent, white background.} \\
\eng{counter / teller} & guichet / guichetier & \eng{Queue at the counter.} \\
\eng{receipt} & reçu & \eng{Keep the receipt.} \\
\eng{form / application form} & formulaire / demande & \eng{Fill in / complete a form.} \\
\hline
\end{tabularx}
\end{center}
\end{propriete}

\begin{aretenir}[Faux amis et pièges lexicaux]
\begin{itemize}
\item \cle{Actually} = \eng{en fait, en réalité} (pas « actuellement » $\rightarrow$ \eng{at the moment / currently}).
\item \cle{Balance} = \eng{solde} (pas \eng{sold} $\rightarrow$ vendu / \eng{sale}).
\item \cle{Remove} = \eng{enlever, retirer (un objet)} ; pour l'argent : \cle{withdraw}.
\item \cle{Agency} = \eng{agence (commerciale, publicité)} ; pour la banque : \cle{branch}.
\item \cle{Signature} = \eng{signature} ; \eng{sign} (verbe).
\end{itemize}
\end{aretenir}

\courssub{Formules fonctionnelles (Functional Language)}

\begin{propriete}[Structure de l'échange bancaire]
\begin{center}
\begin{tabularx}{\linewidth}{|l|X|X|}
\hline
\rowcolor{popblueL}
\textbf{Étape} & \textbf{Formules types (Employee $\leftrightarrow$ Customer)} & \textbf{Fonction} \\
\hline
\textbf{1. Accueil} & \eng{Good morning. Welcome to [Bank Name]. How can I help you?} & \textit{Ouvrir l'échange, offrir assistance} \\
& \eng{Good morning. I would like to... / I need to...} & \textit{Exprimer son besoin} \\
\hline
\textbf{2. Identification} & \eng{Could I have your account number / ID card, please?} & \textit{Demander pièces justificatives} \\
& \eng{Here is my... / My account number is...} & \textit{Fournir informations} \\
\hline
\textbf{3. Demande d'infos} & \eng{How much would you like to withdraw / deposit?} & \textit{Préciser montant (Employé)} \\
& \eng{What documents are required? / How long does it take?} & \textit{Questionner conditions (Client)} \\
& \eng{Could you tell me the fee / interest rate?} & \textit{Question indirecte polie} \\
\hline
\textbf{4. Donne d'infos} & \eng{The minimum deposit is... / The fee is...} & \textit{Énoncer faits, chiffres} \\
& \eng{It will be ready in one week / immediately.} & \textit{Indiquer délai} \\
\hline
\textbf{5. Relance / Vérification} & \eng{Sorry, could you repeat that, please?} & \textit{Demander répétition} \\
& \eng{Did you say [X] or [Y]?} & \textit{Vérifier chiffre / mot} \\
& \eng{What do you mean by “...”?} & \textit{Demander clarification} \\
\hline
\textbf{6. Conclusion} & \eng{Thank you very much for your help. Goodbye.} & \textit{Remercier, clore} \\
& \eng{You're welcome. Have a nice day!} & \textit{Répondre, saluer} \\
\hline
\end{tabularx}
\end{center}
\end{propriete}

\begin{methode}[Comment formuler une question polie indirecte]
\begin{enumerate}
\item Repérez la question directe : \eng{What is the fee?} / \eng{How long does it take?}
\item Ajoutez une introduction polie : \eng{Could you tell me...} / \eng{Do you know...} / \eng{I would like to know...}
\item \textbf{Inversez l'ordre sujet-verbe} (supprimez l'auxiliaire \emph{do/does/did}) : \eng{...what the fee \textbf{is}.} / \eng{...how long it \textbf{takes}.}
\item Terminez par un point (pas de point d'interrogation si la phrase principale est affirmative : \eng{I would like to know...}).
\end{enumerate}
\end{methode}

\begin{exemplebox}[Application]
\eng{Direct:} \eng{How much does the card cost?}\\
\eng{Indirect:} \eng{Could you tell me how much the card \textbf{costs}?} (pas \eng{does cost})
\end{exemplebox}


\courssub{Prononciation, accentuation et intonation}

\begin{propriete}[Règles d'accentuation tonique (Word Stress)]
\begin{itemize}
\item \textbf{Noms / Adjectifs 2 syllabes :} accent sur la \textbf{1ère} syllabe. 
\eng{BAnking}, \eng{CUStomer}, \eng{DOcument}, \eng{BAlance}, \eng{REcEIPT}, \eng{CHEqUE}.
\item \textbf{Verbes 2 syllabes (préfixe + racine) :} accent sur la \textbf{2ème} syllabe. 
\eng{dePOsit}, \eng{withDRAW}, \eng{enQUIRE}, \eng{reCEIPT} (verbe rare), \eng{reFUND}.
\item \textbf{Mots composés :} accent sur le \textbf{premier élément}. 
\eng{BANK account}, \eng{CREDIT card}, \eng{PASSport photo}, \eng{CURrent account}.
\item \textbf{Terminaison \emph{-tion} / \emph{-sion} :} accent sur la syllabe \textbf{avant} la terminaison. 
\eng{withdraw\underline{A}tion}, \eng{depos\underline{I}tion}, \eng{applic\underline{A}tion}.
\end{itemize}
\end{propriete}

\begin{popfigure}
\begin{tikzpicture}[
    node distance=0.8cm and 1.2cm,
    box/.style={draw=popblue, thick, rounded corners, fill=popblueL, text width=2.8cm, align=center, minimum height=1.2cm},
    arrow/.style={-Stealth, thick, popdark}
]
\node[box, align=center] (n1){\textbf{1ère syllabe}\\\eng{BAnking}\\\eng{CUStomer}\\\eng{DOcument}\\\eng{BAlance}};
\node[box, right=of n1, align=center] (n2){\textbf{2ème syllabe}\\\eng{dePOsit}\\\eng{withDRAW}\\\eng{enQUIRE}\\\eng{appliCATION}};
\node[box, right=of n2, align=center] (n3){\textbf{Composés}\\\eng{BANK account}\\\eng{CREDIT card}\\\eng{PASSport photo}};
\draw[arrow] (n1) -- (n2) node[midway, above, font=\small] {Noms/Adj vs Verbes};
\draw[arrow] (n2) -- (n3) node[midway, above, font=\small] {Règle composé};
\end{tikzpicture}
\legende{Schémas d'accentuation tonique du vocabulaire bancaire}
\end{popfigure}

\begin{propriete}[Intonation des questions (Sentence Stress \& Intonation)]
\begin{center}
\begin{tabularx}{\linewidth}{|X|X|X|}
\hline
\rowcolor{popblueL}
\textbf{Type de question} & \textbf{Intonation} & \textbf{Exemples} \\
\hline
\emph{Yes/No questions} (auxiliaire initial) & \textbf{Montante} (\eng{$\nearrow$}) & \eng{Do you have your ID card? $\nearrow$} \\
& & \eng{Would you like a receipt? $\nearrow$} \\
\hline
\emph{Wh- questions} (mots interrogatifs) & \textbf{Descendante} (\eng{$\searrow$}) & \eng{How much do you want? $\searrow$} \\
& & \eng{What documents do I need? $\searrow$} \\
\hline
\emph{Question de vérification / choix} & \textbf{Montante sur l'alternative} & \eng{Did you say fifteen $\nearrow$ or fifty $\searrow$?} \\
\hline
\end{tabularx}
\end{center}
\end{propriete}

\begin{exercice}[Entraînement à la prononciation]
Marquez l'accent tonique (') sur les mots en gras et tracez la flèche d'intonation ($\nearrow$ ou $\searrow$) pour chaque phrase.
\begin{enumerate}
\item The \textbf{customer} wants to \textbf{deposit} money.
\item \textbf{How much} is the annual \textbf{fee}?
\item \textbf{Do you} need a \textbf{receipt}?
\item \textbf{Did you say} \textbf{fifteen} or \textbf{fifty}?
\end{enumerate}
\end{exercice}

\begin{corrige}[Exercice 2]
\begin{enumerate}
\item The \textbf{'customer} wants to de\textbf{'posit} money. (Déclarative $\searrow$)
\item \textbf{'How much} is the annual \textbf{'fee}? (Wh- $\searrow$)
\item \textbf{Do you} need a re\textbf{'ceipt}? (Yes/No $\nearrow$)
\item \textbf{Did you say} \textbf{'fifteen} $\nearrow$ or \textbf{'fifty} $\searrow$? (Choix)
\end{enumerate}
\end{corrige}

\courssub{Dialogues modèles et jeux de rôle guidés (Guided Practice)}

\courssub{Situation de communication}

\textbf{Contexte.} La classe est aménagée en agence bancaire : un bureau fait office de \cle{guichet} (\eng{counter}), une rangée de chaises représente la \cle{file d'attente} (\eng{queue}). Nous sommes à la \eng{Buea Community Bank}, une petite agence au pied du mont Cameroun. Un client entre pour accomplir une mission précise ; l'employé de banque dispose d'une fiche d'informations qu'il doit communiquer avec exactitude et politesse.

\begin{textebox}[Buea Community Bank — Notice at the entrance]
Welcome to Buea Community Bank.\\
Opening hours: Monday to Friday, 8 a.m. to 3.30 p.m.\\
Services available at the counter: opening savings and current accounts, deposits, withdrawals, balance enquiries, bank cards.\\
Please take a ticket and wait for your number to be called.\\
Our staff will be happy to help you in English or in French.\\
\emph{(Bienvenue à la Buea Community Bank. Horaires : du lundi au vendredi, de 8 h à 15 h 30. Services au guichet : ouverture de comptes épargne et courants, dépôts, retraits, demandes de solde, cartes bancaires. Veuillez prendre un ticket et attendre que votre numéro soit appelé.)}
\end{textebox}

\textbf{Documents de jeu.} L'enseignant prépare deux types de cartes :

\begin{itemize}
\item \textbf{Cartes « Client »} (la mission) : 
\begin{itemize}
\item \eng{Open a savings account with 25,000 FCFA.}
\item \eng{Withdraw 50,000 FCFA and ask for your balance.}
\item \eng{Ask about the cost and conditions of a bank card.}
\item \eng{Deposit 100,000 FCFA into your current account.}
\end{itemize}
\item \textbf{Cartes « Employee »} (la fiche d'informations) : 
\begin{itemize}
\item \eng{Savings account: Min. deposit 10,000 FCFA. ID card + 1 photo. Book issued immediately. Interest: 3\% p.a.}
\item \eng{Withdrawal: Max 200,000 FCFA/day. ID required. Receipt free.}
\item \eng{Bank card: Cost 5,000 FCFA/year. Ready in 1 week. Needs ID, proof of address, 2 photos.}
\item \eng{Deposit: No limit. Counted at counter. Receipt issued.}
\end{itemize}
\end{itemize}

\courssub{Consignes de l'activité intégrée}

\begin{enumerate}
\item \textbf{Préparation du lexique (5 min).} En binôme, relisez la notice et retrouvez l'équivalent français de : \eng{savings account}, \eng{withdrawal}, \eng{balance enquiry}, \eng{counter}, \eng{opening hours}.
\item \textbf{Tirage au sort des rôles.} L'enseignant distribue une carte « Client » et une carte « Employee » par binôme. Vous avez 3 minutes pour préparer mentalement vos répliques \emph{sans écrire le dialogue en entier} : notez seulement 3 questions du client et 3 informations de l'employé.
\item \textbf{Jeu de rôle au guichet.} Le client salue, explique sa mission, pose \textbf{au moins trois questions}, utilise \textbf{au moins une relance} (\eng{Could you repeat that, please?} / \eng{Sorry, did you say fifteen or fifty?}) et \textbf{conclut poliment} (\eng{Thank you very much for your help. Goodbye.}). L'employé accueille, donne les informations de sa fiche, vérifie la compréhension (\eng{Is that clear?} / \eng{Do you have any other questions?}).
\item \textbf{Évaluation par les pairs.} Pendant chaque scène, les observateurs remplissent la grille :
\begin{center}
\begin{tabular}{|l|c|}
\hline
\textbf{Critère} & \textbf{Points} \\ \hline
Politesse / Registre (Would like, Could you) & /1 \\
Trois questions posées (variété wh-/yes-no) & /2 \\
Relance utilisée correctement & /1 \\
Prononciation (accent tonique) + Intonation & /2 \\
Conclusion rituelle correcte & /1 \\ \hline
\textbf{Total} & \textbf{/7} \\ \hline
\end{tabular}
\end{center}
\item \textbf{Reprise collective.} Les deux meilleurs binômes rejouent leur scène. L'enseignant corrige au tableau : accentuation (\eng{BAnking}, \eng{withDRAwal}, \eng{enQUIry}) ; intonation (montante $\nearrow$ fermé, descendante $\searrow$ ouvert) ; auxiliaires (\eng{How much \textbf{does} it cost?}).
\item \textbf{Trace écrite.} Recopiez le dialogue modèle corrigé et surlignez les formules fonctionnelles.
\end{enumerate}

\courssub{Ressources à mobiliser (Boîte à outils)}

\begin{aretenir}[Boîte à outils du dialogue bancaire]
\begin{itemize}
\item \textbf{Saluer et ouvrir :} \eng{Good morning, Madam. How can I help you?} / \eng{Good morning. I would like to open a savings account, please.}
\item \textbf{Demander une information :} \eng{How much money do I need to open the account?} ; \eng{What documents should I bring?} ; \eng{How long will it take?} ; \eng{Could you tell me the annual fee for the card?}
\item \textbf{Donner une information :} \eng{You need a minimum deposit of 25,000 FCFA.} ; \eng{The card costs 5,000 FCFA per year.} ; \eng{It will be ready in one week.}
\item \textbf{Relancer / Vérifier :} \eng{Sorry, could you repeat that, please?} ; \eng{Pardon? Did you say five thousand or fifteen thousand?} ; \eng{What do you mean by “proof of address”?}
\item \textbf{Conclure :} \eng{Thank you very much for your help.} ; \eng{You're welcome. Have a nice day!}
\item \textbf{Prononciation clé :} \eng{'BAnking}, \eng{'CUStomer}, \eng{'DOcument} ; \eng{de'POsit}, \eng{with'DRAWal}, \eng{en'QUIry}. Intonation $\nearrow$ (fermé) vs $\searrow$ (ouvert).
\end{itemize}
\end{aretenir}

\begin{attention}[Pièges fréquents des francophones — \textbf{À éviter absolument}]
\begin{itemize}
\item \textbf{Vocabulaire :} Ne dites pas \eng{I want to \textbf{remove} money} $\rightarrow$ \eng{to \textbf{withdraw} money}.
\item \textbf{Faux ami :} \eng{Actually} = « en fait » (pas « actuellement » $\rightarrow$ \eng{at the moment}). Ex. \eng{At the moment, I have 30,000 FCFA.}
\item \textbf{Lexique :} Le solde = \eng{balance} (« ba-lance »), jamais \eng{sold}.
\item \textbf{Grammaire :} 3ème personne du singulier \textbf{-s} obligatoire : \eng{The card \textbf{costs} 5,000 FCFA.} / \eng{It \textbf{takes} one week.}
\item \textbf{Article :} Pas de \emph{the} superflu : \eng{I want to withdraw \textbf{money}} (général), pas \eng{the money} (spécifique connu).
\item \textbf{Question indirecte :} \eng{Could you tell me how much it \textbf{costs}?} (pas \eng{does cost}).
\item \textbf{Préposition :} \eng{Open an account \textbf{with} 25,000 FCFA.} / \eng{Pay \textbf{by} card / \textbf{in} cash.}
\end{itemize}
\end{attention}

\courssub{Entraînement au Bac et production finale}

\begin{exoresolu}[Dialogue modèle : Opening a savings account — Corrigé commenté]
Énoncé : À partir de la carte « Client » \eng{Open a savings account with 25,000 FCFA} et de la fiche « Employee » (dépôt minimum 10,000 FCFA ; pièce d'identité + photo ; carnet remis immédiatement ; intérêt 3 \% par an), produisez le dialogue complet.
\tcblower
\textbf{Dialogue modèle :}

\eng{— Good morning, Sir. Welcome to Buea Community Bank. How can I help you?}\\
\eng{— Good morning. I would like to open a savings account, please.}\\
\eng{— Certainly. How much money do you want to deposit?}\\
\eng{— I have 25,000 FCFA. Is that enough?}\\
\eng{— Yes, it is. The minimum deposit is only 10,000 FCFA.}\\
\eng{— Good. What documents do I need?}\\
\eng{— You need your national ID card and one passport photo.}\\
\eng{— Sorry, could you repeat that, please? A passport photo?}\\
\eng{— That's right. One recent passport photo.}\\
\eng{— All right. And how much interest does the account earn?}\\
\eng{— It earns 3 per cent per year.}\\
\eng{— Did you say three per cent or thirteen per cent?}\\
\eng{— Three per cent per year.}\\
\eng{— Perfect. When will my savings book be ready?}\\
\eng{— You will receive it immediately, as soon as we finish the form.}\\
\eng{— Thank you very much for your help.}\\
\eng{— You're welcome, Sir. Have a nice day!}

\textbf{Commentaire des choix de langue :}
\begin{itemize}
\item \textbf{Ouverture et politesse :} \eng{How can I help you?} et \eng{I would like to...} (plus courtois que \eng{I want}) instaurent le registre formel attendu.
\item \textbf{Trois questions authentiques :} \eng{Is that enough?} (fermée, $\nearrow$), \eng{What documents do I need?} et \eng{How much interest does the account earn?} (ouvertes \emph{wh-}, $\searrow$, auxiliaire \emph{do/does} bien placé).
\item \textbf{Deux relances naturelles :} \eng{Could you repeat that, please?} et \eng{Did you say three per cent or thirteen per cent?} — exploitation de la confusion \emph{-teen / -ty}, stratégie de vérification excellente.
\item \textbf{Informations chiffrées claires :} \eng{The minimum deposit is 10,000 FCFA} ; \eng{3 per cent per year} — phrases courtes (SVO), faciles à prononcer.
\item \textbf{Conclusion rituelle :} \eng{Thank you very much for your help} / \eng{You're welcome}, formules de clôture standard.
\end{itemize}
\end{exoresolu}

\begin{exercice}[Simulation d'épreuve orale du Baccalauréat — \textit{Interaction}]
\textbf{Consigne :} Vous êtes le client. L'examinateur est l'employé de la \eng{Yaoundé Central Bank}. Votre carte mission :
\begin{quote}
\eng{You want to get a bank card. Ask about: cost, required documents, waiting time. You have 15,000 FCFA. Check if it is enough. Use one clarification.}
\end{quote}
Préparez vos 4--5 tours de parole (notes seulement). Jouez la scène avec l'examinateur.
\end{exercice}

\begin{corrige}[Exercice 3 — Proposition de production attendue]
\textbf{Dialogue simulé (Candidate $\leftrightarrow$ Examiner)} \\
\eng{C: Good morning. I would like to get a bank card, please.}\\
\eng{E: Good morning. Certainly. We have debit cards available.}\\
\eng{C: How much does the card cost per year?}\\
\eng{E: The annual fee is 5,000 FCFA.}\\
\eng{C: What documents are required?}\\
\eng{E: We need your ID card, a proof of address and two passport photos.}\\
\eng{C: Sorry, what do you mean by “proof of address”?}\\
\eng{E: An electricity bill or a rent receipt with your name on it.}\\
\eng{C: How long does it take to get the card?}\\
\eng{E: It will be ready in one week.}\\
\eng{C: I have 15,000 FCFA. Is that enough for the first year?}\\
\eng{E: Yes, that covers the fee and leaves 10,000 FCFA on the account.}\\
\eng{C: Perfect. Thank you very much for your help. Goodbye.}\\
\eng{E: You're welcome. Have a nice day!}

\textbf{Points d'évaluation Bac :} Interaction (réactivité), Lexique précis (\eng{proof of address, annual fee}), Grammaire (questions indirectes, auxiliaires), Phonologie (accent \eng{withDRAWal}, intonation $\searrow$ sur \emph{wh-}), Stratégie (relance \eng{what do you mean...}).
\end{corrige}

\courssub{Bilan de la leçon}

Cette leçon a permis de réinvestir, dans une situation de communication authentique et contextualisée (banque camerounaise, FCFA), l'ensemble des acquis : le \cle{lexique} des services bancaires (\eng{savings account}, \eng{deposit}, \eng{withdrawal}, \eng{balance}, \eng{fee}), les \cle{formules fonctionnelles} pour demander, donner une information, relancer et conclure, ainsi que les règles d'\cle{accentuation tonique} et d'\cle{intonation}. La grille d'évaluation par les pairs a développé l'écoute active : pour noter, les observateurs devaient repérer questions, relances et qualité phonologique. La reprise collective a ciblé les erreurs récurrentes : accent tonique (\eng{withDRAwal} vs \eng{WITHdrawal}), omission de l'auxiliaire (\eng{How much it costs?} $\rightarrow$ \eng{How much \textbf{does} it cost?}), confusion \eng{actually} / \eng{at the moment}. Ces points seront consolidés en révision pour l'oral du Baccalauréat où le jeu de rôle fonctionnel est incontournable.

\newpage

\coursec{Méthodes et exercices résolus}

\coursec{Lesson 16 — Speaking: Exchanges information about using banking services}

\courssub{Découverte : un dialogue à la banque}

\begin{textebox}[Dialogue à la banque — Ouverture d'un compte épargne]
\eng{Good morning, madam. How can I help you today?}

\eng{Good morning. I would like to open a savings account, please.}

\eng{Certainly. Do you have your national ID card and two passport photos with you?}

\eng{Yes, I have them here. How much is the minimum deposit?}

\eng{The minimum deposit is 10,000 francs CFA.}

\eng{And what is the interest rate?}

\eng{It is 3.5\% per year. You will also get a free ATM card.}

\eng{That sounds good. Do I get a cheque book with this account?}

\eng{No, cheque books are only available with current accounts. But you can make withdrawals and transfers with your ATM card or via mobile money.}

\eng{I see. Could you repeat the documents required, please?}

\eng{Of course. Your ID card, two passport photos, and a proof of address (electricity bill or tenancy agreement).}

\eng{Thank you very much. I will bring them tomorrow.}

\eng{You are welcome. Have a nice day, madam.}
\end{textebox}

\begin{methode}[Réussir un échange d'informations à la banque]
Pour mener un dialogue bancaire en anglais, suis cette démarche en cinq étapes :
\begin{enumerate}
\item \cle{Saluer et se présenter} : commence par une formule de politesse et dis clairement ce que tu veux : \eng{Good morning, I would like to open a savings account.} « Bonjour, je voudrais ouvrir un compte épargne. »
\item \cle{Poser des questions précises} : utilise des questions fermées (\eng{Do you offer...?}) pour obtenir oui/non, et des questions ouvertes avec \eng{what}, \eng{how much}, \eng{how long} pour obtenir des détails : \eng{What documents do I need?} / \eng{How much is the minimum deposit?}
\item \cle{Demander une précision ou une répétition} : \eng{Sorry, could you repeat that, please?} / \eng{What do you mean by “overdraft”?}
\item \cle{Donner et vérifier l'information} : reformule pour confirmer : \eng{So, if I understand correctly, I need my ID card and a passport photo.} « Donc, si je comprends bien, j'ai besoin de ma carte d'identité et d'une photo d'identité. »
\item \cle{Conclure poliment} : \eng{Thank you very much for your help. Have a nice day!}
\end{enumerate}
Réflexe : à l'oral, privilégie les formes contractées (\eng{I'd like}, \eng{I'll}) et les formules de politesse avec \eng{could} et \eng{would}, plus courtoises que \eng{can} et \eng{want}.
\end{methode}

\courssub{Lexique des services bancaires}

\begin{propriete}[Vocabulaire thématique — Banking Services]
\begin{center}
\begin{tabularx}{\linewidth}{|X|X|X|}
\hline
\rowcolor{popblueL} \textbf{Français} & \textbf{English} & \textbf{Remarque / Contexte camerounais} \\
\hline
un compte (bancaire) & \eng{account} & \eng{savings account} (compte épargne), \eng{current account} (compte courant / chèque) \\
\hline
un dépôt & \eng{deposit} & \eng{make a deposit} ; \eng{minimum deposit} (dépôt minimum) \\
\hline
un retrait & \eng{withdrawal} & \eng{make a withdrawal} ; attention à l'orthographe : un seul \eng{l} final \\
\hline
le solde & \eng{balance} & \eng{check my balance} ; un seul \eng{l} \\
\hline
un prêt & \eng{loan} & \eng{apply for a loan} ; \eng{loan officer} (chargé de clientèle prêts) \\
\hline
taux d'intérêt & \eng{interest rate} & \eng{annual interest rate} ; \eng{fixed / variable rate} \\
\hline
frais (bancaires) & \eng{fee(s)} / \eng{charge(s)} & \eng{monthly maintenance fee} ; \eng{transaction fee} \\
\hline
distributeur automatique & \eng{ATM} (Automated Teller Machine) & \eng{ATM card} ; \eng{PIN code} (code secret) \\
\hline
carte bancaire & \eng{bank card} / \eng{debit card} & \eng{blocked card} (carte bloquée) ; \eng{lost card} (carte perdue) \\
\hline
virement / transfert & \eng{transfer} & \eng{mobile money transfer} (très usité au Cameroun : MTN MoMo, Orange Money) \\
\hline
reçu & \eng{receipt} & prononcé « ri-site » (\eng{p} muet) ; \eng{Can I have a receipt, please?} \\
\hline
découvert & \eng{overdraft} & \eng{overdraft facility} (autorisation de découvert) \\
\hline
pièce d'identité & \eng{ID card} / \eng{identity card} & \eng{National ID card} (CNI camerounaise) \\
\hline
justificatif de domicile & \eng{proof of address} & facture d'électricité (\eng{electricity bill}), quittance de loyer (\eng{tenancy agreement}) \\
\hline
photo d'identité & \eng{passport photo} / \eng{passport-sized photo} & souvent 2 ou 3 exigées \\
\hline
argent liquide & \eng{cash} & \eng{deposit cash} / \eng{withdraw cash} — \textbf{jamais} « liquid money » \\
\hline
chéquier & \eng{cheque book} & orthographe britannique \eng{cheque} (US : \eng{check}) \\
\hline
épargne & \eng{savings} & \eng{savings passbook} (livret d'épargne) \\
\hline
microfinance & \eng{microfinance institution} / \eng{MFI} & ex. : \eng{MC2}, \eng{CAMECC}, \eng{UCCGN} au Cameroun \\
\hline
\end{tabularx}
\end{center}
\end{propriete}

\courssub{Formules fonctionnelles}

\begin{propriete}[Formules clés pour l'échange bancaire]
\begin{center}
\begin{tabularx}{\linewidth}{|l|X|}
\hline
\rowcolor{popblueL} \textbf{Fonction} & \textbf{Formules types (anglais / français)} \\
\hline
\cle{Demander un service} & \eng{I would like to open an account.} / \eng{Could I open an account?} / \eng{Can you help me with...?} \\
\hline
\cle{Poser une question fermée} & \eng{Do you offer student accounts?} / \eng{Is the ATM working today?} / \eng{Can I withdraw money here?} \\
\hline
\cle{Poser une question ouverte} & \eng{What documents do I need?} / \eng{How much is the fee?} / \eng{How long does a transfer take?} / \eng{What is the interest rate?} \\
\hline
\cle{Demander une précision} & \eng{Sorry, could you repeat that, please?} / \eng{What do you mean by...?} / \eng{Could you speak more slowly, please?} \\
\hline
\cle{Donner une information} & \eng{The minimum deposit is 10,000 francs.} / \eng{You need your ID card and two photos.} / \eng{The bank closes at 4 p.m.} \\
\hline
\cle{Vérifier / Reformuler} & \eng{So, I need to bring my ID card, right?} / \eng{If I understand correctly, the fee is 2,000 francs.} / \eng{Let me check: the rate is 3.5\%.} \\
\hline
\cle{Relancer} & \eng{And what about the interest rate?} / \eng{Any other fees I should know about?} / \eng{What about mobile money transfers?} \\
\hline
\cle{Exprimer un problème} & \eng{My ATM card is blocked.} / \eng{I lost my card.} / \eng{I entered the wrong PIN three times.} / \eng{There is a mistake on my statement.} \\
\hline
\cle{Conclure / Remercier} & \eng{Thank you very much for your help.} / \eng{That's all for today.} / \eng{Have a nice day!} / \eng{Goodbye!} \\
\hline
\end{tabularx}
\end{center}
\end{propriete}

\courssub{Poser des questions et répondre de façon adaptée}

\begin{methode}[Construire ses questions : fermées ou ouvertes ?]
\begin{enumerate}
\item \cle{Question fermée} (réponse \eng{yes/no}) : auxiliaire + sujet + verbe. \eng{Do you accept mobile money transfers?} / \eng{Is the bank open on Saturdays?} / \eng{Can I have a receipt?}
\item \cle{Question ouverte} (information détaillée) : mot interrogatif + auxiliaire + sujet + verbe. \eng{What time does the bank close?} / \eng{How long does a transfer to Bamenda take?} / \eng{Why is my card blocked?}
\item \cle{Choisir le bon mot interrogatif} :
\begin{itemize}
\item \eng{what} (chose, nature) — \eng{What type of account?}
\item \eng{when / what time} (moment) — \eng{When does it open?}
\item \eng{where} (lieu) — \eng{Where is the nearest ATM?}
\item \eng{how much} (prix, quantité indénombrable) — \eng{How much is the fee?}
\item \eng{how many} (quantité dénombrable) — \eng{How many photos do I need?}
\item \eng{how long} (durée) — \eng{How long does it take?}
\item \eng{why} (raison) — \eng{Why is it blocked?}
\end{itemize}
\item \cle{Répondre de façon adaptée} : réponse courte d'abord (\eng{Yes, we do.} / \eng{No, it isn't.}), puis développement (\eng{We accept transfers from Monday to Friday.} / \eng{It opens at 8 a.m.}).
\item \cle{Relancer} si la réponse est incomplète : \eng{And what about the interest rate?} / \eng{Could you tell me more about the fees?}
\end{enumerate}
\end{methode}

\begin{exoresolu}[Compléter un dialogue à la banque]
\textbf{Énoncé.} Complete the dialogue with appropriate questions or answers.

\textbf{Customer:} Good morning. (1) .............................................. a savings account?

\textbf{Clerk:} Certainly, madam. Do you have your national ID card with you?

\textbf{Customer:} Yes, I do. (2) .............................................. the minimum deposit?

\textbf{Clerk:} It is 10,000 francs CFA.

\textbf{Customer:} (3) .............................................. a cheque book with this account?

\textbf{Clerk:} Yes, but you will pay a small fee of 2,000 francs per year.

\textbf{Customer:} Thank you very much for your help.
\tcblower
\textbf{Solution détaillée.}
\begin{enumerate}
\item La réponse du guichetier (\eng{Certainly, madam}) montre qu'il s'agit d'une demande polie de service. Structure à retenir : \eng{Could I + base verbale} ou \eng{I would like to + base verbale}. Réponse : \eng{Could I open a savings account?} ou \eng{I would like to open a savings account. Can you help me?}
\item La réponse donne un montant (\eng{10,000 francs CFA}), donc la question porte sur la quantité d'argent : question ouverte avec \eng{how much} + auxiliaire \eng{is} (le dépôt est singulier). Réponse : \eng{How much is the minimum deposit?} ou \eng{What is the minimum deposit?}
\item La réponse commence par \eng{Yes}, donc question fermée avec inversion de l'auxiliaire. Le verbe \eng{get} (obtenir) ou \eng{have} exige l'auxiliaire \eng{do} au présent, sujet \eng{I}. Réponse : \eng{Do I get a cheque book with this account?} ou \eng{Can I have a cheque book with this account?}
\end{enumerate}
\textbf{Conclusion} : toujours lire la réponse qui suit le blanc : elle révèle le type de question attendu (fermée si la réponse commence par \eng{yes/no}, ouverte avec \eng{how much/what} si la réponse donne une information chiffrée).
\end{exoresolu}

\begin{exoresolu}[Formuler les questions correspondant aux réponses]
\textbf{Énoncé.} Write the question that matches each answer.
\begin{enumerate}
\item .............................................. ? — The bank opens at 8 a.m.
\item .............................................. ? — No, we do not charge any fees for students.
\item .............................................. ? — A withdrawal at the ATM costs 300 francs.
\item .............................................. ? — You need two passport photos and your ID card.
\end{enumerate}
\tcblower
\textbf{Solution détaillée.}
\begin{enumerate}
\item La réponse indique une heure : mot interrogatif \eng{what time} (ou \eng{when}), auxiliaire \eng{does} car le sujet \eng{the bank} est à la troisième personne du singulier, verbe à la base verbale \eng{open}. Réponse : \eng{What time does the bank open?}
\item La réponse commence par \eng{No, we do not} : question fermée avec \eng{do you}. Attention : \eng{any} dans la réponse négative, mais on peut garder \eng{any} dans la question. Réponse : \eng{Do you charge any fees for students?}
\item La réponse donne un prix : \eng{how much} + \eng{does} (sujet \eng{a withdrawal}, singulier) + base verbale \eng{cost}. Réponse : \eng{How much does a withdrawal at the ATM cost?}
\item La réponse énumère des choses nécessaires : \eng{what} + \eng{do} (sujet \eng{you}) + \eng{need}. Réponse : \eng{What do I need to open an account?} (ou \eng{What documents do I need?})
\end{enumerate}
\textbf{Conclusion} : méthode infaillible : souligner l'information donnée dans la réponse, choisir le mot interrogatif qui la cible, puis accorder l'auxiliaire avec le sujet.
\end{exoresolu}

\courssub{Mener un court dialogue : le jeu de rôle guidé}

\begin{methode}[Préparer et jouer un rôle à la banque]
\begin{enumerate}
\item \cle{Lire la consigne} et identifier ton rôle (client ou employé) et ton objectif (ouvrir un compte, retirer de l'argent, signaler un problème de carte).
\item \cle{Préparer 3 à 5 répliques} avec les structures apprises : salutation, demande polie, deux questions, conclusion.
\item \cle{Réagir à l'imprévu} : si tu ne comprends pas, dis \eng{Sorry, could you say that again, please?} Ne reste jamais silencieux.
\item \cle{Soigner la prononciation et l'intonation} : voix montante pour les questions fermées (\eng{Do you accept deposits?} $\nearrow$), voix descendante pour les questions en \eng{wh-} (\eng{What is the interest rate?} $\searrow$). Accentue les mots porteurs de sens : \eng{I would LIKE to OPEN an acCOUNT.}
\item \cle{Conclure et remercier} : \eng{Thank you for your help. Goodbye!}
\end{enumerate}
Lexique utile : \eng{account} (compte), \eng{deposit} (dépôt), \eng{withdrawal} (retrait), \eng{balance} (solde), \eng{loan} (prêt), \eng{interest rate} (taux d'intérêt), \eng{ATM} (distributeur), \eng{fee} (frais).
\end{methode}

\begin{exoresolu}[Jeu de rôle : signaler un problème de carte bancaire]
\textbf{Énoncé.} You are a customer at a bank in Douala. Your ATM card is blocked. Write a short dialogue (at least six exchanges) with the bank clerk.
\tcblower
\textbf{Solution détaillée.} Plan : salutation $\rightarrow$ exposition du problème $\rightarrow$ questions du guichetier $\rightarrow$ vérification d'identité $\rightarrow$ solution proposée $\rightarrow$ remerciement.

\eng{Customer: Good morning, sir. My ATM card is blocked and I cannot withdraw any money.}

\eng{Clerk: Good morning. I am sorry to hear that. May I see your ID card, please?}

\eng{Customer: Yes, here it is. Could you tell me why my card is blocked?}

\eng{Clerk: Let me check... You entered the wrong PIN code three times yesterday.}

\eng{Customer: Oh, I see! What should I do to unblock it?}

\eng{Clerk: You just need to fill in this form, and we will give you a new card within one week.}

\eng{Customer: How much will the new card cost?}

\eng{Clerk: It costs 5,000 francs.}

\eng{Customer: All right. Thank you very much for your help.}

\eng{Clerk: You are welcome. Have a nice day!}

\textbf{Justifications grammaticales} : \eng{May I see...?} (demande polie avec inversion de \eng{may}) ; \eng{you entered} (prétérit pour un événement passé daté : \eng{yesterday}) ; \eng{will give} (futur pour une décision/promesse) ; \eng{within one week} (préposition de délai). Le dialogue respecte l'alternance des tours de parole et les formules de politesse exigées.
\end{exoresolu}

\courssub{Prononciation, accentuation et intonation}

\begin{methode}[Bien prononcer le vocabulaire bancaire]
\begin{enumerate}
\item \cle{Accentue la bonne syllabe} (notation entre guillemets français) :
\begin{itemize}
\item \eng{DEposit} (« dé-po-zite ») — nom, accent sur 1\textsuperscript{re} syllabe.
\item \eng{withDRAwal} (« ouiz-dro-ol ») — accent sur 2\textsuperscript{e} syllabe (après préfixe \eng{with-}).
\item \eng{BALance} (« ba-lence ») — accent sur 1\textsuperscript{re} syllabe.
\item \eng{INterest} (« in-treste ») — accent sur 1\textsuperscript{re} syllabe.
\item \eng{acCOUNT} (« a-kaounte ») — accent sur 2\textsuperscript{e} syllabe (appréfixe \eng{ac-}).
\item \eng{MAnager} (« mé-neu-djé ») — accent sur 1\textsuperscript{re} syllabe.
\end{itemize}
\item \cle{Attention au \eng{th}} de \eng{three}, \eng{thousand}, \eng{withdrawal} : langue entre les dents, souffle doux, jamais « s » ni « z » à la française.
\item \cle{Chiffres et montants} : distingue \eng{thirteen} (« ceur-tine », accent final, son long) et \eng{thirty} (« ceur-ti », accent initial, son court) — confusion classique au guichet !
\item \cle{Intonation montante} ($\nearrow$) pour les questions fermées, \cle{descendante} ($\searrow$) pour les questions ouvertes (\eng{wh-}) et les affirmations.
\item \cle{Lie les mots} dans les formules fréquentes : \eng{I'd like to} se prononce « aïe laïk tou » en un seul souffle ; \eng{How much is it?} $\rightarrow$ « ha-mou-tchi-zit ».
\end{enumerate}
\end{methode}

\begin{exoresolu}[Prononciation et intonation]
\textbf{Énoncé.} 1) Indicate the stressed syllable in: \eng{deposit}, \eng{withdrawal}, \eng{account}, \eng{manager}. 2) Say whether the intonation rises ($\nearrow$) or falls ($\searrow$): a) \eng{Do you have an account here?} b) \eng{Where is the nearest ATM?} c) \eng{The bank closes at 4 p.m.}
\tcblower
\textbf{Solution détaillée.}
\begin{enumerate}
\item \eng{DEposit} (première syllabe) ; \eng{withDRAwal} (deuxième syllabe) ; \eng{acCOUNT} (deuxième syllabe) ; \eng{MAnager} (première syllabe). Astuce : les noms de deux syllabes portent souvent l'accent sur la première syllabe (\eng{DEposit}, \eng{MAnager}), mais les mots formés avec un préfixe (\eng{ac-}, \eng{with-}) reportent l'accent après le préfixe.
\item a) $\nearrow$ : question fermée (réponse \eng{yes/no}), la voix monte. b) $\searrow$ : question ouverte en \eng{wh-}, la voix descend. c) $\searrow$ : phrase affirmative, la voix descend en fin d'énoncé.
\end{enumerate}
\textbf{Conclusion} : une mauvaise intonation peut transformer une question en affirmation et créer un malentendu au guichet ; à l'oral du Bac, l'intonation correcte est un critère d'évaluation explicite.
\end{exoresolu}

\courssub{Entraînement au Bac : production orale guidée}

\begin{methode}[Réussir l'épreuve orale sur les services bancaires]
\begin{enumerate}
\item \cle{Écoute la consigne} de l'examinateur et reformule mentalement ta tâche.
\item \cle{Structure ta prise de parole} : salutation $\rightarrow$ objectif $\rightarrow$ 2 ou 3 questions $\rightarrow$ réaction aux réponses $\rightarrow$ conclusion polie.
\item \cle{Utilise le lexique précis} : dis \eng{withdraw money} (et non « take money »), \eng{open an account}, \eng{transfer money}, \eng{check my balance}.
\item \cle{Emploie les temps justes} : présent simple pour les horaires et habitudes (\eng{The bank opens at 8.}), prétérit pour un incident passé (\eng{I lost my card yesterday.}), \eng{would like} pour les demandes.
\item \cle{Reste naturel} : si tu bloques, utilise \eng{Well...}, \eng{Let me see...}, puis reprends.
\end{enumerate}
\end{methode}

\begin{exoresolu}[Sujet type Bac : interaction à la microfinance]
\textbf{Énoncé.} You want to save money for your university studies. You go to a microfinance institution in Yaoundé. Ask the officer at least three questions about their savings services, then react to the answers.
\tcblower
\textbf{Solution détaillée.} Production attendue (modèle) :

\eng{Good afternoon, madam. My name is Ngo Bell and I am a student at Lycée Leclerc. I would like to save money for my university studies.}

\eng{First, what types of savings accounts do you offer for students?}

\eng{... Oh, a youth account? That sounds interesting. How much is the minimum deposit?}

\eng{... Only 5,000 francs? That is perfect for me. And what is the interest rate per year?}

\eng{... I see. One last question: can I withdraw my money at any time, or is there a notice period?}

\eng{... Thank you very much for this information. I will come back tomorrow with my documents. Goodbye!}

\textbf{Analyse} : le candidat salue et se présente (contexte clair), pose trois questions ouvertes/fermées variées (\eng{what types}, \eng{how much}, \eng{can I}), réagit à chaque réponse (\eng{That sounds interesting}, \eng{That is perfect for me}), et conclut poliment avec une décision (\eng{I will come back}). Temps corrects : présent simple pour les conditions du service, \eng{will} pour la décision future. Ce sont exactement les critères de la grille d'évaluation orale : pertinence, correction linguistique, prononciation, fluidité.
\end{exoresolu}

\courssub{Contexte culturel : Banque et microfinance au Cameroun}

\begin{savaistu}[Le système bancaire camerounais]
Au Cameroun, le secteur bancaire est supervisé par la \eng{COBAC} (Commission Bancaire de l'Afrique Centrale) et la \eng{BEAC} (Banque des États de l'Afrique Centrale). La monnaie est le \eng{Franc CFA} (XAF), arrimé à l'euro (1 € = 655,957 FCFA).
\begin{itemize}
\item \cle{Banques commerciales} : Afriland First Bank, Société Générale Cameroun, BICEC, UBA, Ecobank, Standard Chartered, etc. Elles offrent comptes courants, épargne, crédits, cartes VISA/Mastercard.
\item \cle{Microfinance (MFI)} : très présentes (MC2, CAMECC, UCCGN, ACEP, etc.), elles ciblent les populations non bancarisées, les étudiants, les petits commerçants. Conditions d'ouverture souvent plus souples (dépôt minimum faible, pièces justificatives allégées).
\item \cle{Mobile Money} : MTN MoMo et Orange Money sont omniprésents. On peut \eng{deposit}, \eng{withdraw}, \eng{transfer} et même \eng{save} (\eng{MoMo Savings}) depuis son téléphone. Les banques proposent désormais la liaison \eng{bank-to-wallet} et \eng{wallet-to-bank}.
\item \cle{Horaires types} : Lundi–Vendredi 8h–15h30 (parfois 16h) ; quelques agences ouvertes le samedi matin 8h–12h. Fermées jours fériés.
\item \cle{Documents standards} : CNI (carte nationale d'identité) ou passeport, 2 photos d'identité, justificatif de domicile de moins de 3 mois (facture ENEO/Camwater ou quittance de loyer), parfois une attestation d'hébergement.
\end{itemize}
\end{savaistu}

\courssub{Exercices d'entraînement}

\begin{exercice}[Mise en situation — Compléter le dialogue]
\textbf{Énoncé.} You are at the bank in Buea. Complete the dialogue with your own words (customer's part).
\begin{enumerate}
\item \textbf{Clerk:} Good morning. How can I help you?
\textbf{You:} ...........................................................................................................................
\item \textbf{Clerk:} Yes, we have a student savings account. Do you have your ID card?
\textbf{You:} ...........................................................................................................................
\item \textbf{Clerk:} The minimum deposit is 5,000 francs.
\textbf{You:} ...........................................................................................................................
\item \textbf{Clerk:} The interest rate is 4\% per year.
\textbf{You:} ...........................................................................................................................
\item \textbf{Clerk:} You need your ID card, two passport photos and a proof of address.
\textbf{You:} ...........................................................................................................................
\end{enumerate}
\end{exercice}

\begin{corrige}[Exercice 1 — Proposition de réponses]
\begin{enumerate}
\item \eng{Good morning. I would like to open a student savings account, please.}
\item \eng{Yes, here is my national ID card.} (ou \eng{Yes, I have it with me.})
\item \eng{How much is the minimum deposit?} (déjà donné, mais l'élève peut demander : \eng{And what is the interest rate?})
\item \eng{What is the interest rate?} / \eng{Is that a good rate?}
\item \eng{What documents do I need exactly?} / \eng{So, I need my ID, two photos and a proof of address. Is that all?}
\end{enumerate}
\end{corrige}

\begin{exercice}[Traduction — Du français vers l'anglais]
\textbf{Énoncé.} Translate into English.
\begin{enumerate}
\item Je voudrais retirer 50 000 francs, s'il vous plaît.
\item Ma carte est bloquée car j'ai tapé le mauvais code trois fois.
\item Quel est le taux d'intérêt pour un compte épargne ?
\item Combien de temps dure un virement vers Douala ?
\item Je n'ai pas de justificatif de domicile avec moi.
\end{enumerate}
\end{exercice}

\begin{corrige}[Exercice 2 — Traduction]
\begin{enumerate}
\item \eng{I would like to withdraw 50,000 francs, please.} (Attention : \eng{withdraw}, pas \eng{take}).
\item \eng{My card is blocked because I entered the wrong PIN three times.} (Prétérit \eng{entered} car action passée datée implicitement).
\item \eng{What is the interest rate for a savings account?}
\item \eng{How long does a transfer to Douala take?}
\item \eng{I don't have a proof of address with me.} (Pas « I have not »).
\end{enumerate}
\end{corrige}

\begin{exercice}[Jeu de rôle libre — Scénario au choix]
\textbf{Énoncé.} Choose ONE situation and write a dialogue of at least 8 exchanges (4 per person).
\begin{enumerate}
\item \textbf{Situation A:} You want to apply for a student loan at a bank in Yaoundé. Ask about conditions, interest rate, repayment period, and required documents.
\item \textbf{Situation B:} You received a money transfer from your uncle in France via Western Union. You go to the bank to collect the cash. The clerk asks for your ID and the transfer reference number (MTCN).
\item \textbf{Situation C:} You notice a mistake on your bank statement (a withdrawal you didn't make). You go to the bank to complain and ask for an investigation.
\end{enumerate}
\end{exercice}

\begin{corrige}[Exercice 3 — Pistes de correction (Situation A)]
\textbf{Customer:} Good morning. I would like to apply for a student loan for my university fees.
\textbf{Clerk:} Good morning. We have a specific student loan. Do you have a guarantor?
\textbf{Customer:} Yes, my father works at the Ministry of Finance. What is the interest rate?
\textbf{Clerk:} It is 5\% per year, fixed. The repayment period starts six months after graduation.
\textbf{Customer:} That sounds fair. How much can I borrow maximum?
\textbf{Clerk:} Up to 2,000,000 francs, depending on your guarantor's income.
\textbf{Customer:} What documents do I need to provide?
\textbf{Clerk:} Your ID, student card, admission letter, your father's pay slips and a guarantee letter.
\textbf{Customer:} Thank you very much. I will bring the file next week.
\textbf{Clerk:} You are welcome. Good luck with your studies!
\end{corrige}

\begin{attention}[Les 5 erreurs qui coûtent des points au Bac]
\begin{enumerate}
\item \cle{Confondre les faux amis} : \eng{actually} signifie « en fait », pas « actuellement » (\eng{currently}) ; une \eng{bank note} est un billet de banque, pas une « note de banque » ; \eng{to demand} est agressif — dis \eng{to ask for} ou \eng{to request}.
\item \cle{Oublier l'auxiliaire dans les questions} : jamais \eng{How much costs a withdrawal?} mais \eng{How much does a withdrawal cost?} ; jamais \eng{You have an account?} à l'écrit, mais \eng{Do you have an account?} (sauf intonation montante à l'oral très familier).
\item \cle{Calquer le français} : « prendre de l'argent » ne se traduit pas par \eng{take money} mais \eng{withdraw money} ; « faire un compte » n'est pas \eng{make an account} mais \eng{open an account} ; « l'argent liquide » se dit \eng{cash}, pas « liquid money ».
\item \cle{Se tromper de temps} : pour un incident passé et daté, utilise le prétérit : \eng{I lost my card yesterday} (et non \eng{I have lost} avec \eng{yesterday}) ; pour les horaires, le présent simple : \eng{The bank closes at 4 p.m.}
\item \cle{Massacrer les mots-clés} : orthographe — \eng{withdrawal} (avec un \eng{l} final), \eng{balance} (un seul \eng{l}), \eng{receipt} (« ri-site », le \eng{p} ne se prononce pas) ; prononciation — ne confonds jamais \eng{thirteen} et \eng{thirty} devant un montant.
\end{enumerate}
\end{attention}

\newpage

\coursec{Exercices}

\courssub{Je vérifie mes connaissances}

\begin{exercice}[QCM : le bon mot bancaire]
Pour chaque phrase, choisis la bonne réponse (A, B ou C).

\begin{enumerate}
\item I would like to \textbf{......} some money from my account. I need cash for the market.
\begin{enumerate}
\item deposit \item withdraw \item borrow
\end{enumerate}
\item My father asked for a \textbf{......} of 500,000 francs to buy a taxi. He will pay it back every month.
\begin{enumerate}
\item loan \item balance \item receipt
\end{enumerate}
\item Could you tell me my account \textbf{......}, please? I want to know how much money I have left.
\begin{enumerate}
\item transfer \item statement \item balance
\end{enumerate}
\item Every month, my aunt \textbf{......} money from Buea to her son in Yaoundé.
\begin{enumerate}
\item transfers \item withdraws \item lends
\end{enumerate}
\item You must fill in this \textbf{......} before you can open an account.
\begin{enumerate}
\item form \item coin \item queue
\end{enumerate}
\end{enumerate}
\end{exercice}

\begin{exercice}[Vrai ou faux ? Justifie ta réponse]
Dis si les affirmations suivantes sont vraies (V) ou fausses (F), puis corrige les affirmations fausses.

\begin{enumerate}
\item \eng{“Withdraw” means to put money into your account.}
\item \eng{“Could you tell me...?” is more polite than “Tell me...”.}
\item \eng{A “bank statement” is a document that shows the operations on your account.}
\item \eng{To end a conversation at the bank, you can say: “Thank you for your help. Goodbye!”}
\item \eng{In English, we say “How much is the balance?” to ask for the amount of money in an account.}
\end{enumerate}
\end{exercice}

\begin{exercice}[Vocabulaire : trouve le mot]
Trouve le mot anglais du vocabulaire bancaire qui correspond à chaque définition.

\begin{enumerate}
\item The person who works behind the counter and serves customers in a bank.
\item The money you pay to the bank for a service (for example, for a transfer).
\item A small piece of paper the bank gives you to prove you have paid or received money.
\item The secret number you use with your bank card.
\item The place in the bank where you can take money with a card, day and night.
\end{enumerate}
\end{exercice}

\begin{exercice}[Conjugaison et formes : complète le dialogue]
Complète les répliques avec le verbe entre parenthèses à la forme correcte.

\begin{textebox}[At the counter]
Clerk: Good morning, Madam. How \textbf{......} (can) I help you?

Customer: Good morning. I \textbf{......} (would like) to open a savings account.

Clerk: Certainly. \textbf{......} you \textbf{......} (bring) an identity card with you?

Customer: Yes, here it is. Yesterday I \textbf{......} (deposit) 20,000 francs in my brother's account, but I \textbf{......} (not / receive) a receipt.

Clerk: Don't worry. I \textbf{......} (check) that for you right now.
\end{textebox}
\end{exercice}

\courssub{Je m'entraîne}

\begin{exercice}[Traduction : à la banque]
Traduis les phrases suivantes en anglais, en utilisant les formules polies étudiées dans la leçon.

\begin{enumerate}
\item Bonjour Monsieur, je voudrais retirer 50,000 francs de mon compte, s'il vous plaît.
\item Pourriez-vous me dire quel est le solde de mon compte ?
\item Combien sont les frais pour un transfert vers Bamenda ?
\item Excusez-moi, je n'ai pas compris le montant. Pouvez-vous répéter, s'il vous plaît ?
\item Merci beaucoup pour votre aide. Au revoir !
\end{enumerate}
\end{exercice}

\begin{exercice}[Remets le dialogue dans l'ordre]
Les répliques du dialogue suivant sont dans le désordre. Numérote-les de 1 à 8 pour reconstituer une conversation logique à la banque (salutation → demande → informations → relance → conclusion).

\begin{enumerate}
\item \textbf{A.} \eng{Good morning, Madam. How can I help you today?}
\item \textbf{B.} \eng{Thank you very much for your help. Have a nice day!}
\item \textbf{C.} \eng{Good morning. I would like to withdraw 50,000 francs, please.}
\item \textbf{D.} \eng{You're welcome. Goodbye, Madam.}
\item \textbf{E.} \eng{Of course. Could I see your identity card, please?}
\item \textbf{F.} \eng{Certainly. Your balance is 120,500 francs.}
\item \textbf{G.} \eng{Here you are. Could you also tell me my balance?}
\item \textbf{H.} \eng{I'm sorry, could you repeat the amount, please? I didn't hear well.}
\end{enumerate}
\end{exercice}

\begin{exercice}[Dialogue à trous : les formules fonctionnelles]
Complète le dialogue avec les formules de la liste suivante (chaque formule ne sert qu'une fois) :

\begin{center}
\eng{Could you tell me — I would like — How can I help you — You're welcome — Could you repeat that, please — Thank you for your help}
\end{center}

\begin{textebox}[At UBA Bank, Douala]
Clerk: Good afternoon, Madam. \textbf{......} today?

Customer: Good afternoon. \textbf{......} to transfer some money to my sister in Bamenda.

Clerk: No problem. \textbf{......} her account number, please?

Customer: It's 1002 4567. How much are the charges?

Clerk: The charges are 1,500 francs.

Customer: \textbf{......} I didn't understand the amount.

Clerk: One thousand five hundred francs, Madam.

Customer: All right. Here is the money. \textbf{......}.

Clerk: \textbf{......}, Madam. Goodbye!
\end{textebox}
\end{exercice}

\begin{exercice}[Appariement : la bonne réaction]
Associe chaque situation (1--6) à la réaction la plus adaptée (A--F).

\begin{enumerate}
\item You did not understand the amount the clerk said.
\item You want to know the cost of the operation.
\item You want to open an account.
\item You want to thank the clerk at the end.
\item You want to know how much money is in your account.
\item The clerk is speaking too fast.
\end{enumerate}

\begin{enumerate}
\item \textbf{A.} \eng{Could you tell me my balance, please?}
\item \textbf{B.} \eng{Sorry, could you speak more slowly, please?}
\item \textbf{C.} \eng{Excuse me, could you repeat the amount, please?}
\item \textbf{D.} \eng{I would like to open a savings account, please.}
\item \textbf{E.} \eng{Thank you very much for your help!}
\item \textbf{F.} \eng{How much are the charges for this service?}
\end{enumerate}
\end{exercice}

\courssub{Je me prépare au Bac}

\begin{exercice}[Compréhension écrite et langue — type Bac]
Lis le texte suivant, puis réponds aux questions en anglais, par des phrases complètes.

\begin{textebox}[Banking services in Cameroon]
Every morning, Mrs Etoundi leaves her shop in the Mokolo market in Yaoundé and walks to the bank on the main avenue. Like many traders, she needs banking services almost every day: she deposits the money from her sales, withdraws cash to pay her suppliers, and sometimes transfers money to her daughter who studies in Buea.

“Before, I kept all my money at home, and I was always afraid of thieves,” she explains. “Now, with my savings account, my money is safe, and I can check my balance whenever I want.”

However, using banking services is not always easy for everyone. Some customers do not understand the documents, which are often written in difficult language. Others complain about the charges, which they find too high. “When I send 100,000 francs to my daughter, the bank takes more than 2,000 francs,” says Mr Nkeng, a taxi driver from Douala. “That is a lot of money for a poor family.”

To help their customers, many banks now employ clerks who speak both English and French, and some banks explain their services in Pidgin too. Mobile banking is also changing habits: with a simple telephone, people can now transfer money without going to the bank at all. For Mrs Etoundi, however, nothing replaces a conversation with a friendly clerk: “When I have a problem, I prefer to speak to a real person, face to face.”
\end{textebox}

\textbf{A. Comprehension questions (answer in your own words)}

\begin{enumerate}
\item Why does Mrs Etoundi go to the bank almost every day? Give two reasons.
\item Why was she afraid before she opened her account?
\item What two problems do some customers have with banking services?
\item What do banks do to help customers who have language problems?
\item According to the text, what advantage does mobile banking offer?
\end{enumerate}

\textbf{B. Language work}

\begin{enumerate}
\item Find in the text the opposite of: (a) \eng{dangerous} ; (b) \eng{easy} ; (c) \eng{cheap}.
\item Find in the text a word or expression that means: (a) \eng{argent liquide} ; (b) \eng{frais bancaires} ; (c) \eng{vérifier son solde}.
\item Put the verbs in brackets in the correct tense: \eng{“Every morning, Mrs Etoundi ...... (leave) her shop and ...... (walk) to the bank. Before, she ...... (keep) her money at home.”}
\end{enumerate}
\end{exercice}

\begin{exercice}[Traduction — type Bac]
Traduis le passage suivant en anglais.

\begin{textebox}[Texte à traduire]
Hier matin, mon père est allé à la banque pour retirer de l'argent. Au guichet, il a dit : « Bonjour Madame, je voudrais retirer 75,000 francs, s'il vous plaît. Pourriez-vous aussi me dire le solde de mon compte ? » L'employée a répondu poliment et lui a donné un reçu. Mon père l'a remerciée et il est rentré à la maison très content.
\end{textebox}
\end{exercice}

\begin{exercice}[Production orale et écrite — jeu de rôle noté]
\textbf{Situation :} \eng{You are at UBA Bank in Douala. You want to withdraw 50,000 francs and check your balance. The examiner is the bank clerk.}

\begin{enumerate}
\item \textbf{Préparation écrite (10 minutes) :} rédige ton dialogue complet (8 à 10 répliques au total). Tu dois obligatoirement utiliser :
\begin{itemize}
\item une formule de salutation et une formule de conclusion ;
\item \eng{I would like...} pour exprimer ta demande ;
\item une question polie avec \eng{Could you tell me...?} ;
\item une formule pour demander de répéter ou pour relancer ;
\item une formule de remerciement.
\end{itemize}
\item \textbf{Passage oral (2 à 3 minutes) :} joue la scène avec l'examinateur. Soigne ta prononciation : accentue correctement \eng{acCOUNT}, \eng{dePOSIT}, \eng{withDRAWal}, \eng{BALance}, \eng{TRANSfer} (nom), \eng{CUS-tomer}, et fais monter l'intonation (↗) dans les questions polies en \eng{Could you...?}.
\item \textbf{Pour aller plus loin :} l'examinateur peut annoncer que le distributeur est en panne (\eng{The ATM is out of order}). Réagis avec une formule adaptée et demande une solution.
\end{enumerate}
\end{exercice}



\newpage

\coursec{Corrigés des exercices}

\begin{corrige}[Exercice 1 — QCM : le bon mot bancaire]
\begin{enumerate}
\item \textbf{B. \eng{withdraw}} — \eng{I would like to withdraw some money from my account.} « Je voudrais retirer de l'argent de mon compte. » \eng{Withdraw} (\cle{retirer}) s'oppose à \eng{deposit} (\cle{déposer}) ; \eng{borrow} signifie « emprunter », ce qui ne convient pas : l'argent appartient déjà au client.
\item \textbf{A. \eng{loan}} — \eng{My father asked for a loan of 500,000 francs.} Un \eng{loan} (\cle{prêt bancaire}) est remboursé par mensualités (\eng{pay it back every month}). \eng{Balance} = \cle{solde} ; \eng{receipt} = \cle{reçu}.
\item \textbf{C. \eng{balance}} — \eng{Could you tell me my account balance, please?} Le \eng{balance} (\cle{solde}) est la somme restant sur le compte (\eng{how much money I have left}). \eng{Statement} = \cle{relevé de compte} ; \eng{transfer} = \cle{virement}.
\item \textbf{A. \eng{transfers}} — \eng{Every month, my aunt transfers money from Buea to her son.} \eng{Transfer} = \cle{virer, transférer} ; attention à la 3\up{e} personne du singulier : \eng{transfer\textbf{s}}. \eng{Withdraw} = \cle{retirer} ; \eng{lend} = \cle{prêter (à quelqu'un)}.
\item \textbf{A. \eng{form}} — \eng{You must fill in this form before you can open an account.} Un \eng{form} (\cle{formulaire}) se « remplit » avec le phrasal verb \eng{fill in}. \eng{Coin} = \cle{pièce de monnaie} ; \eng{queue} = \cle{file d'attente}.
\end{enumerate}
\end{corrige}

\begin{corrige}[Exercice 2 — Vrai ou faux ?]
\begin{enumerate}
\item \textbf{Faux.} \eng{Withdraw} signifie « retirer de l'argent ». Le verbe qui signifie « mettre de l'argent sur son compte » est \eng{deposit} : \eng{I deposited 20,000 francs into my account.} « J'ai déposé 20 000 francs sur mon compte. »
\item \textbf{Vrai.} \eng{Could you tell me...?} est une formule polie (conditionnel atténué), adaptée à un inconnu ou à un employé de banque. L'impératif \eng{Tell me...} est direct, voire impoli avec un inconnu.
\item \textbf{Vrai.} Un \eng{bank statement} (\cle{relevé de compte}) est un document qui liste toutes les opérations du compte : dépôts, retraits, virements, frais.
\item \textbf{Vrai.} \eng{Thank you for your help. Goodbye!} est une formule de conclusion standard et polie, parfaite pour clore un échange au guichet.
\item \textbf{Faux.} On dit \eng{What is my balance?} ou, plus poliment, \eng{Could you tell me my balance, please?} \eng{How much is the balance?} n'est pas idiomatique : \eng{how much} s'emploie pour un prix (\eng{How much are the charges?}), \eng{what} pour un montant de solde.
\end{enumerate}
\end{corrige}

\begin{corrige}[Exercice 3 — Vocabulaire : trouve le mot]
\begin{enumerate}
\item \eng{Bank clerk} (\cle{employé(e) de banque, guichetier(-ière)}) ; on trouve aussi \eng{teller} en anglais américain et \eng{cashier}.
\item \eng{Charges} ou \eng{fees} (\cle{frais bancaires}) ; on parle aussi de \eng{commission}. Prononciation : « tchar-djiz ».
\item \eng{Receipt} (\cle{reçu}). Attention à la prononciation : « ri-\textbf{ssite} » — le \eng{p} ne se prononce pas.
\item \eng{PIN} (\cle{Personal Identification Number}), aussi appelé \eng{PIN code} ou \eng{secret code} (\cle{code confidentiel}).
\item \eng{ATM} (\cle{Automated Teller Machine}), aussi \eng{cash machine} ou \eng{cashpoint} en Grande-Bretagne (\cle{distributeur automatique de billets}).
\end{enumerate}
\end{corrige}

\begin{corrige}[Exercice 4 — Conjugaison et formes : complète le dialogue]
\begin{textebox}[At the counter — Corrigé]
Clerk: Good morning, Madam. How \textbf{can} I help you?

Customer: Good morning. I \textbf{would like} to open a savings account.

Clerk: Certainly. \textbf{Did} you \textbf{bring} an identity card with you?

Customer: Yes, here it is. Yesterday I \textbf{deposited} 20,000 francs in my brother's account, but I \textbf{did not receive} (didn't receive) a receipt.

Clerk: Don't worry. I \textbf{will check} ('ll check) that for you right now.
\end{textebox}
\emph{Justifications grammaticales :}
\begin{itemize}
\item \eng{How can I help you?} : formule figée d'offre d'aide, modal \eng{can} invariable.
\item \eng{Would like} : forme invariable (conditionnel de politesse), toujours suivie de l'infinitif avec \eng{to} : \eng{I would like to open...}
\item \eng{Did you bring...?} : question au prétérit (auxiliaire \eng{did} + base verbale \eng{bring} ; verbe irrégulier \eng{bring -- brought -- brought}).
\item \eng{Deposited} : prétérit régulier, justifié par le marqueur temporel \eng{Yesterday}. Attention au doublement : \eng{deposit} $\rightarrow$ \eng{deposited} (un seul \eng{t}, la syllabe accentuée est \eng{-POS-}).
\item \eng{Did not receive} : prétérit négatif, auxiliaire \eng{didn't} + base verbale.
\item \eng{Will check} : futur avec \eng{will}, décision prise sur l'instant (\eng{right now}) pour rassurer le client.
\end{itemize}
\end{corrige}

\begin{corrige}[Exercice 5 — Traduction : à la banque]
\begin{enumerate}
\item \eng{Good morning, Sir. I would like to withdraw 50,000 francs from my account, please.} — Points de vigilance : \eng{withdraw... \textbf{from} my account} (préposition \eng{from}) ; pas de futur après \eng{would like}.
\item \eng{Could you tell me what my balance is, please?} ou plus simplement \eng{Could you tell me my balance, please?} — Attention à l'ordre des mots dans l'interrogative indirecte : \eng{what my balance \textbf{is}} (sujet avant verbe), jamais \eng{what is my balance} après \eng{Could you tell me}.
\item \eng{How much are the charges for a transfer to Bamenda?} ou \eng{What are the fees for a transfer to Bamenda?} — \eng{Charges} est pluriel, d'où \eng{are} ; \eng{how much} pour un prix.
\item \eng{Excuse me, I didn't understand the amount. Could you repeat, please?} ou \eng{Sorry, could you repeat the amount, please? I didn't catch it.} — Prétérit \eng{didn't understand} pour « je n'ai pas compris » (action ponctuelle achevée).
\item \eng{Thank you very much for your help. Goodbye!} — Formule figée : \eng{thank you \textbf{for} your help}.
\end{enumerate}
\end{corrige}

\begin{corrige}[Exercice 6 — Remets le dialogue dans l'ordre]
L'ordre logique est : \textbf{1. A} $\rightarrow$ \textbf{2. C} $\rightarrow$ \textbf{3. E} $\rightarrow$ \textbf{4. G} $\rightarrow$ \textbf{5. F} $\rightarrow$ \textbf{6. H} $\rightarrow$ \textbf{7. B} $\rightarrow$ \textbf{8. D}.

\begin{textebox}[Dialogue reconstitué]
Clerk: Good morning, Madam. How can I help you today?

Customer: Good morning. I would like to withdraw 50,000 francs, please.

Clerk: Of course. Could I see your identity card, please?

Customer: Here you are. Could you also tell me my balance?

Clerk: Certainly. Your balance is 120,500 francs.

Customer: I'm sorry, could you repeat the amount, please? I didn't hear well.

Customer: Thank you very much for your help. Have a nice day!

Clerk: You're welcome. Goodbye, Madam.
\end{textebox}
\emph{Logique de la conversation :} salutation du guichetier (A) $\rightarrow$ demande du client (C) $\rightarrow$ vérification d'identité (E) $\rightarrow$ remise du document et nouvelle question (G) $\rightarrow$ information donnée (F) $\rightarrow$ demande de répétition (H, qui fait suite à l'annonce du montant en F) $\rightarrow$ remerciement (B) $\rightarrow$ réponse et congé (D).
\end{corrige}

\begin{corrige}[Exercice 7 — Dialogue à trous : les formules fonctionnelles]
\begin{textebox}[At UBA Bank, Douala — Corrigé]
Clerk: Good afternoon, Madam. \textbf{How can I help you} today?

Customer: Good afternoon. \textbf{I would like} to transfer some money to my sister in Bamenda.

Clerk: No problem. \textbf{Could you tell me} her account number, please?

Customer: It's 1002 4567. How much are the charges?

Clerk: The charges are 1,500 francs.

Customer: \textbf{Could you repeat that, please?} I didn't understand the amount.

Clerk: One thousand five hundred francs, Madam.

Customer: All right. Here is the money. \textbf{Thank you for your help}.

Clerk: \textbf{You're welcome}, Madam. Goodbye!
\end{textebox}
\emph{Points de vigilance :} \eng{How can I help you} est la formule d'accueil du guichetier ; \eng{Could you repeat that, please?} est la formule de relance quand on n'a pas compris ; \eng{You're welcome} (« je vous en prie ») répond toujours à un remerciement — ne jamais le confondre avec \eng{welcome} (« bienvenue »).
\end{corrige}

\begin{corrige}[Exercice 8 — Appariement : la bonne réaction]
\begin{center}
\begin{tabular}{|l|l|}
\hline
\rowcolor{popblueL} Situation & Réaction \\
\hline
1. Montant incompris & \textbf{C} — \eng{Excuse me, could you repeat the amount, please?} \\
\hline
2. Coût de l'opération & \textbf{F} — \eng{How much are the charges for this service?} \\
\hline
3. Ouvrir un compte & \textbf{D} — \eng{I would like to open a savings account, please.} \\
\hline
4. Remercier à la fin & \textbf{E} — \eng{Thank you very much for your help!} \\
\hline
5. Connaître son solde & \textbf{A} — \eng{Could you tell me my balance, please?} \\
\hline
6. Le guichetier parle trop vite & \textbf{B} — \eng{Sorry, could you speak more slowly, please?} \\
\hline
\end{tabular}
\end{center}
\emph{Astuce méthodologique :} repérer les mots-clés de chaque situation (\eng{amount}, \eng{cost}, \eng{open}, \eng{thank}, \eng{balance}, \eng{too fast}) et les faire correspondre aux mots-clés des répliques (\eng{repeat the amount}, \eng{charges}, \eng{open an account}, \eng{thank you}, \eng{balance}, \eng{more slowly}).
\end{corrige}

\begin{corrige}[Exercice 9 — Bac : compréhension écrite et langue]
\textbf{Barème indicatif (sur 20) :} Partie A : 10 points (2 pts par question : 1 pt pour l'information exacte, 1 pt pour la correction de la phrase) ; Partie B : 10 points (question 1 : 3 pts ; question 2 : 3 pts ; question 3 : 4 pts).

\textbf{A. Comprehension questions — réponses modèles}
\begin{enumerate}
\item \eng{Mrs Etoundi goes to the bank almost every day because she deposits the money from her sales and withdraws cash to pay her suppliers.} (On accepte aussi : \eng{she sometimes transfers money to her daughter in Buea}.) — « Elle va à la banque pour déposer l'argent de ses ventes et retirer des liquides pour payer ses fournisseurs. »
\item \eng{Before she opened her account, she was afraid of thieves because she kept all her money at home.} — « Avant, elle gardait tout son argent à la maison et avait peur des voleurs. »
\item \eng{Some customers do not understand the documents because they are written in difficult language, and others complain that the charges are too high.} — « Certains ne comprennent pas les documents ; d'autres trouvent les frais trop élevés. »
\item \eng{To help these customers, many banks employ clerks who speak both English and French, and some banks explain their services in Pidgin.} — « Les banques emploient du personnel bilingue et certaines expliquent leurs services en pidgin. »
\item \eng{Mobile banking allows people to transfer money without going to the bank at all.} — « La banque mobile permet de virer de l'argent sans se déplacer. »
\end{enumerate}

\textbf{B. Language work}
\begin{enumerate}
\item (a) \eng{safe} (« en sécurité », l. 6) ; (b) \eng{difficult} (l. 9) ; (c) \eng{high} dans \eng{too high} (l. 11) — on accepte aussi \eng{a lot of money} comme idée, mais le mot attendu est \eng{high}.
\item (a) \eng{cash} (l. 4) ; (b) \eng{charges} (l. 10) ; (c) \eng{check my balance} (l. 7).
\item \eng{Every morning, Mrs Etoundi \textbf{leaves} her shop and \textbf{walks} to the bank. Before, she \textbf{kept} her money at home.} — \eng{Leaves} et \eng{walks} : présent simple, 3\up{e} personne du singulier (habitude signalée par \eng{every morning}) ; \eng{kept} : prétérit du verbe irrégulier \eng{keep -- kept -- kept} (passé révolu signalé par \eng{before}).
\end{enumerate}

\textbf{Points de vigilance :} répondre par des phrases complètes avec sujet et verbe ; ne pas recopier le texte mot à mot sans l'adapter à la question (\eng{in your own words}) ; attention à l'accord à la 3\up{e} personne (\eng{she deposit\textbf{s}}, \eng{she withdraw\textbf{s}}) ; ne pas confondre \eng{safe} (adjectif, « en sécurité ») et \eng{safety} (nom).
\end{corrige}

\begin{corrige}[Exercice 10 — Bac : traduction]
\textbf{Barème indicatif (sur 10) :} 5 phrases $\times$ 2 pts (1 pt pour le sens, 1 pt pour la correction grammaticale : temps, prépositions, ordre des mots).

\begin{textebox}[Proposition de traduction]
Yesterday morning, my father went to the bank to withdraw some money. At the counter, he said: “Good morning, Madam. I would like to withdraw 75,000 francs, please. Could you also tell me my balance?” The clerk replied politely and gave him a receipt. My father thanked her and went home very happy.
\end{textebox}

\textbf{Points de vigilance :}
\begin{itemize}
\item « Est allé » $\rightarrow$ prétérit irrégulier \eng{went} (\eng{go -- went -- gone}), jamais \eng{goed}.
\item « Pour retirer » $\rightarrow$ infinitif de but : \eng{to withdraw} (pas de \eng{for withdraw}).
\item « Au guichet » $\rightarrow$ \eng{at the counter} (préposition \eng{at}).
\item « Le solde de mon compte » $\rightarrow$ \eng{my balance} ou \eng{the balance of my account} ; pas de \eng{how much}.
\item « A répondu » $\rightarrow$ \eng{replied} ou \eng{answered} ; « lui a donné » $\rightarrow$ \eng{gave him} (\eng{give -- gave -- given}).
\item « Est rentré à la maison » $\rightarrow$ \eng{went home} (sans préposition : jamais \eng{went to home}).
\item « Très content » $\rightarrow$ \eng{very happy} ou \eng{very pleased}.
\end{itemize}
\end{corrige}

\begin{corrige}[Exercice 11 — Bac : production orale et écrite, jeu de rôle noté]
\textbf{Barème indicatif (sur 20) :} respect de la situation et des consignes (5 pts) ; exactitude grammaticale et vocabulaire bancaire (5 pts) ; richesse des formules fonctionnelles (5 pts) ; prononciation, accentuation et intonation à l'oral (5 pts).

\begin{textebox}[Modèle de production écrite]
Clerk: Good morning, Madam. How can I help you today?

You: Good morning. I would like to withdraw 50,000 francs from my account, please.

Clerk: Certainly. Could I see your identity card and your bank card, please?

You: Of course, here you are. Could you also tell me my balance, please?

Clerk: Let me check... Your balance is 125,000 francs.

You: I'm sorry, could you repeat the amount, please? I didn't catch it.

Clerk: One hundred and twenty-five thousand francs, Madam.

You: Thank you. And could I have a receipt, please?

Clerk: Here is your receipt and your money.

You: Thank you very much for your help. Goodbye!

Clerk: You're welcome, Madam. Have a nice day!
\end{textebox}

\textbf{Accentuation tonique à soigner (syllabe forte en majuscules) :}
\begin{center}
\begin{tabular}{|l|l|l|}
\hline
\rowcolor{popblueL} Mot & Accentuation & Prononciation approchée \\
\hline
account & ac\textbf{COUNT} & « a-\textbf{kaounte} » \\
\hline
deposit & de\textbf{POS}it & « di-\textbf{po}-zite » \\
\hline
withdrawal & with\textbf{DRAW}al & « ouiz-\textbf{dro}-ol » \\
\hline
balance & \textbf{BAL}ance & « \textbf{ba}-lens » \\
\hline
transfer (nom) & \textbf{TRANS}fer & « \textbf{trans}-feur » \\
\hline
customer & \textbf{CUS}tomer & « \textbf{keuss}-to-meur » \\
\hline
\end{tabular}
\end{center}

\textbf{Intonation :} faire monter la voix (↗) à la fin des questions polies en \eng{Could you...?} (\eng{Could you tell me my balance, please? ↗}) ; la faire descendre (↘) dans les questions en \eng{wh-} (\eng{How much are the charges? ↘}) et dans les salutations de conclusion (\eng{Goodbye! ↘}).

\textbf{Pour aller plus loin — réaction modèle si le distributeur est en panne :}
\eng{“Oh, the ATM is out of order? That's a problem. What can I do? Can I withdraw the money at the counter instead?”} « Oh, le distributeur est en panne ? C'est embêtant. Que puis-je faire ? Puis-je retirer l'argent au guichet à la place ? »

\textbf{Points de vigilance :} ne pas oublier la salutation d'ouverture ni la formule de congé ; employer \eng{I would like} (et non \eng{I want}, trop direct) ; placer \eng{please} en fin de demande ; à l'oral, ne pas prononcer le \eng{p} de \eng{receipt} et bien marquer l'accent tonique des mots bancaires.
\end{corrige}

\newpage

\coursec{Fiche bilan}

\courssub{Fiche Bilan — Lesson 16 : Exchanges Information about Using Banking Services}

\begin{popfigure}
\begin{tikzpicture}[
  mindmap,
  concept color=popblue,
  level 1 concept/.append style={level distance=3.5cm, sibling angle=360/7},
  level 2 concept/.append style={level distance=2.5cm},
  every node/.style={concept, execute at begin node=\small\bfseries, align=center},
  grow cyclic,
  text=white
]
\node[concept, minimum size=2.2cm, align=center] (center){Banking\\Services}
  child[concept color=popblue] {node[concept, minimum size=1.8cm, align=center]{Key\\Vocabulary}
    child[concept color=popblue!70] {node[concept, minimum size=1.4cm] {Account types}}
    child[concept color=popblue!70] {node[concept, minimum size=1.4cm] {Transactions}}
    child[concept color=popblue!70] {node[concept, minimum size=1.4cm] {People \& Places}}
  }
  child[concept color=poporange] {node[concept, minimum size=1.8cm, align=center]{Opening an\\Account}
    child[concept color=poporange!70] {node[concept, minimum size=1.4cm] {ID \& Proof of address}}
    child[concept color=poporange!70] {node[concept, minimum size=1.4cm] {Forms \& Signature}}
    child[concept color=poporange!70] {node[concept, minimum size=1.4cm] {Initial deposit}}
  }
  child[concept color=popgreen] {node[concept, minimum size=1.8cm, align=center]{Daily\\Transactions}
    child[concept color=popgreen!70] {node[concept, minimum size=1.4cm] {Deposit / Withdraw}}
    child[concept color=popgreen!70] {node[concept, minimum size=1.4cm] {Transfer / Standing order}}
    child[concept color=popgreen!70] {node[concept, minimum size=1.4cm] {Balance \& Statement}}
  }
  child[concept color=poppurple] {node[concept, minimum size=1.8cm, align=center]{Functional\\Phrases}
    child[concept color=poppurple!70] {node[concept, minimum size=1.4cm] {Asking for info}}
    child[concept color=poppurple!70] {node[concept, minimum size=1.4cm] {Giving info}}
    child[concept color=poppurple!70] {node[concept, minimum size=1.4cm] {Clarifying / Confirming}}
  }
  child[concept color=poppink] {node[concept, minimum size=1.8cm] {Pronunciation}
    child[concept color=poppink!70] {node[concept, minimum size=1.4cm] {Word stress}}
    child[concept color=poppink!70] {node[concept, minimum size=1.4cm] {Sentence intonation}}
    child[concept color=poppink!70] {node[concept, minimum size=1.4cm] {Figures / Amounts}}
  }
  child[concept color=popgold] {node[concept, minimum size=1.8cm, align=center]{Roleplay\\Scenarios}
    child[concept color=popgold!70] {node[concept, minimum size=1.4cm] {At the counter}}
    child[concept color=popgold!70] {node[concept, minimum size=1.4cm] {ATM issues}}
    child[concept color=popgold!70] {node[concept, minimum size=1.4cm] {Phone / App banking}}
  }
  child[concept color=popdark] {node[concept, minimum size=1.8cm, align=center]{Common\\Mistakes}
    child[concept color=popdark!70] {node[concept, minimum size=1.4cm] {False friends}}
    child[concept color=popdark!70] {node[concept, minimum size=1.4cm] {Prepositions}}
    child[concept color=popdark!70] {node[concept, minimum size=1.4cm] {Figures / Dates}}
  };
\end{tikzpicture}
\legende{Carte mentale récapitulative : Lesson 16 — Banking Services}
\end{popfigure}

\begin{bilan}

\textbf{1. Lexique clé (Key Vocabulary)}

\begin{center}
\begin{tabularx}{\linewidth}{|X|X|X|}
\hline
\rowcolor{popblueL} \textbf{English} & \textbf{Français} & \textbf{Contexte / Exemple} \\
\hline
\cle{current account} (UK) / \cle{checking account} (US) & compte courant & \eng{I'd like to open a current account.} « Je voudrais ouvrir un compte courant. » \\
\cle{savings account} & compte d'épargne / livret & \eng{The interest rate on this savings account is 3\%.} \\
\cle{deposit} (n./v.) & dépôt / déposer & \eng{I want to make a deposit of 50,000 francs.} \\
\cle{withdrawal} (n.) / \cle{withdraw} (v.) & retrait / retirer & \eng{The daily withdrawal limit is 200,000 FCFA.} \\
\cle{transfer} (n./v.) & virement / virer & \eng{Can I transfer money to a UBA account in Douala?} \\
\cle{standing order} & virement permanent & \eng{Set up a standing order for the rent.} \\
\cle{balance} & solde & \eng{Could you check my balance, please?} \\
\cle{bank statement} & relevé bancaire & \eng{I need my bank statement for the last three months.} \\ \hline
\end{tabularx}
\end{center}
\textbf{Phrases utiles :} \eng{I would like to open an account.} / \eng{How can I transfer money abroad?} / \eng{Could you repeat that, please?}
\end{bilan}

\begin{aretenir}[Checklist avant le Bac]
\begin{itemize}
\item[$\square$] Je connais le vocabulaire bancaire de base (\eng{account, deposit, withdraw, transfer, balance, statement}).
\item[$\square$] Je sais ouvrir une conversation au guichet et demander une information poliment (\eng{Could you tell me\ldots ?}).
\item[$\square$] Je sais reformuler et confirmer une information (\eng{So, you mean\ldots ?}).
\item[$\square$] Je prononce correctement les montants et les dates.
\item[$\square$] J'évite les faux amis (\eng{actually} $\neq$ actuellement, \eng{to assist} $\neq$ assister).
\item[$\square$] Je maîtrise les prépositions usuelles (\eng{at the bank, in cash, by card}).
\end{itemize}
\end{aretenir}


\findecours

\end{document}
